\documentclass[10pt]{article}
\usepackage{arydshln}
\usepackage{graphicx}
\usepackage{tcolorbox}
\usepackage[a4paper, margin=1in]{geometry}
\usepackage{amsmath, amsfonts, amssymb}
\usepackage[colorlinks=true,linkcolor=black,citecolor=black,urlcolor=black]{hyperref}
\usepackage{natbib}
\usepackage{url}
\usepackage{enumitem}
\usepackage{multicol}

% Bibliography section
% \bibliographystyle{plain}
% \bibliography{refs}

% Optional: include all references
% \nocite{*}

% In-text citation
% \cite{ref of cite}

% ---- Macros ----
\newcommand{\boxednote}[1]{%
    \fbox{%
        \parbox{\dimexpr\linewidth-2\fboxsep-2\fboxrule\relax}{#1}%
    }%
}
\newcommand{\myhrule}[1]{%
  \vspace{#1}%
  \noindent\hrule
  \vspace{#1}%
}
\makeatletter
\newcommand{\bern}{%
  \begin{enumerate}%
  \renewcommand{\labelenumi}{\roman{enumi}.}%
}
\newcommand{\eern}{%
  \end{enumerate}%
}
\makeatother
\newcommand{\FlaThreeByThreeTL}[9]{%
\left(
\begin{array}{cc|c}
#1 & #2 & #3\\
#4 & #5 & #6\\
\hline
#7 & #8 & #9
\end{array}
\right)%
}
\tcbuselibrary{breakable}
\newcommand{\boxedtext}[1]{%
\begin{tcolorbox}[
    colback=white,
    colframe=black,
    sharp corners,
    boxrule=0.5pt,
    breakable
]%
#1
\end{tcolorbox}%
}
\newcommand{\bt}[1]{%
\begin{tcolorbox}[
    colback=white,
    colframe=black,
    sharp corners,
    boxrule=0.5pt,
    breakable
]%
#1
\end{tcolorbox}%
}
\newcommand{\FlaTwoByOne}[2]{%
\left(
\begin{array}{c}
#1\\
\hline
#2
\end{array}
\right)%
}
\newcommand{\FlaThreeByOneB}[3]{%
\left(
\begin{array}{c}
#1\\
\hline
#2\\
#3
\end{array}
\right)%
}
\newcommand{\FlaTwoByTwoSingleLine}[4]{%
\left(
\begin{array}{cc}
#1 & #2\\
#3 & #4
\end{array}
\right)%
}
\newcommand{\complexpm}{\boxed{\pm}}
\newcommand{\FlaOneByTwoSingleLine}[2]{%
\left(
\begin{array}{cc}
    #1 & #2
\end{array}
\right)%
}
\newcommand{\FlaTwoByOneSingleLine}[2]{%
\left(
\begin{array}{cc}
#1\\
#2
\end{array}
\right)%
}
\newcommand{\FlaThreeByThreeBR}[9]{%
\left(
\begin{array}{c|cc}
#1 & #2 & #3 \\ \hline
#4 & #5 & #6 \\
#7 & #8 & #9
\end{array}
\right)%
}
\newcommand{\FlaTwoByTwo}[4]{%
\left(
\begin{array}{c|c}
#1 & #2 \\
\hline
#3 & #4
\end{array}
\right)
}
\newcommand{\FlaOneByThreeR}[3]{%
\left(
\begin{array}{c|cc}
#1 & #2 & #3
\end{array}
\right)%
}
\newcommand{\FlaOneByTwo}[2]{%
\left(
\begin{array}{c|c}
#1 & #2
\end{array}
\right)%
}
\newcommand{\Span}{\operatorname{Span}}
\newcommand{\pbreak}{\vspace{30pt}}
\newcommand{\ud}[1]{\underline{#1}}
\newcommand{\ital}[1]{\textit{#1}}
\newcommand{\fb}[1]{\fbox{#1}}
\newcommand{\vocab}[1]{\fbox{\textbf{\textit{#1}}}}
\newcommand{\eps}{\varepsilon}
\newcommand{\mb}{\mathbf}
\newcommand{\bs}{\boldsymbol}
\newcommand{\mdef}{\overset{\mathrm{def}}{=}}
\newcommand{\h}[1]{\textbf{\fbox{#1}}}
\newcommand{\real}{\mathbb{R}}
\newcommand{\R}{\mathbb{R}}
\newcommand{\C}{\mathbb{C}}
\newcommand{\Null}{\mathcal{N}}
\newcommand{\amp}{&}
\newcommand{\tpose}{^{\top}}
\newcommand{\ident}[1][]{\mathbf{I}_{#1}}
\newcommand{\defeq}{\overset{\cdot}{=}}
\newcommand{\sbsi}[4]{%
	\begin{center}
   		\begin{minipage}{0.5\textwidth}
      			\centering
      			\includegraphics[width=#2\textwidth]{#1}
    		\end{minipage}%
    		\begin{minipage}{0.5\textwidth}
      			\centering
      			\includegraphics[width=#4\textwidth]{#3}
		\end{minipage}
 	\end{center}
}
\newcommand{\ci}[2]{%
  \begin{center}
	\includegraphics[width=#2\textwidth]{#1}
  \end{center}
}
\newcommand{\cii}[4]{%
  \begin{center}
    \includegraphics[width=#2\textwidth]{#1}\hfill
    \includegraphics[width=#4\textwidth]{#3}
  \end{center}
}
\newcommand{\newbp}{%
  \newpage
  \mbox{}
  \newpage
}
\newcommand{\formbox}[1]{%
  \begin{center}
    \fbox{\parbox{0.95\linewidth}{\[
      #1
    \]}}
  \end{center}
}
\newcommand{\euler}{\mathrm{e}}
\newcommand{\expect}{\mathbb{E}}
\newcommand{\bi}{\begin{itemize}}
\newcommand{\ei}{\end{itemize}}
\newcommand{\be}{\begin{enumerate}}
\newcommand{\ee}{\end{enumerate}}
\newcommand{\al}[1]{\ensuremath{\begin{aligned}#1\end{aligned}}}
\DeclareMathOperator*{\argmax}{arg\,max}
\DeclareMathOperator*{\argmin}{arg\,min}
\newcommand{\rightarrowname}[1]{\ensuremath{\xrightarrow{\text{#1}}}}
\newcommand{\leftarrowname}[1]{\ensuremath{\xleftarrow{\text{#1}}}}
\newcommand{\lt}{<}
\newcommand{\gt}{>}
\setlength{\parindent}{0pt}
\newcommand{\tab}{\hspace*{2em}}
\newcommand{\pause}{}
% ------------------------------
\DeclareMathSizes{10}{13}{8}{8}

% optionally:
% \setcounter{secnumdepth}{0} % removes section numbers
\renewcommand{\arraystretch}{1.5}
% \DeclareMathSizes{<text size>}{<math size>}{<script size>}{<scriptscript size>}

\title{Sam's Teach Yourself C++17 Notes}
\author{}
\date{}

\begin{document}

\maketitle

\tableofcontents

\newpage
\section{Getting Started}

\subsection{History of C++}
\begin{itemize}
	\item Initially developed by Bjarne Stroustroup at Bell Labs in 1979.
	\item It was meant to be a successor to C.
	\item C++ was designed to be an \textbf{object-oriented} language with concepts like
	\begin{itemize}
		\item inheritance
		\item abstraction
		\item polymorphism
		\item encapsulation
	\end{itemize}
	\item In 1998, the first standward of C++ was ratified by the ISO Committee. \\ (This is when we get C++14 for example for the year it is ratified (2014) again and so on)
	\item C++ can be for operating systems, web services, databases, and much more.
\end{itemize}

\newpage
\subsection{Programming a C++ Application}
When you start Notepad on Windows or a terminal on Linux, you are running an \vocab{executable} of that program. The \textbf{executale} is the finished product that can be run.

\subsubsection{Steps to Generating an Executable}
\begin{enumerate}
	\item \textbf{Writing} (or programming) C++ code using a text editor/IDE.
\item \textbf{Compiling} code using a C++ \vocab{compiler} that converts it to a \textbf{machine language} version contained in "\vocab{object files}" \\ 
		(gcc (gnu compiler collection) or clang (llvm) or msvc (microsoft visual c++), etc)
		\begin{quote}
		gcc example: g++ main.cpp -o main\\ 
		clang example: clang++ main.cpp -o main\\ 
		msvc example: cl main.cpp
		\end{quote}
	\item \textbf{Linking} the output of the compiler using a \vocab{linker} to get an \vocab{executable} (.exe for Windows, for example)
\end{enumerate}

\begin{verbatim}
main.cpp
   |
   v
Compiler (g++, clang++, MSVC)
   |
   v
main.o (Linux/macOS) or main.obj (Windows)
   |
   v
Linker
   |
   v
program.exe (Windows)
program     (Linux/macOS)
\end{verbatim}

\pbreak

\textbf{Compilation} is the step where code in C++ (typically text files with .cpp extension) is converted into \vocab{byte code} that the processor can \textbf{execute}.\\ 

The \vocab{Compiler} converts one code file at a time, generating an \textbf{object file} with a \textbf{.o} (linux) or \textbf{.obj} (Windows) extension, \textbf{ignoring dependencies} that this cpp file may have on code in another file.\\ 

The \vocab{Linker} joins the dots and resolves these \textbf{dependencies}.\\ 

In the event of successful \textbf{linkage}, it creates an \textbf{executable} for the programmer to execute and distribute.\\ 

This whole process is called \textit{\textbf{building an exectuable}}.

\pbreak 

Recall \textbf{.c} refers to C files and not C++ files.\\ 

Stick to vim little boy.

\newpage
\subsubsection{Programming Your First C++ Application}
Note all code will be in ../code dir.\\ 

\fbox{\textbf{Listing 1.1} Hello.cpp, the Hello World Program}\\ 

Using linux, you can create the executable via 
\begin{verbatim}
g++ -o hello Hello.cpp
\end{verbatim}

This command instructs g++ to create an \textbf{exectuable} named \verb|hello|\\ 

clang equivalent command:
\begin{verbatim}
clang++ -o hello Hello.cpp
\end{verbatim}

\pbreak

Both g++ and clang++ support many of the same options or \vocab{flags}, for example:

\begin{verbatim}
g++ -std=c++20 -Wall -Wextra -O2 -o hello Hello.cpp
clang++ -std=c++20 -Wall -Wextra -O2 -o hello Hello.cpp
\end{verbatim}

\begin{itemize}
    \item \texttt{-std=c++20}: Compiles the program using the C++20 language standard. Other common standards include \texttt{c++11}, \texttt{c++14}, \texttt{c++17}, and \texttt{c++23}.

    \item \texttt{-Wall}: Enables a broad set of commonly useful compiler warnings. While it does not enable \emph{all} warnings, it catches many common programming mistakes.

    \item \texttt{-Wextra}: Enables additional warnings beyond those provided by \texttt{-Wall}. Using both flags together is considered good practice during development.

    \item \texttt{-O2}: Enables a moderate level of compiler optimizations to improve the performance of the generated executable while keeping compilation times reasonable. Other optimization levels include:
    \begin{itemize}
        \item \texttt{-O0}: No optimization (default for debugging).
        \item \texttt{-O1}: Basic optimizations.
        \item \texttt{-O2}: Recommended general-purpose optimization.
        \item \texttt{-O3}: More aggressive optimizations that may increase executable size.
        \item \texttt{-Ofast}: Enables aggressive optimizations that may sacrifice strict standards compliance.
    \end{itemize}

    \item \texttt{-g}: Includes debugging information in the executable so that a debugger (such as \texttt{gdb} or \texttt{lldb}) can display source code lines, variables, stack frames, and other useful debugging information. This flag does \emph{not} make the program slower by itself; it simply embeds additional symbolic information.

    \item \texttt{-o hello}: Specifies the name of the output executable. Without this flag, the default executable name is typically \texttt{a.out} on Linux/macOS or \texttt{a.exe} on Windows (depending on the compiler).
\end{itemize}

\pbreak

Some more feature of compilers include \vocab{Preprocessor Definitions}\\ 

The \texttt{-D} compiler flag defines a preprocessor macro at compile time. For example,

\begin{verbatim}
g++ -DDEBUG -o hello Hello.cpp
\end{verbatim}

is equivalent to having

\begin{verbatim}
#define DEBUG
\end{verbatim}

before preprocessing begins. This allows for \emph{conditional compilation}:

\begin{verbatim}
#ifdef DEBUG
    std::cout << "Debug mode enabled\n";
#endif
\end{verbatim}

The code inside the \texttt{\#ifdef DEBUG} block is only compiled when the \texttt{DEBUG} macro is defined. This technique is commonly used to enable debugging, logging, assertions, or other development-only features.

\pbreak

Remember to actually execute the executable by calling it like ./hello on Linux or hello.exe on Windows.\\ 

A thing to remember is through the evolution of C++, the same program can be compiled using different compilers on different operating systems.

\newpage
\subsection{QA}
\begin{itemize}
	\item \textbf{Q: Can I ignore warning messages from my compiler?}

	A: These warnings show to the programmer a better/safer way to go about things.\\ 
	It is up to them to decide to ignore or not. Generally, it's best to address them for production level code.\\

	\item \textbf{Q: How does an interpreted language differ from a compiled language?} 

		A: Languages such as Python or Javascript are \textbf{interpreted} (for the most part). There is \textbf{no compilation step}.

		An \vocab{interpreted language} use an \textbf{interpreter} that directly reads the script text file (code) and performs the desired actions.

		Consequently, you need to have the interpreter installed on a machine where the script needs to be executed; consequently, \textbf{performance usually takes a hit} as the interpreter works as a \underline{\textbf{runtime translator}} between the microprocessor and the code written.\\

	\item \textbf{Q: What are runtime errors, and how are they different from compile-time errors?}

		A: Errors that happen when you execute your application are called \vocab{runtime errors}.

	\vocab{Compile-time} errors do not reach the end user and are typically caused by \textbf{syntax} (e.g., missing semicolons, unmatched braces) or \textbf{semantic} (e.g., undeclared variables, type mismatches) errors. They prevent the compiler from successfully generating an executable.
\end{itemize}

\newpage 

\subsection{Quiz}
\begin{enumerate}
	\item What is the difference betwwen an interpreter and compiler?

	A: An \textbf{interpreter} is a tool that interprets what you code (or an \textbf{intermediate byte code}) and performs certain actions. 

	A \textbf{compiler} takes your code as an input and generates an \textbf{object file}. In the case of C++, after \textbf{compiling} and \textbf{linking} have have an \textbf{executable} that can run directly by the processor without need for any further interpretation.\\ 

	\item What does the linker do?

		A: A compiler takes a C++ file as input and generates an \textbf{object file} in machine language. 

		Often your code has \textbf{dependencies} on libraries and functions in other code files. Creating these \textbf{links} and generating an executable that integrates all dependencies is the job of the \textbf{linker}.\\

	\item What are the stages in the normal development cycle?

		A: Code. Compile to create object file. Link to create executable. Execute. Debug. Fix any errors. Repeat.

		In many/most cases compiling and linking are one step.
\end{enumerate}

\newpage
\subsection{Exercies}
\begin{enumerate}
	\item \fbox{E1}
	\item \fbox{E3}
\end{enumerate}

\newpage
\section{The Anatomy of a C++ Program}
C++ programs are comprised into \vocab{classes} comprising \vocab{member functions} and \vocab{member variables/fields}.

\pbreak

\subsection{Parts of the Hello World Program}

Analyze ts:\\

\fbox{\textbf{Listing 2.1} HellowWorldAnalysis.cpp: Analyze a Simple C++ Program.}\\ 

The C++ program can be classified into two parts:
\begin{enumerate}
	\item the \vocab{preprocessor directives} that start with a \# and
	\item the \vocab{main body} of the program that starts with \verb|int main()|.
\end{enumerate}

\newpage 

\subsection{Preprocessor Directive \#include}
The \vocab{preprocessor} is a tool that runs before the actual \textbf{compilation} starts.

\textbf{Preprocessor directives} are commands to the \textbf{preprocessor} and always start with a \# symbol.

\begin{quote}
\texttt{\#include <iostream>} \\
\texttt{\#define PI 3.14159} \\
\texttt{\#ifdef DEBUG} \\
\texttt{\#endif} \\
\texttt{\#ifndef HEADER\_H} \\
\texttt{\#endif}
\end{quote}

\pbreak

In our code Listing 2.1, the \verb|#include <filename>| tells the preprocessor to take the contents of the file (\verb|iostream| in our example) and inc9ulude them at the line where the directive is made.\\

\verb|iostream| is a \vocab{standard header file} that enables the usage of \verb|std::cout| to display "Hello World" on the screen.\\

The compiler was able to compile the code that contains \verb|std::cout| b/c we instructed the preprocessor to include the definition of \verb|std::cout|

\pbreak 

\textbf{Note:} We can include our own header files using
\verb|#include "relative/path/to/FileB.h"|.\\
Use \textbf{double quotes} (\verb|""|) for project headers and \textbf{angle brackets} (\verb|<>|) for standard library headers (e.g., \verb|<iostream>|, \verb|<vector>|).\\ 

The underlying rule is:

\#include "..." $\rightarrow$ searches the current/project directories first, then system include paths.\\
\#include $<$...$>$ $\rightarrow$ searches only the system include paths.\\

The syntax ("" vs $<>$) does not define “own vs standard”.\\
The include search path (set with -I) is what actually lets the compiler find your own headers in different folders.

\newpage
\subsection{The Body of Your Program main()}
Following the \textbf{preprocessor directive(s)} is the \vocab{body} of the program characterized by the function \textbf{main()}.\\ 

The execution of a C++ program always starts here. It is a standardized convention that \textbf{main()} is of type \textbf{int}, which stands for \textbf{integer}.

\pbreak

int main() can take many forms in C/C++:

\begin{verbatim}
int main() {
    return 0;
}
\end{verbatim}
Standard C++ entry point with no command-line arguments (also valid in C, though \texttt{int main(void)} is preferred).

\begin{verbatim}
int main()
{
    return 0;
}
\end{verbatim}
Same as above; only the brace style differs.

\begin{verbatim}
int main(void)
{
    return 0;
}
\end{verbatim}
Standard C entry point explicitly specifying that no arguments are accepted.

\begin{verbatim}
int main(int argc, char *argv[])
{
    return 0;
}
\end{verbatim}
Standard entry point that receives command-line arguments.

\begin{verbatim}
int main(int argc, char **argv)
{
    return 0;
}
\end{verbatim}
Equivalent to the previous form since array parameters decay to pointers.

\begin{verbatim}
int main(int argc, char *argv[], char *envp[])
{
    return 0;
}
\end{verbatim}
Implementation-defined form that also receives environment variables.

\begin{verbatim}
int wmain(int argc, wchar_t *argv[])
{
    return 0;
}
\end{verbatim}
Microsoft-specific Unicode entry point using wide-character command-line arguments.

\begin{verbatim}
void main()
{
    /* Non-standard */
}
\end{verbatim}
Accepted by some compilers but not part of the C or C++ standards.

\newpage
\subsection{Returning a Value}
Functions must have a \textbf{return type}. If a function's return type is \verb|void|, it does not return a value, so you do not need to write a \verb|return| statement (although you may write \verb|return;| to exit the function early).\\ 

\textbf{main()} is a function of type \textbf{int}, so it would always return an \textbf{int} or integer.\\ 

\textbf{main()} having a return type of type int can be useful for OSes to determine the state of which the program had exited. For example, in many cases one application is launched by another and the \vocab{parent application} (that launches it) wants to know if the \vocab{child application} (that was launched) has completed its task successfully.\\ 

The programmer can use the \textbf{return value} of \textbf{main()} to convery a success or error state to the \textbf{parent application}.

\pbreak

Conventionally, \textbf{main()} returns:
\begin{itemize}
	\item $\mathbf{0}$ for \textbf{success}
	\item $\mathbf{-1}$ for \textbf{error}
\end{itemize}
However, the return type is an \textbf{integer} and you can still have flexibility to convey many different states of success or failure using the available range of integer return values.\\ 

Side note: C/C++ is case-sensitive, so Int and int aren't the same. Int would have to be a class/struct/typedef/etc and wouldn't be the same as the standard int type.

\newpage
\subsection{Namespaces}
The reason why you used \verb|std::cout| and not just \verb|cout| is that the \vocab{artifact} (cout) is apart of the \vocab{standard (std) \underline{namespace}}.\\ 

So, what are \textbf{namespaces}?\\ 

First, assume that you didn't use the \textbf{namespace qualifier} in invoking \verb|cout| and assume \verb|cout| existed in \textbf{two locations} known to the \textbf{compiler} --- which one should the compiler invoke?\\ 

This trivially causes a conflict and the \textbf{compilation fails}.

\pbreak

\textbf{Namespaces} to the rescue. Namespaces are \textbf{names} given to parts of code that reduc the potential for \textbf{naming conflicts}.\\ 

By invoking \textbf{std::cout}, you are telling the compiler to use that one unique \verb|cout| that is available in the \textbf{std namespace}.\\ 

To use the namespace you simply do:
\begin{verbatim}
using namespace <namespace_name>;

// like 

using namespace std;
\end{verbatim}

\pbreak 

To create a \textbf{namespace}, you do:

\begin{verbatim}
namespace MyNamespace
{
    int x = 10;

    void hello()
    {
        std::cout << "Hello!\n";
    }
}
\end{verbatim}

You access members using the scope resolution operator ::

\begin{verbatim}
int main()
{
    std::cout << MyNamespace::x << '\n';
    MyNamespace::hello();
}
\end{verbatim}

\pbreak

Instead of writing \verb|<namespace_name>::whatever| every time, you can write:

\begin{verbatim}
using namespace MyNamespace;
\end{verbatim}

to make all names in \verb|MyNamespace| directly accessible.

Alternatively, if you only want to bring a single name into scope, use:

\begin{verbatim}
using MyNamespace::whatever;
\end{verbatim}

\newpage 

\textbf{std} stands for the \textbf{standard} namespace and you invokve functions, etc there that have been \textbf{ratified} by the ISO Standards Committee.

\pbreak

\fbox{\textbf{Listing 2.2} The using namespace Declaration}\\

\fbox{\textbf{Listing 2.3} Another Demo of the \textit{using} keyword} \\ 


However, with all this said:\\
\fbox{\textbf{Avoid using namespace ...; in global scope, especially in headers and large programs.}}
this limits damage/errors to a smaller scope.

\pbreak

\begin{itemize}
	\item The difference btw \verb|using namespace std| and \verb|using std::cout or std::endl| is that the \textbf{former} allows \textbf{all artifacts} in the \textbf{std namespace} (cout, cin, etc) to be used without explicit inclusiong of the namespace.
	\item With the \textbf{latter}, the namespace would explicitly be restricted to only \textbf{std::cout} and \textbf{std::endl}.
\end{itemize}

\newpage
\subsection{Functions in C++}
\begin{itemize}
	\item We've already written a function before: \textbf{int main()}.
	\item Functions have a \textbf{return type} like \textbf{int} for \textbf{int main()}.
	\item Other return types include:
	\begin{itemize}
		\item \textbf{void} — the function returns no value.
		\item \textbf{float} — returns a single-precision floating-point number.
		\item \textbf{double} — returns a double-precision floating-point number.
		\item \textbf{char} — returns a single character.
		\item \textbf{bool} — returns a boolean value (true/false).
		\item \textbf{std::string} — returns a string object (C++).
		\item \textbf{int* / double* / etc.} — returns a pointer to a value.
		\item \textbf{custom types (struct/class)} — returns user-defined data types.
	\end{itemize}
	\item We also have \vocab{lambda functions}, which are anonymous functions (they do not have an explicit name).
	\item \fbox{\textbf{Listing 2.4}, Declaring, Defining, and Calling a Function}
\end{itemize}

\pbreak

\begin{itemize}
	\item \vocab{Function declarations} are, for example:
	\begin{verbatim}
int DemoConsoleOutput();
	\end{verbatim}
	This is a function declaration without a body (i.e. no definition). It informs the compiler that a function with this name, return type, and parameter list exists somewhere in the program.

	This allows the function to be called before its actual definition appears in the file. The \textbf{compiler} assumes the \textbf{definition} will be provided later (possibly in the same file or in another translation unit).

	\item \vocab{Function definitions} provide the actual implementation of the function. For example:
	\begin{verbatim}
int DemoConsoleOutput()
{
    std::cout << "Hello\n";
    return 0;
}
	\end{verbatim}
	A function definition includes the function body, where the actual logic is written and executed when the function is called.

	Function definitions allocate the actual code that runs at runtime, whereas declarations only describe the function’s interface.

	Using separate declarations and definitions helps with:
	\begin{itemize}
		\item better code organization (especially in large projects),
		\item separation of interface (.h files) and implementation (.cpp files),
		\item allowing functions to be used across multiple files.
	\end{itemize}
\end{itemize}

\pbreak

\begin{itemize}
	\item \fbox{\textbf{Listing 2.5} Using the Return Value of a Function}
		\item Functions 
			\begin{itemize}

				\item can take \textbf{parameters} (values defined in the function signature that receive input from the caller as \textbf{arguments}), 
				\item can be \textbf{recursive} (a function that calls itself and requires a base case to stop), 
				\item can contain multiple return statements (though if not inside a condition, the compiler may issue a warning), 
				\item can be \vocab{overloaded} (multiple functions with the same name but different parameter lists),
				\item can be \textbf{overwritten} (more correctly: overridden in inheritance, where a derived class replaces a base class function), 
				\item can be \textbf{expanded in-line} by the compiler (the compiler replaces the function call with the function body to reduce call overhead), 
				\item and much more.
			\end{itemize}
\item This is discussed in \textbf{Lesson 7}.
\end{itemize}

\newpage
\subsection{Basic Input Using std::cin and Output using std::cout}
\begin{itemize}
	\item We use the \textbf{console} to write and read information.\\ 
	\item use \textbf{std::cout} (standard "see-out") to \textbf{write} text to the console.
		\begin{quote}
		\textbf{std::cout} is accompanied by $<<$
		\begin{verbatim}
std::cout << "Hello World" << std::endl;
		\end{verbatim}
		\end{quote}
	\item use \textbf{std::cin} (standard "see-in") to \textbf{read} text from the console.
		\begin{quote}
		\textbf{std::cin} is accompanied by $>>$
		\begin{verbatim}
std::cin >> Variable;
		std::cin >> Variable1 >> Variable2;
		\end{verbatim}
		\end{quote}

	\item \fbox{\textbf{Listing 2.6} Use cin and cout}
\end{itemize}

\newpage

\subsection{QA}
\begin{itemize}
	\item \textbf{Q: What does \#include do}?

		A: This is a \textbf{directive} to the \textbf{preprocessor} when you call your compiler. This specific directive causes the contents of the file named in $<>$ after \#include to be \textbf{inserted} at that line as if it were typed at that location in your source code.\\
	
	\item \textbf{Q: When do you need to program command-line arguments?}

	A: To supply \textbf{options} that allow the user to alter the behavior of a program.

	For example, \verb|ls| in Linux or \verb|dir| in Windows enables you to see the contents within the current directory or folder. To view files in another directory, you specify the path of the same using command-line arguments, as in \verb|ls /| or \verb|dir \|.
\end{itemize}

\newpage
\subsection{Quiz}
\begin{itemize}
	\item \textbf{Q: What is the problem in declaring} \verb|Int main()|?

	A: \verb|int main()| is what it must be.\\ 

	\item \textbf{Q: Can comments be longer than one line}?

	A: Yes via \verb|/* */|
\end{itemize}

\newpage
\subsection{Exercises}
\begin{enumerate}
	\item e1.cpp
\end{enumerate}

\newpage
\section{Using Variables, Declaring Constants}
\begin{itemize}
	\item \vocab{Variables} are points in memory that store data in a finite amount of time.
	\item \vocab{Constants} are points in memory that define \textbf{artifacts} thare are \underline{\textbf{not allowed}} to \textbf{change}.
\end{itemize}

\pbreak

\subsection{What Is a Variable?}
\subsubsection{Memory and Addressing}
\begin{itemize}
	\item All computers contain a \vocab{microprocessor} and a certain amount of \textbf{memory} for \textbf{temporary storage} called \vocab{Random Access Memory (RAM)}.
	\item In addition, many devices also allow for data to be \textbf{persisted} on a \vocab{storage device} such as the \textbf{hard disk}.
	\item The \textbf{microprocessor} executes your application, and in doing so works with the \textbf{RAM} to fetch the application \textbf{binary code} to be executed as well as the \textbf{data} associated with it.
\end{itemize}

\pbreak

\begin{itemize}
	\item \textbf{RAM} can be thought as a \textbf{storage} akin to a row of lockers in a hallway, each locker having a \textbf{number} --- that is, an \vocab{address}.
	\item To access a location in \textbf{memory}, say location 578, the processor needs to be asked via an \textbf{instruction} to \textbf{read} a value from there or \textbf{write} a value to it.

	\item There exist \vocab{virtual addresses}, which allow each process to have its \textbf{own address space} \underline{independent} of \textbf{physical memory}. These \textbf{virtual addresses} are translated into \textbf{physical addresses} by the operating system and hardware.

	\item The main purpose of \textbf{virtual addresses} is to provide \textbf{memory isolation}, \textbf{memory protection}, and the illusion that each process has its own large, contiguous block of memory, even though \textbf{physical memory is shared} among many processes.
\end{itemize}

\newpage 

\subsubsection{Declaring Variables to Access and Use Memory}
\begin{itemize}
 \item For variables, it's best to initialize it no matter if you use it, so that you know what kind of value it holds. 
	\item If you don't, you wouldn't know the value, which may cause problems ahead.
	\item \fbox{\textbf{Listing 3.1}}

	\item The compiler does its job of mapping the variable to a location \textbf{in memory} and takes care of the associated \textbf{memory-address} bookkepping for you.
	\item The programmer works with the human-friendly names, while the \textbf{compiler} manages \textbf{memory-addressing} and creates the instructions for the microprocessor to execute in working with the \textbf{RAM}.
\end{itemize}

\newpage

\subsubsection{Declaring and Initializing Multiple Variables of a Type}
\begin{itemize}
	\item You can condense the declaration of variables of the same type to one line of code:
		\begin{verbatim}
int firstNum = 0, secondNum = 0, thirdNum = 0; // and so on
		\end{verbatim}

	\item Note: data stored in variables is data stored in RAM. It is not persistent and is lost upon application termination. To persist it, as of today, you must write it to a file on disk.
\end{itemize}

\newpage

\subsubsection{Scope of a Variable}
\begin{itemize}

	\item When outside the \vocab{scope} of a variable, the \textbf{variable names} will \underline{not} be recognized by the compiler and your program won't compile.
\item Beyond its scope a variable is an \textbf{unidentified entity} that the compiler knows nothing of.

\item \fbox{\textbf{Listing 3.2} Scope}
\end{itemize}

\newpage
\subsubsection{Global Variables}
\begin{itemize}
	\item Not best practice for production projects in source files.
\item All functions within the source file can access global variables declared in that file (provided they are in scope). Functions in other source files cannot access them unless the variables have external linkage (e.g., are declared with \verb|extern| in a header and defined in exactly one source file).

	Example:
	\begin{quote}
		\begin{verbatim}
		// globals.h
		extern int g_value;

		// globals.cpp
		int g_value = 42;
		\end{verbatim}
	\end{quote}

\item \fbox{\textbf{Listing 3.3} Global variables}
\end{itemize}

\newpage 

\subsection{Common Compiler-Supported C++ Variable Types}
\begin{itemize}
    \item \verb|bool| -- stores \verb|true| or \verb|false|.
    \item \verb|char| -- typically 1 byte, range: $-128$ to $127$ (or $0$ to $255$ if unsigned).
    \item \verb|signed char| -- 1 byte, range: $-128$ to $127$.
    \item \verb|unsigned char| -- 1 byte, range: $0$ to $255$.
    \item \verb|short| (or \verb|short int|) -- typically 2 bytes, range: $-32,\!768$ to $32,\!767$.
    \item \verb|unsigned short| -- typically 2 bytes, range: $0$ to $65,\!535$.
    \item \verb|int| -- typically 4 bytes, range: $-2,\!147,\!483,\!648$ to $2,\!147,\!483,\!647$.
    \item \verb|unsigned int| -- typically 4 bytes, range: $0$ to $4,\!294,\!967,\!295$.
    \item \verb|long| -- typically 8 bytes on Linux/macOS and 4 bytes on Windows.
    \item \verb|unsigned long| -- unsigned version of \verb|long|.
    \item \verb|long long| -- typically 8 bytes, range: $-9.22\times10^{18}$ to $9.22\times10^{18}$.
    \item \verb|unsigned long long| -- typically 8 bytes, range: $0$ to $1.84\times10^{19}$.
    \item \verb|float| -- typically 4 bytes, approximately $\pm 3.4\times10^{38}$ (about 7 decimal digits of precision).
    \item \verb|double| -- typically 8 bytes, approximately $\pm 1.7\times10^{308}$ (about 15--16 decimal digits of precision).
    \item \verb|long double| -- extended precision floating-point type (size and precision are implementation-defined).
    \item \verb|void| -- represents the absence of a value (e.g., functions that return nothing).
\end{itemize}

\subsection*{Fixed-Width Integer Types (\texttt{<cstdint>})}
\begin{itemize}
    \item \verb|int8_t| -- exactly 8 bits, range: $-128$ to $127$.
    \item \verb|uint8_t| -- exactly 8 bits, range: $0$ to $255$.
    \item \verb|int16_t| -- exactly 16 bits, range: $-32,\!768$ to $32,\!767$.
    \item \verb|uint16_t| -- exactly 16 bits, range: $0$ to $65,\!535$.
    \item \verb|int32_t| -- exactly 32 bits, range: $-2,\!147,\!483,\!648$ to $2,\!147,\!483,\!647$.
    \item \verb|uint32_t| -- exactly 32 bits, range: $0$ to $4,\!294,\!967,\!295$.
    \item \verb|int64_t| -- exactly 64 bits, range: $-9,\!223,\!372,\!036,\!854,\!775,\!808$ to $9,\!223,\!372,\!036,\!854,\!775,\!807$.
    \item \verb|uint64_t| -- exactly 64 bits, range: $0$ to $18,\!446,\!744,\!073,\!709,\!551,\!615$.
\end{itemize}

\newpage 
\subsubsection{Signed vs Unsigned integers}
\begin{itemize}
	\item \vocab{Sign} implies \textbf{positive} or \textbf{negative}.
	\item All numbers you work with using a computer are stored in the \textbf{memory} in the form of \vocab{bits} and \vocab{bytes}.

	\item A \textbf{memory location} that is \textbf{1 byte} long contains \textbf{8 bits}.

	\item Each \textbf{bit} can either be $\mathbf{0}$ or $\mathbf{1}$.

	\item Thus, a \textbf{memory location} that is \textbf{1 byte} long can contain a maximum of \textbf{2} to the power of \textbf{8 values} --- that is \textbf{256} unique values.

	\item Similarly, a \textbf{memory location} that is \textbf{16 bits} large (or \textbf{2 bytes}) can contain \textbf{2} to the power of \textbf{16 values} --- that is, \textbf{65, 536} unique values.

	\item Think of it as just
		\[
			2^{\text{number of bits}} = \text{quantity of how many numbers we can represent}
		\]
\end{itemize}

\pbreak

\begin{itemize}
	\item If the values are \vocab{unsigned} (i.e., only non-negative), then an \textbf{\(n\)-bit} integer can represent values from \fbox{$0$ to $2^n - 1$}.

    \item For example:
    \begin{itemize}
        \item \textbf{1 byte} ($8$ bits): $0$ to $2^8 - 1 = 255$.
        \item \textbf{2 bytes} ($16$ bits): $0$ to $2^{16} - 1 = 65,\!535$.
    \end{itemize}

    \item For an 8-bit unsigned integer, the place values are:
    \[
        2^7 \quad 2^6 \quad 2^5 \quad 2^4 \quad 2^3 \quad 2^2 \quad 2^1 \quad 2^0
    \]

    If every bit is set to 1,
    \[
        11111111_2 = 2^7 + 2^6 + \cdots + 2^1 + 2^0 = 2^8 - 1 = 255.
    \]
\end{itemize}

\pbreak

\begin{itemize}
    \item To represent \textbf{negative numbers}, signed integers are typically stored using \vocab{two's complement}. Rather than sacrificing the leftmost bit, it is assigned a weight of $-2^{n-1}$, while all remaining bits retain their usual positive weights.

    \item For an \textbf{\(n\)-bit signed integer}, the representable range is
    \[
        -2^{n-1} \text{ to } 2^{n-1} - 1.
    \]

    \item Thus, an \textbf{8-bit signed integer} has the range
    \[
        -2^7 \text{ to } 2^7 - 1 = -128 \text{ to } 127.
    \]

    \item Since there are $2^8 = 256$ unique bit patterns, exactly half represent negative values and half represent non-negative values.

    \item The place values for an 8-bit signed integer are
    \[
        -2^7 \quad 2^6 \quad 2^5 \quad 2^4 \quad 2^3 \quad 2^2 \quad 2^1 \quad 2^0.
    \]

    Therefore,
    \[
        \verb|10000000| = -2^7 = -128,
    \]
    and
    \[
        \verb|11111111| = -128 + 64 + 32 + 16 + 8 + 4 + 2 + 1 = -1.
    \]

\item Recoginize the smallest/minimum is when only the \textbf{MSB} is 1 and the rest are 0, which gives our msb weight of $-2^{n-1}$.
\end{itemize}

\newpage 

\subsubsection{Signed integer types: short, int, long, and long long}

\begin{itemize}
    \item These are the fundamental \textbf{signed integer types} in C++, used to store whole numbers (both positive and negative).
    \item They primarily differ in the amount of memory they occupy and, consequently, the range of values they can represent.
    \item In general, \verb|short| stores smaller integers than \verb|int|, while \verb|long| and \verb|long long| can store progressively larger values.
    \item The exact size of each type is \textbf{implementation-defined} (depends on the compiler and platform), although modern systems typically use:
    \begin{itemize}
        \item \verb|short|: 2 bytes
        \item \verb|int|: 4 bytes
        \item \verb|long|: 4 bytes (Windows) or 8 bytes (Linux/macOS)
        \item \verb|long long|: 8 bytes
    \end{itemize}
    \item When an exact integer size is required (e.g., exactly 32 or 64 bits), prefer the fixed-width integer types from \verb|<cstdint>| such as \verb|int32_t| and \verb|int64_t|.
\end{itemize}

\newpage
\subsubsection{Unsigned integer types: unsigned short, unsigned int, unsigned long, unsigned long long}

\begin{itemize}
    \item These are the fundamental \textbf{unsigned integer types} in C++, used to store only \textbf{non-negative} whole numbers.
    \item Since no bits are reserved for representing negative values, unsigned integers can represent approximately twice the positive range of their signed counterparts of the same size.
    \item They have the same typical sizes as their signed versions:
    \begin{itemize}
        \item \verb|unsigned short|: 2 bytes
        \item \verb|unsigned int|: 4 bytes
        \item \verb|unsigned long|: 4 bytes (Windows) or 8 bytes (Linux/macOS)
        \item \verb|unsigned long long|: 8 bytes
    \end{itemize}
    \item An unsigned \textbf{\(n\)-bit} integer has the range
    \[
        0 \text{ to } 2^n - 1.
    \]
    \item Unsigned integers are useful when negative values are impossible (e.g., sizes, counts, and bit manipulation), but care should be taken when mixing signed and unsigned types, as implicit conversions can produce unexpected results.
\end{itemize}

\newpage

\subsubsection{Integer Overflow}

\begin{itemize}
    \item \vocab{Integer overflow} occurs when the result of an arithmetic operation is outside the range that can be represented by the integer's type.

    \item For example, an 8-bit signed integer can only store values from $-128$ to $127$. Attempting to compute
    \[
        127 + 1
    \]
    produces a value that cannot be represented.

    \item For \textbf{unsigned integers}, overflow is \textbf{well-defined} in C++. Arithmetic wraps around modulo $2^n$, where $n$ is the number of bits. For example, an 8-bit unsigned integer satisfies
    \[
        255 + 1 = 0.
    \]

    \item For \textbf{signed integers}, overflow results in \vocab{undefined behavior}. This means the C++ standard does not specify what should happen---the program may appear to wrap around, produce an unexpected value, or behave unpredictably.

    \item To avoid overflow, choose an integer type with a sufficiently large range (e.g., \verb|int64_t| instead of \verb|int32_t|) or check whether an operation would exceed the type's limits before performing it.
\end{itemize}

\pbreak

\begin{itemize}
	\item \fbox{\textbf{Listing 3.4} Signed and Unsigned overflow}

	\item The output shows the wrap around effect of incrementing and overflowing for both types.
\end{itemize}

\newpage

\subsubsection{Floating-Point types \texttt{float} and \texttt{double}}

\begin{itemize}
    \item \vocab{Floating-point types} are used to store \textbf{real numbers} (numbers with a fractional or decimal component), such as \verb|3.14| or \verb|-0.001|.

    \item The two most commonly used floating-point types are:
    \begin{itemize}
        \item \verb|float| -- typically occupies \textbf{4 bytes} and provides about \textbf{7 decimal digits} of precision.
        \item \verb|double| -- typically occupies \textbf{8 bytes} and provides about \textbf{15--16 decimal digits} of precision.
    \end{itemize}

    \item Floating-point numbers are generally represented according to the \vocab{IEEE 754} standard, which stores a number using three components:
    \begin{itemize}
        \item a \textbf{sign bit},
        \item an \textbf{exponent},
        \item and a \textbf{significand (mantissa)}.
    \end{itemize}

    \item Since only a finite number of bits are available, many decimal values (such as \verb|0.1|) cannot be represented exactly. As a result, floating-point arithmetic may introduce small rounding errors.

    \item Floating-point types should not be used when exact decimal precision is required (e.g., financial calculations). Instead, use an appropriate decimal or fixed-point representation.

    \item In general, prefer \verb|double| over \verb|float| unless reduced memory usage or compatibility with specific hardware is important.
\end{itemize}

\pbreak

\begin{itemize}

\item Under the hood, \verb|float| and \verb|double| are stored using the \vocab{IEEE 754 floating-point format}. A number is represented in \textbf{scientific notation in base 2}:
\[
(-1)^{\text{sign}} \times 1.\text{fraction} \times 2^{\text{exponent}}
\]

\item A \verb|float| (32 bits) is split into:
\begin{itemize}
    \item 1 bit for the \textbf{sign}
    \item 8 bits for the \textbf{exponent}
    \item 23 bits for the \textbf{fraction (mantissa)}
\end{itemize}

\item A \verb|double| (64 bits) is split into:
\begin{itemize}
    \item 1 bit for the \textbf{sign}
    \item 11 bits for the \textbf{exponent}
    \item 52 bits for the \textbf{fraction (mantissa)}
\end{itemize}

\item For example, the value $5.75$ is stored as:
\[
5.75_{10} = 101.11_2 = 1.0111 \times 2^2
\]
So the stored components are:
\begin{itemize}
    \item sign = 0 (positive)
    \item exponent = 2 (stored with a bias)
    \item fraction = 0111 (remaining bits after the leading 1)
\end{itemize}

\end{itemize}

\pbreak

For the above example: $5.75$
\begin{itemize}
    \item Whole part: $5$
    \[
        5 = 101_2
    \]

    \item Fractional part: $0.75$
    \begin{itemize}
        \item $0.75 \times 2 = 1.5 \rightarrow 1$
        \item $0.5 \times 2 = 1.0 \rightarrow 1$
    \end{itemize}
    \[
        0.75 = 0.11_2
    \]

    \item Combine whole and fraction:
    \[
        5.75 = 101.11_2
    \]

    \item Normalize (move binary point left 2 places):
    \[
        101.11_2 = 1.0111 \times 2^2
    \]
\end{itemize}

\pbreak

\begin{itemize}
\item \textbf{Why we multiply by 2 repeatedly for the 2nd step?}

We are converting a decimal fraction (base 10) into a binary fraction (base 2).

In binary, each fractional bit represents a negative power of two:
\[
\frac{1}{2}, \quad \frac{1}{4}, \quad \frac{1}{8}, \quad \frac{1}{16}, \dots
\]

We therefore need a method to determine which of these powers of two are present in the number.

Multiplying by 2 repeatedly allows us to extract these bits one at a time by shifting the fraction:
\begin{itemize}
    \item the integer part after multiplication gives the next binary digit,
    \item the remaining fractional part is used for the next step.
\end{itemize}

\end{itemize}

\pbreak

\textbf{Converting back to base 10:}

We start from the binary result:
\[
5.75 = 101.11_2
\]

Split into integer and fractional parts:
\[
101.11_2 = 101_2 + 0.11_2
\]

Convert each part separately:

\begin{itemize}
    \item Integer part:
    \[
        101_2 = 1 \cdot 2^2 + 0 \cdot 2^1 + 1 \cdot 2^0 = 4 + 0 + 1 = 5
    \]

    \item Fractional part:
    \[
        0.11_2 = 1 \cdot 2^{-1} + 1 \cdot 2^{-2} = \frac{1}{2} + \frac{1}{4} = 0.75
    \]
\end{itemize}

Therefore:
\[
101.11_2 = 5.75_{10}
\]

\newpage

\subsection{Size of a Variable using sizeof}

\begin{itemize}
	\item \textbf{Size} is the amount of \textbf{memory} that the compiler reserves when a variable is declared.
	\item The size of a variable depends on its \textbf{type}.
	\item We can see the size via \verb|sizeof(<type>)| and this outputs the \fbox{\textbf{size of type in bytes}}.

	\item \fbox{\textbf{Listing 3.5} Finding the Size of Standard C++ Variable Types}

\item Note these sizes are \textbf{dependent} on the 32 or 64-bit \textbf{compiler} (not necessarily an OS).
\end{itemize}

\newpage

\subsection{Automatic Type Inference Using auto}
\begin{itemize}
	\item Compilers supporting C++11 and beyond give the option of \textbf{not} having to \textbf{explicitly} specify the variable type when using the keyword \vocab{auto}

		\begin{verbatim}
auto coinFlippedHeads = true;
		\end{verbatim}

	\item We leave the task of defining an exact type for variable \verb|coinFlippedHeads| to the \textbf{compiler}.

	\item The \textbf{compiler} checks the nature of the value the variable is initialized to and then decides the best possible type that suits this variable.

	\item In the above example, the compiler determines \verb|bool| as the type that suits variabled \verb|coinFlippedHeads| and \textbf{internally} treats \verb|coinFlippedHeads| as a \verb|bool|.

	\item \fbox{\textbf{Listing 3.6} Using auto}

	\item When you \textbf{don't initialize} a variable of type \textbf{auto}, you get a \textbf{compile error} like:
		\begin{verbatim}
error: declaration of ‘auto willGetError’ has no initializer
    9 |         auto willGetError;
		\end{verbatim}
\end{itemize}

\newpage
\subsection{typdef and/or using to Substitute a Variable's Type}
\begin{itemize}
	\item You can subsitute a variables types to something that you might find convenient.

	\item You would use the keyword \vocab{typedef} for that.

	\item Here is an example where a programmer wants to call an \textbf{unsigned int} a descriptive \verb|STRICTLY_POSITIVE_INTEGER|.

		\begin{verbatim}
typedef unsigned int STRICTLY_POSITIVE_INTEGER
STRICTLY_POSITIVE_INTEGER numEggsInBasket = 4532;
		\end{verbatim}

	\item Note, an equivalent keyword now in practice in modern C++ (introducd in C++11) is the keyword \vocab{using}:

	\begin{verbatim}
// Modern C++ equivalent
using STRICTLY_POSITIVE_INTEGER = unsigned int;
STRICTLY_POSITIVE_INTEGER numEggsInBasket = 4532;
	\end{verbatim}

\item When compiled, the first line tells the compiler that a \verb|STRICTLY_POSITIVE_INTEGER| is an \textbf{unsigned int}. At later stages when the compiler encounters the already defined type \verb|STRICTLY_POSITIVE_INTEGER|, it substitutes it for \textbf{unsigned int} and continues compilation.

\item \textbf{typedef} or \textbf{type substitution} is useful for when dealing with complex types, for example types that use \textbf{templates}.
\end{itemize}

\newpage

\subsection{What is a Constant?}
\begin{itemize}
	\item C++ enables you to define a variable as \textbf{constant} that means it \textbf{cannot change (its value)} after \textbf{declaration}.

	\item Assignment to a constant cause compilation errors.

	\item Constants can be:
		\begin{itemize}
			\item Literal constants
			\item Declared constants using the \verb|const| keyword
			\item Constant expressions using the \verb|constexpr| keyword (new since C++11)
			\item Enumerated constants using the \verb|enum| keyword
			\item Defined constants that are not recommended and deprecated.
		\end{itemize}
\end{itemize}

\pbreak

\subsubsection{Literal constants}
\begin{itemize}
	\item \textbf{Literal constants} can be of many types --- \textbf{integer, string, and so on}.
	\item In our Hello World program, here "Hello World" is a \vocab{string literal constant} (or just \textbf{string constant}) meaning they cannot be modified.
	\begin{quote}
	String literals are \textbf{immutable} in C++. Attempting to modify a string literal results in \textbf{undefined behavior}. For example:

	\begin{verbatim}
	const char* str = "Hello";
	str[0] = 'J';   // Error: cannot modify a string literal

	char* ptr = "Hello";
	ptr[0] = 'J';   // Undefined behavior
	\end{verbatim}

	If you need a modifiable sequence of characters, initialize a character array instead:

	\begin{verbatim}
	char str[] = "Hello";
	str[0] = 'J';   // OK

	std::cout << str;   // Prints: Jello
	\end{verbatim}

	Here, the compiler creates a writable character array by copying the contents of the string literal.
	\end{quote}

	\pbreak

\begin{quote}

	Note: for \verb|const char* str = "Hello";|
	the pointer itself is \emph{not} placed in read-only memory---the \textbf{string literal} is.\\

For example:

\begin{verbatim}
const char* str = "Hello";
\end{verbatim}

In this declaration:

\begin{itemize}
    \item \verb|"Hello"| is a \vocab{string literal}, and its characters are typically stored in read-only memory.
    \item \verb|str| is a pointer variable (usually stored on the stack if it is a local variable). It stores the address of the first character of the string literal.
    \item The \verb|const| qualifier prevents the characters from being modified through \verb|str|.
\end{itemize}

In contrast,

\begin{verbatim}
char str[] = "Hello";
\end{verbatim}

creates a writable character array by copying the contents of the string literal into the array. Since the array is a separate copy, modifying it is valid:

\begin{verbatim}
str[0] = 'J';   // OK
\end{verbatim}

\end{quote}

\item For \vocab{int literals}, we have also been using them all along in our programs. An integer literal such as \verb|42| or \verb|0| is simply a fixed value written directly in the source code.

Unlike variables, an integer literal is not a named object in memory that we can modify. Instead, it is used by the compiler as a constant value when evaluating expressions or initializing variables.

For example:

\begin{verbatim}
int x = 42;
x = x + 1;
\end{verbatim}

Here, \verb|42| is an integer literal that is copied into the variable \verb|x|. We are not modifying the literal itself; we are modifying the variable \verb|x| that stores a copy of that value.

\item This is why the following declaration makes sense:

\begin{verbatim}
const char* p = "Hello";
\end{verbatim}

The string literal \verb|"Hello"| corresponds to a real array of characters stored in memory (typically in a read-only section of the program). Therefore, \verb|p| can hold the address of its first character.

In contrast, the following is invalid:

\begin{verbatim}
int* p = &42;   // invalid
\end{verbatim}

This does not work because \verb|42| is an integer literal, not a memory object. It is a value used directly by the compiler, and it does not have a valid, addressable memory location. Therefore, taking its address using \verb|&| is not allowed.

\item \vocab{Binary literals} were introduced in \textbf{C++14}. They allow writing numbers in base 2 using the prefix \verb|0b| or \verb|0B|. For example:

\begin{verbatim}
int x = 0b1010;   // 10 in decimal
\end{verbatim}

\item \vocab{Octal literals} have existed since early C++. They use a leading zero \verb|0|. For example:

\begin{verbatim}
int x = 075;   // 61 in decimal
\end{verbatim}

Here, \verb|075| is interpreted as an octal (base-8) number rather than a decimal.

\item \vocab{User-defined literals} were introduced in \textbf{C++11}. They allow programmers to define custom suffixes for literals, enabling more readable and expressive code.

For example, we can define a literal to represent distances in meters:

\begin{verbatim}
long double operator"" _m(long double x) {
    return x;
}

long double distance = 5.0_m;
\end{verbatim}

Here, \verb|5.0_m| is a user-defined literal, where the suffix \verb|_m| converts the number into a meaningful unit.

User-defined literals can be used with integers, floating-point numbers, characters, and strings, and they are commonly used in libraries to improve readability (e.g., time durations, units, or complex numbers).

\end{itemize}

\newpage 

\subsubsection{Declaring Variables as Constants Using const}
\begin{itemize}
	\item \verb|const| is meant to be used \underline{\textbf{before}} the variable type is declared.

	\begin{verbatim}
const type-name constant-name = value;
	\end{verbatim}

	\item \fbox{\textbf{Listing 3.7} Declaring a Constant Called pi}

	\item Here is the compile-error if you try to modify a const
	\begin{verbatim}
const double pi = 22.0 / 7;
pi = 345;

error: assignment of read-only variable ‘pi’
   11 |         pi = 345;
	\end{verbatim}

\item It's good programming practice to define variables that are \textbf{not supposed to change} as \verb|const| since it shows the programmer has thought about ensuring the correctness and constantness of data (acts as protection).
\end{itemize}

\newpage 

\subsubsection{Constant Expression Using constexpr}

\begin{itemize}
	\item Keyword \vocab{constexpr} allows function-like declartions of \textbf{constants}:
	\begin{verbatim}
	constexpr double GetPi() {return 22.0 / 7;}
	\end{verbatim}

\item One \textbf{constexpr} can use another:
	\begin{verbatim}
constexpr double TwicePi() {return 2 * GetPi();}
	\end{verbatim}

\item \verb|constexpr| make look like a \textbf{function}, \underline{however}, allows for optimization possibilities from the compiler's view. So long as a compiler is capable of evaluting a \textbf{constant expression} to a \textbf{constant}, it can be used in statements and expressions at places where a constant is expected.\\

\item In the preceding example, \verb|TwicePi()| is a \textbf{constexpr} that uses a \textbf{constant expression} \verb|GetPi()|.

	This will possibly trigger a \textbf{compile-time optimization} wherein \underline{every usage} of \verb|TwicePi()| is simply \textbf{replaced} by $6.28571$ by the compiler, and \textbf{\underline{not}} the code that would calculate $2 \times 22 / 7$ when executed.\\

\item We can also combine \vocab{constexpr} with \vocab{const} when we want to emphasize that a value is both \textbf{compile-time computable} and \textbf{read-only after initialization}. While \verb|constexpr| already implies constness for the expression, adding \verb|const| can still be useful for clarity in variable declarations.

For example:

\begin{verbatim}
constexpr double GetPi() { return 22.0 / 7; }

const double pi1 = GetPi();          // runtime or compile-time constant
constexpr double pi2 = GetPi();      // guaranteed compile-time constant
\end{verbatim}

Here, \verb|pi2| must be evaluated at compile time, while \verb|pi1| is simply a read-only value initialized from a constant expression. In practice, \verb|constexpr| is preferred when compile-time evaluation is desired, since it already enforces immutability and constant-expression rules.\\ 

\item We can also place \vocab{const} at the end of a member function declaration to indicate that the function does not modify the state of the object. This is different from \verb|const| variables or \verb|constexpr| functions.

For example:

\begin{verbatim}
class Circle {
public:
    double GetRadius() const {
        return radius;
    }

    void SetRadius(double r) {
        radius = r;
    }

private:
    double radius;
};
\end{verbatim}

Here, \verb|GetRadius() const| guarantees that no data members of the object will be modified inside the function. If we attempt to modify any member variable inside a \verb|const| function, the compiler will produce an error.

This \verb|const| qualifier at the end of a function is commonly used for \textbf{getter functions} to ensure they are safe and do not change object state.\\

\item \fbox{\textbf{Listing 3.8} using constexpr}

\item The two \textbf{constexpr}'s may look like functions, but they are not exactly.

	\textbf{Functions} are invoked at program \textbf{execution time}.

	\textbf{Constant expressions} are already substituted to their hard values by the compiler at \textbf{compile time}.

\item So here in the listing, the \verb|constexpr| saves \textbf{time and computation} to the same calculation being contained in a function.

\item Note if the block/body/defintion of the \textbf{constexpr} relies on user input/numerics not known at compile time, then the compiler treats the \textbf{constexpr} as a \textbf{regular function}, i.e. no compile-time optimization.
\end{itemize}

\newpage
\subsubsection{Enumerations}
\begin{itemize}
	\item Sometimes you want a \textbf{variable} whose values are \textbf{restricted} to a certain \textbf{set} of values defined by you.
	\item \vocab{Enumerations} are the tool and are characterized by the \textit{keyword} \vocab{enum}.

	\item Example:

	\begin{verbatim}
enum RaindbowColor
{
	Violet = 0,
	Indigo,
	Blue,
	Green,
	Yellow,
	Orange,
	Red
};
	\end{verbatim}
\item A \textbf{cardinal direction example}:
\begin{verbatim}
enum CardinalDirections
{
	North,
	South,
	East,
	West
};
\end{verbatim}

\item \textbf{Enumerations} are \textbf{user-defined} types.
\item Variables of this \textbf{type} can be assigned to a \textbf{range of values} restricted to the \textbf{enumerators} contained in the \textbf{enumeration}.
\item So, if defined a variable that contains the colors of a rainbow, you declare the variable like this:
	\begin{verbatim}
RainbowColors MyFavColor = Blue;
	\end{verbatim}

\item This works for a traditional (unscoped) enumeration because the enumerator names (such as \verb|Blue|) are placed directly into the surrounding scope.

If the enumeration were declared as a \verb|enum class| (a scoped enumeration, introduced in C++11), the enumerators would remain inside the scope of the enumeration. In that case, we would need to qualify the name:

\begin{verbatim}
RainbowColors MyFavColor = RainbowColors::Blue;
\end{verbatim}

Scoped enumerations help prevent name collisions by requiring the enumeration name when referring to its values.\\ 

\item C++ provides two kinds of enumerations: the traditional \vocab{enum} and the newer \vocab{enum class} (introduced in \textbf{C++11}). The main difference is that an \verb|enum class| is \textbf{scoped}, meaning its enumerator names remain inside the enumeration and do not pollute the surrounding scope.

For example, with a traditional enumeration:

\begin{verbatim}
enum RainbowColors { Red, Orange, Yellow, Green, Blue };

RainbowColors color = Blue;   // OK
\end{verbatim}

The name \verb|Blue| is placed directly into the surrounding scope, so it can be used without qualification. However, this can lead to name conflicts if another enumeration also contains an enumerator named \verb|Blue|.

With a scoped enumeration:

\begin{verbatim}
enum class RainbowColors { Red, Orange, Yellow, Green, Blue };

RainbowColors color = RainbowColors::Blue;   // OK
\end{verbatim}

Here, the enumerator must be qualified with the enumeration name (\verb|RainbowColors::Blue|). This avoids name collisions and provides better type safety, making \verb|enum class| the preferred choice in modern C++.\\ 

\item Generally, the compiler converts the \textbf{enumerator} such as \textit{Violet} and so on into integers.

	Each \textbf{enumerated value} specified is \textbf{one more than} the \textbf{previous value} if not explicitly initialized.

	You have the choice of specifying a \textbf{starting value}, and if this is \underline{\textbf{not specified}}, the compiler takes it as $\mathbf{0}$.

\item You can still specify the explicit value against each of the enumerated constants by initializing them.

	\begin{verbatim}
enum ErrorCode
{
    Success = 0,
    Warning = 100,
    Error = 200
};
\end{verbatim}

\item \fbox{\textbf{Listing 3.9} Enumerated Values}
\end{itemize}

\newpage

\subsubsection{Defining Constants Using \#define}
\begin{itemize}
	\item For C++ this kind of defining is \underline{not recommended}.
	\item This is legacy support (C defines $\rightarrow$ OK in C but we're in C++).

	\item We can defned pi like before as 
		\begin{verbatim}
#define pi 3.14286
		\end{verbatim}

	\item \textbf{\#define} is a \vocab{preprocessor macro}.
	\item What is done here is that all mentions of \verb|'pi'| henceworth are replaced by $3.14286$ for the \textbf{compiler} to process.

	\item Note: the compiler neither knows nor cares about the \textbf{actual type} of the \textbf{constant} in question (it's just a text replacement (read: non-intelligent replacement) done by the \textbf{preprocessor}).
\end{itemize}

\newpage
\subsection{Keywords that Cannot be Used as Variable or Constant Names}

The following are examples of \textbf{reserved keywords} in C++. 

Since these words have special meanings to the compiler, they cannot be used as identifiers (such as variable, constant, or function names).

\begin{multicols}{3}
\begin{itemize}
    \item asm
    \item auto
    \item bool
    \item break
    \item case
    \item catch
    \item char
    \item class
    \item const
    \item constexpr
    \item continue
    \item default
    \item delete
    \item do
    \item double
    \item else
    \item enum
    \item explicit
    \item export
    \item extern
    \item false
    \item float
    \item for
    \item friend
    \item goto
    \item if
    \item inline
    \item int
    \item long
    \item mutable
    \item namespace
    \item new
    \item nullptr
    \item operator
    \item private
    \item protected
    \item public
    \item register
    \item return
    \item short
    \item signed
    \item sizeof
    \item static
    \item struct
    \item switch
    \item template
    \item this
    \item throw
    \item true
    \item try
    \item typedef
    \item typeid
    \item typename
    \item union
    \item unsigned
    \item using
    \item virtual
    \item void
    \item volatile
    \item while
\end{itemize}
\end{multicols}

\newpage

\subsection{QA}
\begin{itemize}
	\item \textbf{Q: Why define constants at all if you can use regular variables instead of them?}

	A: \textbf{Constants}, especially those declared using the keyword \verb|const|, are your way of telling the \textbf{compiler} that the \textbf{value} of a particular variable be \textbf{fixed} and \textbf{not allowed} to \textbf{change}.

	It ensures/enforces correctness and constantness and quality of your code, so that another programmer reading/writing your program would be aware that this const variable is to not be changed.\\ 

\item \textbf{Q: Why should I initialize a variable?}

A: If you don't, you would not know for sure what the variable \textbf{contains} for a starting value.

It would contain junk data at that memory location of the variable, and you cannot check against non-zero values of it cause you wouldn't know that it's not initialized with this check.\\ 

\item \textbf{Q: Why shouldn't I use global variables frequently?}

	A: \textbf{Global variables} can be \textbf{read} and \textbf{assigned} \underline{globablly}. The \textbf{latter} is the problem as they can be changed globally.

	Trivially, other programmers can mistakenly overwrite your global variables and that would effect your code---even in another .cpp file.\\ 

\item \textbf{Q: What happens if I decrement from zero contained in an unsigned int}?

	A: You see a \textbf{wrapping effect}. Decrementing an \textbf{unsigned integer} that contains 0 by 1 means that it wraps to the \textbf{highest} value it can hold!
\end{itemize}

\newpage

\subsection{Quiz}
\begin{itemize}
    \item What is the difference btw a \textbf{signed} and an \textbf{unsigned integer}?

    A:
    A \textbf{signed} integer can store both positive and negative values, whereas an \textbf{unsigned} integer can store only non-negative values (0 and positive numbers). Since unsigned integers do not need to represent negative values, they can represent a larger maximum positive value.\\

    \item Why should you \textbf{not use} \#define to declare constant?

    A:
    Because \verb|#define| performs a simple text substitution and is not type-safe. Modern C++ prefers \verb|const| or \verb|constexpr| because they are checked by the compiler, have a type, and are easier to debug.\\

    \item Why would you initialize a variable?

    A:
    Initializing a variable gives it a known value before it is used, helping to prevent bugs caused by uninitialized (garbage) values.

    \item Consider the enum below. What is the value of Queen?\\

\begin{verbatim}
enum YourCards { Ace, Jack, Queen, King };
\end{verbatim}

    A:
    \verb|Queen| has the value \verb|2|. By default, enumeration values begin at 0 and increase by 1.

    \item What is wrong with this variable name?\\

\begin{verbatim}
int Integer = 0;
\end{verbatim}

    A:
    Nothing is technically wrong with the name. \verb|Integer| is not a reserved C++ keyword, so it is a valid identifier. However, it is not a descriptive variable name and could be confused with the type \verb|int|, so a more meaningful name is recommended.\\

\end{itemize}

\newbp
\section{Managing Arrays and Strings}
\subsection{What is an Array?}
\begin{itemize}
	\item An \vocab{array} is a group of elements forming a \textbf{complete} unit. Example: an array of solar panels.

	\begin{itemize}
		\item is a collection of elements.
		\item all elements contained in an array are of the \textbf{same kind}.
		\item this collection forms a \textbf{complete set}.
	\end{itemize}

\item In C++, \textbf{arrays} enable you to \textbf{store} data elements of a type in \textbf{memory}, in a \underline{\textbf{sequential}} and ordered fashion (cache friendly).
\end{itemize}

\newpage

\subsubsection{The Need for Arrays}
\begin{itemize}
	\item Imagine you want to store values for 500 integers.
	\item The best way to do this is with an array:
		\begin{verbatim}
int arr[500];
		\end{verbatim}

	\item There are several common ways to initialize all integers in an array to \texttt{0}:

		\begin{enumerate}
			\item Brace initialization (recommended)
			\begin{verbatim}
int arr[500] = {0};
			\end{verbatim}

			\item Empty brace initialization (C++11 and later)
			\begin{verbatim}
int arr[500] = {};
			\end{verbatim}

			\item Using \texttt{memset}
			\begin{verbatim}
int arr[500];
memset(arr, 0, sizeof(arr));
			\end{verbatim}

			\item Using \texttt{std::fill}
			\begin{verbatim}
int arr[500];
std::fill(arr, arr + 500, 0);
			\end{verbatim}

			\item Using a range-based \texttt{for} loop
			\begin{verbatim}
int arr[500];
for (int &x : arr)
    x = 0;
			\end{verbatim}

			\item Using a traditional \texttt{for} loop
			\begin{verbatim}
int arr[500];
for (int i = 0; i < 500; i++)
    arr[i] = 0;
			\end{verbatim}

    \item Partially initialize the first two elements
    \begin{verbatim}
int arr[5] = {10, 20};
// arr = {10, 20, 0, 0, 0}
    \end{verbatim}
		\end{enumerate}

	\item \textbf{Note:} \texttt{memset} should only be used to initialize an integer array to \texttt{0}. It should \emph{not} be used to initialize an array of integers to arbitrary values (e.g., \texttt{1}, \texttt{-1}, or \texttt{100}), since it operates on bytes rather than integer values.\\ 

	memset = sets raw memory bytes \\
int initialization = sets typed values\\ 

If you do \verb|memset(arr, -1, sizeof(arr))| --- this only “works as expected” in some cases because it produces all bits set:

0xFF 0xFF 0xFF 0xFF = -1 (two’s complement)\\

	\item Such \textbf{arrays} are called \vocab{static arrays} because the \textbf{number of elements} they contain as well as the \textbf{memory} the array consumes is \underline{\textbf{fixed}} at \textbf{compile-time}.
\end{itemize}

\newpage 

\subsubsection{How Data Is Stored in an Array}
\begin{itemize}
	\item Indices in C++ start with $\mathbf{0}$.
	\item Memory of the elements in the array are equally allocated.
	\item The amount of \textbf{memory} reserved by the compiler for the array of 5 ints is:
		\begin{verbatim}
sizeof(int) * 5
		\end{verbatim}

	\item In general, the amount of memory reserved by the compiler for an array in \textbf{bytes} is
		\begin{verbatim}
Bytes consumed by an array = sizeof(element-type) * Number of Elements
		\end{verbatim}
\end{itemize}

\newpage
\subsubsection{Accessing Data Stored in an Array}
\begin{itemize}
\item When you access an element at index $N$, the compiler \underline{does not} "search" for it.
		Instead, it computes the exact memory address using \textbf{arithmetic}.

	\item Arrays in C/C++ are stored in \textbf{contiguous memory}, meaning each element is placed directly next to the previous one.

	\item The compiler treats the \textbf{array name} as a \textbf{\underline{pointer} to the first element}.
	That is, the array name holds the memory address of index \textbf{0}.

	\item To access element $A[N]$, the compiler:
	\begin{itemize}
		\item Starts from the base address $A[0]$
		\item Adds an offset of $N \times \texttt{sizeof(element)}$
	\end{itemize}

	\begin{verbatim}
address of A[N] = base_address + N * sizeof(element)
	\end{verbatim}

	\item This is why arrays and pointers are closely related:
	\texttt{A} is essentially the address of \texttt{A[0]}, so:
	\begin{verbatim}
A[i] == *(A + i)
	\end{verbatim}

	\item In other words, pointer arithmetic is what makes array indexing possible.

	\item The C++ compiler \underline{does not} check if the index is within the actual defined bounds of the array.

	\item Accessing an array beyond its bounds results in unpredicatble behavior and raises security risks to your program like crashing and should be avoided at all costs.

	\item \fbox{\textbf{Listing 4.1} Arrays}
\end{itemize}

\newpage

\subsubsection{Modifying Data Stored in an Array}
\begin{itemize}
	\item \fbox{\textbf{Listing 4.2} Assigning Values to Elements in Array}
\end{itemize}

\newpage

\subsection{Multidimensional Arrays}
\begin{itemize}
	\item A multidimensional array is an array of arrays to comprise in total an \vocab{n-dimensional array}.
	\item The most common type is a \textbf{2D array}, which can be thought of as a table with rows and columns.

	\item Declaration example:
	\begin{verbatim}
int grid[3][4];
	\end{verbatim}
	This creates a 2D array with 3 rows and 4 columns.

	\item Initialization example:
	\begin{verbatim}
int grid[2][3] = {
    {1, 2, 3},
    {4, 5, 6}
};
	\end{verbatim}

	\item Accessing elements uses two indices:
	\begin{verbatim}
grid[row][col]
	\end{verbatim}

	\item Example:
	\begin{verbatim}
grid[0][1] = 2;
grid[1][2] = 6;
	\end{verbatim}

	\item Memory layout:
	Multidimensional arrays are stored in \textbf{row-major order}, meaning rows are stored contiguously in memory.

	\item Example layout:
	\begin{verbatim}
{1, 2, 3, 4, 5, 6}
	\end{verbatim}

	\item Pointer relationship:
	\begin{verbatim}
grid[i][j] == *(*(grid + i) + j)
	\end{verbatim}

	\item Common use cases:
	\begin{itemize}
		\item Matrices in math
		\item Grids (games, maps)
		\item Tables of data
	\end{itemize}

\item \fbox{\textbf{Listing 4.3} Multidimensional Array}
\end{itemize}

\newpage

\subsection{Dynamic Arrays}
\begin{itemize}
	\item A dynamic array is an array that can change size during runtime, unlike a normal C-style array which has a fixed size.

	\item In C++, the most common way to use dynamic arrays is through \textbf{vectors}.

	\item A vector is part of the Standard Template Library (STL) and automatically manages memory for you.

	\item Declaration example:
	\begin{verbatim}
#include <vector>

vector<int> nums;
	\end{verbatim}

	\item Adding elements dynamically:
	\begin{verbatim}
nums.push_back(10);
nums.push_back(20);
nums.push_back(30);
	\end{verbatim}

	\item Accessing elements (same as arrays):
	\begin{verbatim}
cout << nums[0];   // 10
cout << nums[1];   // 20
	\end{verbatim}

	\item Size of the vector:
	\begin{verbatim}
cout << nums.size();
	\end{verbatim}

	\item Example full usage:
	\begin{verbatim}
vector<int> nums = {1, 2, 3};

for (int i = 0; i < nums.size(); i++)
{
    cout << nums[i] << " ";
}
	\end{verbatim}

	\item Key advantages:
	\begin{itemize}
		\item Automatically resizes when elements are added
		\item No need to manually manage memory
		\item Safer than raw pointers and arrays
	\end{itemize}

	\item Internally, vectors use dynamic memory allocation and may reallocate to a larger block when capacity is exceeded.
\end{itemize}

Side Notes:

\subsection*{Fixed Arrays using \texttt{std::array}}
\begin{itemize}
	\item \texttt{std::array} is a container from the C++ Standard Library that represents a fixed-size array.

	\item It is defined in the header:
	\begin{verbatim}
#include <array>
	\end{verbatim}

	\item Unlike \texttt{std::vector}, the size of a \texttt{std::array} is fixed at compile time.

	\item Declaration example:
	\begin{verbatim}
std::array<int, 5> nums = {1, 2, 3, 4, 5};
	\end{verbatim}

	\item Accessing elements works like normal arrays:
	\begin{verbatim}
cout << nums[0];   // 1
cout << nums[3];   // 4
	\end{verbatim}

	\item Bounds-checked access can be done using \texttt{at()}:
	\begin{verbatim}
cout << nums.at(0);
	\end{verbatim}

	\item Getting the size:
	\begin{verbatim}
cout << nums.size();
	\end{verbatim}

	\item Iterating through a \texttt{std::array}:
	\begin{verbatim}
for (int i = 0; i < nums.size(); i++)
{
    cout << nums[i] << " ";
}
	\end{verbatim}

	\item Range-based loop (modern C++):
	\begin{verbatim}
for (int x : nums)
{
    cout << x << " ";
}
	\end{verbatim}

	\item Key advantages over raw arrays:
	\begin{itemize}
		\item Knows its own size (\texttt{.size()})
		\item Safer access with \texttt{.at()}
		\item Works well with STL algorithms
		\item No pointer decay unless explicitly used
	\end{itemize}

	\item Key difference from \texttt{std::vector}:
	\begin{itemize}
		\item \texttt{std::array}: fixed size
		\item \texttt{std::vector}: dynamic size
	\end{itemize}
\end{itemize}

\newpage

\subsection{C-style Character Strings}
\begin{itemize}
	\item \textbf{C-style strings} are a special case of an array of characters.

	\item We have seen examples of C-style strings in the form of \textbf{string literals}:
		\begin{verbatim}
std::cout << "Hellow World";
		\end{verbatim}

	\item This is equivalent to using the \textbf{array declaration}:
	\begin{verbatim}
char sayHello[] = {'H', 'e', 'l', 'l', 'o', ' ', 'W', 'o', 'r', 'l', 'd', '\0'};
std::cout << sayHello << std::endl;
	\end{verbatim}

	\item Equivalent way using a string literal (recommended):
\begin{verbatim}
char sayHello[] = "Hello World";
std::cout << sayHello << std::endl;
\end{verbatim}

\item Equivalent using a pointer to a string literal:
\begin{verbatim}
const char* sayHello = "Hello World";
std::cout << sayHello << std::endl;
\end{verbatim}

\item Equivalent using std::string (modern C++ strings):
\begin{verbatim}
#include <string>

std::string sayHello = "Hello World";
std::cout << sayHello << std::endl;
\end{verbatim}

\item \fbox{\textbf{Listing 4.6} Risky app of C-style strings}

\item If you forget the \vocab{null terminator} (\verb|\0|) or it gets overwritten, then certain functions that rely on this \textbf{null-terminator} will go beyond the intended data and can cause data security issues and potentially a crash.

\item All strings must end in \verb|\0|.

	\pbreak

\item Example: \verb|strlen()| computes the length of a string by walking the character buffer and counts the number of characters crossed until it reachs a \textbf{null-terminator} that indicates the \textbf{end of the string}.

\item This behavior of \verb|strlen()| makes it \textbf{dangerous} as it can easily walk past the bounds of the character array if this null-terminator isn't present.

\pbreak

\item C programs use these kind of functions (\verb|strlen()| (length of a string), \verb|strcpy()| (copy a string), \verb|strcat()| (concatenate a string), etc) that take in as input strings and can be dangerous as they seek the \textbf{null-terminator} and can exceed the boundaries potentially.


\end{itemize}

\newpage

Side note:
\subsubsection{Difference between \texttt{'\textbackslash n'} and \texttt{std::endl}}

\begin{itemize}
	\item Both \texttt{'\textbackslash n'} and \texttt{std::endl} are used to move output to a new line, but they are not the same.

	\subsubsection*{\texttt{'\textbackslash n'} (newline character)}
	\begin{itemize}
		\item Inserts only a newline character into the output stream.
		\item Does \textbf{not} flush the output buffer.
		\item More efficient and faster.
	\end{itemize}

	\begin{verbatim}
cout << "Hello\n";
	\end{verbatim}

	\subsubsection*{\texttt{std::endl}}
	\begin{itemize}
		\item Inserts a newline character.
		\item Forces a \textbf{flush} of the output buffer.
		\item Slower because flushing is an expensive operation.
	\end{itemize}

	\begin{verbatim}
cout << "Hello" << endl;
	\end{verbatim}

	\subsubsection*{Key difference}
	\begin{itemize}
		\item \texttt{'\textbackslash n'} only moves to the next line.
		\item \texttt{std::endl} moves to the next line \textbf{and flushes output immediately}.
	\end{itemize}

	\subsubsection*{Comparison table}
	\begin{center}
	\begin{tabular}{|c|c|c|}
	\hline
	Feature & \texttt{'\textbackslash n'} & \texttt{std::endl} \\
	\hline
	New line & Yes & Yes \\
	Flush buffer & No & Yes \\
	Performance & Faster & Slower \\
	\hline
	\end{tabular}
	\end{center}

	\subsubsection*{When to use}
	\begin{itemize}
		\item Use \texttt{'\textbackslash n'} for normal output (recommended most of the time).
		\item Use \texttt{std::endl} when you need immediate output (e.g., debugging).
	\end{itemize}
\end{itemize}

\newpage

\subsection{C++ Strings: std::string}

\begin{itemize}
	\item C++ standard strings are an efficient and \textbf{safer} way to deal with text input.
	\item \verb|std::string| is \underline{\textbf{not static}} in \textbf{size} like a \textbf{char array} implementation of a \textbf{C-style string} is.
\item \verb|std::string| can \textbf{dynamically resize} when more data is added, unlike fixed-size C-style character arrays. It is only slightly slower in theory due to overhead, but in practice and with compiler optimizations, it's minute (small difference in speed).
\item \fbox{\textbf{Listing 4.7} using std::string}
\end{itemize}

\newpage 

\subsection{QA}

\begin{itemize}
	\item \textbf{Q: Why take the trouble to initialize a static array's elements?}

		A: Unless initialized, the array, unlike a variable of any other tpye, contains \textbf{junk} and \textbf{unpredictable values} as the \textbf{memory} at that location was left untouched after the last operations.

		Initializing ensures a predictable initial state.\\

	\item \textbf{Q: Would you need to initialize the elements n a dynamic array for the same principles?}

	A: No, not necessarily. A \textbf{dynamic array} in cpp is smart. Elements are typically defaulted to a value of 0 if left uninitialized.\\

\item \textbf{Q: Does the length of the string \underline{include} the null-terminator}?

	A: No, it doesn't. The length of the string "Hello World" is 11, including the space character in btw the words and \underline{\textbf{excluding}} the \textbf{null-terminator}.\\

\item \textbf{Q: If i want to still use C-style arrays, what should be its size?}

	A: The size should be the length of the string + the null-terminator. So, basically account for a + 1 to the size when allocating the array.

	"Hello World" should be allocated in an array of size 11 + 1 = 12.
\end{itemize}

\newpage

\subsection{Quiz}
\begin{enumerate}
	\item If you need to allow the user to input strings, would you use C-style strings?

	A: No. As shown, the input can be greater than your bounds of the statically allocated C-style string.

	This will lead to undefined behavior if not checked.\\

\item How many characters are in \verb|'\0'| as seen by the compiler?

	A: One character: the null-terminator.
\end{enumerate}

\newbp

\section{Working with Expressions, Statements, and Operators}
At its heart, a program is just a sequence of steps whether it be sequentially or in parallel.\\
These steps are programmed into expressions, statements, and use operators to perform calculations/actions.

\subsection{Statements}
\begin{itemize}
	\item Languages --- spoken or programmed --- are made up of \vocab{statements} that are executed one after another.

	\item Example:
		\begin{verbatim}
std::cout << "Hello World" << std::endl;
		\end{verbatim}

	\item All statements in C++ end in a \textbf{;}
	\item Extra spaces between statements \underline{\textbf{do not}} matter (some languages it does matter like Python)
\item In the front end of a compiler, the \textbf{lexical analyzer (lexer)} typically ignores whitespace that separates tokens. For example, the following declarations are treated identically:

\begin{verbatim}
char c;
char      c;
char
c;
char     


c;
\end{verbatim}

\item The lexer produces the same sequence of tokens for each:

\begin{verbatim}
KEYWORD(char)  IDENTIFIER(c)  SEMICOLON(;)
\end{verbatim}

\item Whitespace is only significant in places where it separates tokens. Outside of string and character literals, extra spaces, tabs, and newlines generally do not change the meaning of a C++ program.

\item Writing whitespaces within a statement can effect certain things like for example string literals:

	\begin{verbatim}
cout << "Hello
	World" << endl; // is invalid
	\end{verbatim}

\item To solve the above we can either do
	\begin{verbatim}
cout << "Hello \
	World" << endl; // split ot two lines is OK
	\end{verbatim}

\item Or this but entails a concatenation by the compiler
		\begin{verbatim}
cout << "Hello "
	"World" << endl; // two string literals is also OK^{*} 
					// (compiler does extra work --- catenation)
	\end{verbatim}

\item The above one with \verb|\| is best to split up a long statement into multiple lines.

\end{itemize}

\newpage 

\subsection{Compound Statements or Blocks}
\begin{itemize}
	\item When you group \textbf{statements} together within \textbf{braces} \verb|{...}|, you create a \vocab{compound statement} or \vocab{block}.

	\begin{verbatim}
{
	int daysInYear = 365;
	cout << "Block contains an int and a cout statement" << endl;
}
	\end{verbatim}

\item A \vocab{block} typically groups many statements to indicate that they belong together.

\item We've seen it already through functions, constexpr's, arthimetic operations, and so on with conditionals, loops, etc
\end{itemize}

\newpage

\subsection{Using Operators}
\begin{itemize}
	\item \vocab{Operators} are tools for you to be able to work with data, transform it, process it, etc.
\end{itemize}

\pbreak

\subsubsection{Assignment Operator (=)}
\begin{itemize}
	\item We've been using this
	\begin{verbatim}
int daysInYear = 365;
	\end{verbatim}

\item The \textbf{assignment operator} replaces the value contained by the \vocab{operand} on the left (called \vocab{l-value}) by that on the right (called \vocab{r-value}).
\end{itemize}

\pbreak

\subsubsection{L-vales and R-values}
\begin{itemize}
	\item \vocab{L-values} often refer to \textbf{locations in memory}.

	\item A \textbf{variable} such as \verb|daysInYear| from the preceding example is actually a \textbf{\underline{handle}} to a \textbf{memory location} and is an \vocab{l-value}.

	\item \vocab{R-values}, on the other hand, can be the very \underline{\textbf{content}} of a \textbf{memory location}.

		\pbreak

	\item an l-value = a reference/handle to a memory location (something you can assign to, like a variable)
	\item an r-value = the value contained in that location (the data itself)

	\pbreak

\item So, all \textbf{l-values} can be used as \textbf{r-values}, but not all \textbf{r-values} can be used as \textbf{l-values}.
\begin{verbatim}
Example:
    daysInYear = 365;   // valid (l-value = r-value)

    365 = daysInYear;   // invalid (r-value cannot be assigned to)
\end{verbatim}

\item So, strictly speaking: an r-value \underline{\textbf{cannot}} be an l-value.

\end{itemize}

\newpage

\subsubsection{Operators to Add (+), Subtract(-), Multiply(* or $<<$), Divide(/ or $>>$), and Modulo(\%)}
\begin{itemize}
	\item These are basic arithmetic operators used to perform mathematical operations on numeric values.

	\item \textbf{Addition (+)}: Computes the sum of two operands.

	\item \textbf{Subtraction (-)}: Computes the difference between two operands.

	\item \textbf{Multiplication (* or $<<$)}: Multiplies two values; in some contexts, \verb|<<| is used for bitwise left shift, which effectively multiplies by powers of 2.

	\item \textbf{Division (/ or $>>$)}: Divides one value by another; \verb|>>| is a bitwise right shift, often corresponding to integer division by powers of 2.

	\item \textbf{Modulo (\%)}: Returns the remainder after division of two integers.

	\pbreak

	\item \textbf{Division (/)}: Performs division between two operands. If both operands are integers, the result is also an integer and the fractional part is truncated (i.e., decimal portion is discarded, effectively rounding toward zero). If at least one operand is a \verb|float| or \verb|double|, floating-point division is performed and the decimal part is preserved.

\begin{verbatim}
Key clarification (division behavior in C++)

7 / 2 = 3              (integer / integer → truncates toward zero)
7 / 2.0 = 3.5          (integer / double → floating-point division)
7.0 / 2 = 3.5          (double / integer → floating-point division)
7.0 / 2.0 = 3.5        (double / double → floating-point division)

-7 / 2 = -3            (truncates toward zero, NOT -4)
-7 / 2.0 = -3.5        (floating-point division)
-7.0 / 2 = -3.5        (floating-point division)

int x = 7 / 2;         // x = 3  (result truncated before assignment)
double y = 7 / 2;      // y = 3.0 (division happens first, still integer division, IMPORTANT)
double z = 7 / 2.0;    // z = 3.5 (floating-point division) 
\end{verbatim}
\end{itemize}

\newpage 

\subsubsection{Operators to Increment(++) and Decrement(--)}

\begin{itemize}
	\item The increment operator \textbf{(++)} increases a variable’s value by 1, and the decrement operator \textbf{(--)} decreases it by 1. These operators are shorthand for \verb|x = x + 1| and \verb|x = x - 1|.

	\item \textbf{Prefix form (++x, $--$x)}: The variable is modified first, and the updated value is then used in the expression.
	
	\item \textbf{Postfix form (x++, x$--$)}: The current value is used first, and then the variable is modified.

	\item Example:
	\begin{verbatim}
int x = 5;
int a = ++x;  // x becomes 6, a = 6
int b = x++;  // b = 6, then x becomes 7
	\end{verbatim}

	\item \textbf{Pointers and increment/decrement}: When used on pointers, \verb|++| and \verb|--| move the pointer to the next or previous memory location of its type (not just +1 byte).

	\item Example:
	\begin{verbatim}
int arr[] = {10, 20, 30};
int* p = arr;

p++;  // moves to arr[1]
p--;  // moves back to arr[0]
	\end{verbatim}

	\item \textbf{Iterators in loops}: In C++ iterators behave like pointers, so \verb|++it| moves to the next element. This is commonly used in \verb|for| loops to traverse containers.

	\item Example:
	\begin{verbatim}
for (vector<int>::iterator it = v.begin(); it != v.end(); ++it)
{
	cout << *it << endl;
}
	\end{verbatim}

	\pbreak

\item \textbf{Prefix} is perferred over \textbf{postfix} due to, at least theoretically, with the \textbf{postfix operators}, the compiler needs to store the initial alue temporarily in the event of it needing to be assigned.

	The effect on performance in these cases is neglible with respect to integers, but in the case of certain classes there might be a point in this argument.

	\vocab{Smart pointers} may optimize away the differences.
\end{itemize}

\newpage

\subsubsection{Equality Operators (==) and (!=)}
\begin{itemize}
	\item (==) checks if two \textbf{operands} are equal.
	\item (!=) checks if two \textbf{operands} are not equal.

	\item The behavior of relational and equality operators can depend on the \textbf{types being compared}. For example, comparing integers vs floating-point numbers may involve implicit type conversion, and comparing pointers checks addresses rather than values.

	\item Example: integer vs floating-point comparison
	\begin{verbatim}
int a = 5;
double b = 5.0;

cout << (a == b) << endl;  // true (a is promoted to double)
	\end{verbatim}

	\item Example: floating-point precision issues
	\begin{verbatim}
double x = 0.1 + 0.2;
double y = 0.3;

cout << (x == y) << endl;  // false (precision error in floating-point math)
	\end{verbatim}

	\begin{verbatim}
Why (0.1 + 0.2) != 0.3 in C++

Even though mathematically:
    0.1 + 0.2 = 0.3

In floating-point (binary) representation:
    0.1 is stored as an approximation
    0.2 is stored as an approximation
    0.3 is also stored as an approximation

So:
    (0.1 + 0.2) produces a slightly different value than 0.3

Example result:
    x = 0.1 + 0.2  ->  0.3000000000000000444
    y = 0.3        ->  0.2999999999999999889

Therefore:
    (x == y) is false due to tiny precision error
\end{verbatim}

	\item Example: pointer comparison (compares addresses, not values)
	\begin{verbatim}
int a = 10, b = 10;

int* p1 = &a;
int* p2 = &b;

cout << (p1 == p2) << endl;  // false (different memory locations)
cout << (a == b) << endl;    // true (values are equal)
	\end{verbatim}

	\item Example: string comparison (C-style strings compare addresses, not content)
	\begin{verbatim}
const char* s1 = "hello";
const char* s2 = "hello";

cout << (s1 == s2) << endl;  // may be true or false (implementation-dependent)
	\end{verbatim}
\end{itemize}

\newpage

\subsubsection{Relational Operators}
\begin{itemize}
	\item Relational operators are used to compare two operands and determine the relationship between them. The result of any relational expression is a \textbf{boolean value} (true or false).

	\item \textbf{Less than ($<$)}: Checks if the left operand is smaller than the right operand.

	\item \textbf{Greater than $(>$)}: Checks if the left operand is larger than the right operand.

	\item \textbf{Less than or equal to $(<=$)}: Checks if the left operand is smaller than or equal to the right operand.

	\item \textbf{Greater than or equal to ($>=$)}: Checks if the left operand is larger than or equal to the right operand.

	\item Relational operators are commonly used in conditions such as \verb|if| statements and loops.

	\item Example:
	\begin{verbatim}
int a = 5, b = 10;

cout << (a < b) << endl;   // true
cout << (a > b) << endl;   // false
cout << (a <= 5) << endl;  // true
cout << (b >= 20) << endl; // false
	\end{verbatim}

	\item Relational comparisons can behave differently depending on types. For example, comparing integers is exact, while floating-point comparisons may be affected by precision errors as demostrated in the previous section.

\item In C++, relational and logical expressions evaluate to a \textbf{boolean value} (\verb|true| or \verb|false|). When printed using \verb|cout|, \verb|false| is displayed as \verb|0| and \verb|true| is displayed as \verb|1|. 

	\textbf{\underline{Any non-zero}} value is treated as \verb|true| in conditional expressions.

	The compiler checks the condition against zero. So, if it's not zero, then the condition is marked true.

\end{itemize}

\newpage

\subsubsection{Logical Operations NOT, AND, OR, and XOR}
\begin{itemize}
	\item Logical operators are used to combine or modify boolean expressions. In C++, any expression evaluates to \verb|true| (1) or \verb|false| (0) when used in a logical context.

	\item \textbf{NOT (!) Operator}: Reverses the boolean value.
	\begin{verbatim}
Operand: false
Result : !false = true (1)

Operand: true
Result : !true = false (0)
	\end{verbatim}

	\item \textbf{AND (\&\&) Operator}: Returns true only if both operands are true.
	\begin{verbatim}
Operand 1: false   Operand 2: true
Result   : false && true = false (0)

Operand 1: true    Operand 2: true
Result   : true && true = true (1)
	\end{verbatim}

	\item \textbf{OR ($||$) Operator}: Returns true if at least one operand is true.
	\begin{verbatim}
Operand 1: false   Operand 2: true
Result   : false || true = true (1)

Operand 1: false   Operand 2: false
Result   : false || false = false (0)
	\end{verbatim}

	\item \textbf{XOR concept (not a built-in logical operator in C++)}: Exclusive OR is true only when operands are different. In C++, it is implemented using bitwise XOR \verb|^| for integers.
	\begin{verbatim}
Operand 1: false (0)   Operand 2: true (1)
Result   : 0 ^ 1 = 1 (true)

Operand 1: true (1)    Operand 2: true (1)
Result   : 1 ^ 1 = 0 (false)
	\end{verbatim}

	\item \textbf{Key idea}: In C++, \verb|false| is represented as 0 and \verb|true| as 1 when printed, but internally they are boolean values, not integers.

	\item \textbf{Short-circuit behavior}:
	\begin{itemize}
		\item \verb|&&| stops evaluating if the first operand is false.
		\item \verb+||+ stops evaluating if the first operand is true.
	\end{itemize}
\end{itemize}

\newpage

\subsubsection{Bitwise Operators: NOT (\texttt{\~}), AND (\&), OR (\texttt{|}), XOR (\texttt{\^})}

\begin{itemize}
	\item Bitwise operators manipulate the individual \textbf{bits} of integral data types (such as \verb|char|, \verb|int|, and \verb|long|). They operate on the binary representation of values rather than their decimal values.

	\item \textbf{Bitwise NOT (\texttt{\~})}: Flips every bit in the operand (0 becomes 1, and 1 becomes 0).

	\item Example:
	\begin{verbatim}
unsigned char x = 10;      // 00001010
unsigned char y = ~x;      // 11110101
	\end{verbatim}

	\item \textbf{Bitwise AND (\&)}: Produces a 1 in each bit position only if the corresponding bits of both operands are 1.

	\item Example:
	\begin{verbatim}
12 & 10

1100
1010
----
1000    // 8
	\end{verbatim}

	\item \textbf{Bitwise OR (\texttt{|})}: Produces a 1 in each bit position if either corresponding bit is 1.

	\item Example:
	\begin{verbatim}
12 | 10

1100
1010
----
1110    // 14
	\end{verbatim}

	\item \textbf{Bitwise XOR (\texttt{\string\^})}: Produces a 1 in each bit position only when the corresponding bits are different.

	\item Example:
	\begin{verbatim}
12 ^ 10

1100
1010
----
0110    // 6
	\end{verbatim}

	\item Bitwise operators are commonly used for bit masking, setting or clearing flags, toggling bits, low-level hardware programming, and performance-sensitive code.
\end{itemize}

\newpage

\subsubsection{Bitwise right shift $>>$ and left shift $<<$}

\begin{itemize}
	\item Bitwise shift operators move the bits of an integral value to the left or right by a specified number of positions.

	\item \textbf{Left Shift (\texttt{<<})}: Shifts all bits to the left. Each left shift by one position is approximately equivalent to multiplying the value by 2, provided no significant bits are discarded (overflow).

	\item Example:
	\begin{verbatim}
5 << 1

00000101   (5)
<< 1
--------
00001010   (10)

5 << 2

00000101   (5)
<< 2
--------
00010100   (20)
	\end{verbatim}

	\item \textbf{Right Shift (\texttt{>>})}: Shifts all bits to the right. Each right shift by one position is approximately equivalent to integer division by 2, discarding any remainder.

	\item Example:
	\begin{verbatim}
20 >> 1

00010100   (20)
>> 1
--------
00001010   (10)

20 >> 2

00010100   (20)
>> 2
--------
00000101   (5)
	\end{verbatim}

	\item Bits shifted out of the value are discarded. For a left shift, zeros are shifted into the least significant bits. For unsigned integers, a right shift shifts zeros into the most significant bits.

	\item Example:
	\begin{verbatim}
1 << 5

00000001
<< 5
--------
00100000   (32)
	\end{verbatim}

	\item Bitwise shift operators are commonly used for efficient multiplication or division by powers of two, bit masking, manipulating binary data, and low-level systems programming.

	\item \textbf{Note}: Although left and right shifts can often be thought of as multiplication or division by powers of two, this relationship is only reliable when no overflow occurs (for left shifts) and when shifting unsigned integers.
\end{itemize}

\newpage

\subsubsection{Compound Assignment Operators}

\begin{itemize}
	\item Compound assignment operators combine an arithmetic or bitwise operation with assignment. They provide a shorthand for updating the value of a variable.

	\item \textbf{Addition Assignment (\texttt{+=})}: Adds the right operand to the left operand, then stores the result in the left operand.
	\begin{verbatim}
x += 5;    // Equivalent to: x = x + 5;
	\end{verbatim}

	\item \textbf{Subtraction Assignment (\texttt{-=})}: Subtracts the right operand from the left operand, then stores the result.
	\begin{verbatim}
x -= 5;    // Equivalent to: x = x - 5;
	\end{verbatim}

	\item \textbf{Multiplication Assignment (\texttt{*=})}: Multiplies the left operand by the right operand, then stores the result.
	\begin{verbatim}
x *= 5;    // Equivalent to: x = x * 5;
	\end{verbatim}

	\item \textbf{Division Assignment (\texttt{/=})}: Divides the left operand by the right operand, then stores the result.
	\begin{verbatim}
x /= 5;    // Equivalent to: x = x / 5;
	\end{verbatim}

	\item \textbf{Modulo Assignment (\texttt{\%=})}: Computes the remainder after division and stores the result.
	\begin{verbatim}
x %= 5;    // Equivalent to: x = x % 5;
	\end{verbatim}

	\item \textbf{Bitwise AND Assignment (\texttt{\&=})}: Performs a bitwise AND with the right operand and stores the result.
	\begin{verbatim}
x &= y;    // Equivalent to: x = x & y;
	\end{verbatim}

	\item \textbf{Bitwise OR Assignment (\texttt{|=})}: Performs a bitwise OR with the right operand and stores the result.
	\begin{verbatim}
x |= y;    // Equivalent to: x = x | y;
	\end{verbatim}

	\item \textbf{Bitwise XOR Assignment (\texttt{\string^=})}: Performs a bitwise XOR with the right operand and stores the result.
	\begin{verbatim}
x ^= y;    // Equivalent to: x = x ^ y;
	\end{verbatim}

	\item \textbf{Left Shift Assignment (\texttt{<<=})}: Shifts the bits of the left operand to the left and stores the result.
	\begin{verbatim}
x <<= 2;   // Equivalent to: x = x << 2;
	\end{verbatim}

	\item \textbf{Right Shift Assignment (\texttt{>>=})}: Shifts the bits of the left operand to the right and stores the result.
	\begin{verbatim}
x >>= 2;   // Equivalent to: x = x >> 2;
	\end{verbatim}

	\item Compound assignment operators are commonly used to write shorter, more readable code while avoiding repetition of the left-hand operand.
\end{itemize}

\newpage

\subsubsection{Using sizeof Operator}
\begin{itemize}
	\item \verb|sizeof| operator tells you the \textbf{amount of memory} in \textbf{\underline{bytes}} consumed by a particular type or a variable.

		\begin{verbatim}
sizeof(variable);
		or
sizeof(type);
		\end{verbatim}

	\item \verb|sizeof()| may look like a function, but it is \textbf{\underline{not a function}}; it's an \textbf{operator}.
	\item Interestingly, this operator \underline{\textbf{cannot be overloaded}}, i.e. \textbf{cannot} be defined by the programmer.
\end{itemize}

\newpage

\subsubsection{Operator Precedence}
\begin{enumerate}
	\item Scope resolution
	\item Member selection, subscripting, increment, decrement
	\item sizeof, prefix increment and decrement, complement, and, not, unary minus and plus, address-of and dereference, new, new[], delete, delete[], ccasting, sizeof()

	\item Member selection for pointer
	\item Multiply, divide, modulo
	\item Add, subtract
	\item Shift (shift left, shift right)
	\item Inequality relational
	\item Equality, inequality
	\item Bitwise AND
	\item Bitwise exclusive OR
	\item Bitwise OR
	\item Logical AND
	\item Logical OR
	\item Conditional
	\item Assignment operators
	\item Comma
\end{enumerate}

\newpage

\subsection{QA}
\begin{itemize}
	\item Trivial questions
\end{itemize}

\newpage

\subsection{Quiz}
\begin{enumerate}
	\item I am writing an app to divide numbers. What's a better suited data type: int or float?

	A: I would say float personally, cause it wouldn't truncate the decimals.
\end{enumerate}

\newbp
\section{Controlling Program Flow}
\begin{itemize}
	\item \textbf{Switch statements} instead of if conditions
	\item The \verb|switch (variable)| must match the \verb|case variable_type:|
	\item If you dont do a \verb|break;| in the cases, you \textbf{fall through} to the next below case.
\begin{verbatim}
#include <iostream>

using namespace std;

int main() {
    int choice;

    cout << "Enter a number (1-3): ";
    cin >> choice;

    switch (choice) {
        case 1:
            cout << "You selected option 1" << endl;
            break;

        case 2:
            cout << "You selected option 2" << endl;
            break;

        case 3:
            cout << "You selected option 3" << endl;
            break;

        default:
            cout << "Invalid option" << endl;
            break;
    }

    return 0;
}
\end{verbatim}
\end{itemize}

\newpage
\section{Organizing Code with Functions}
Don't need this personally.

\newbp
\section{Pointers and References}
\subsection{What is a Pointer?}
\begin{itemize}
	\item A \vocab{pointer} is also a \textbf{variable} --- one that \textbf{stores} an \underline{\textbf{address}} in \textbf{memory}.
	\item Just the same way as a \textbf{variable} of type \textbf{int} is used to contain an \textbf{integer value}, a \fbox{\textbf{pointer variable}} is used to contain a \textbf{memory address}.

	\item Thus, a \textbf{pointer} is a \textbf{variable}, and like all variables a \textbf{pointer} occupies \textbf{space in memory}.

	\item The \textbf{value} contained in a \textbf{pointer} is interpreted as a \underline{\textbf{memory address}}.
	\item Again, a \textbf{pointer} is a special variable that \textit{points} to a \textbf{location in memory}.

	\pbreak

\item Note: \textbf{memory addresses} are in \vocab{hexadecimal (base 16)} and start with an \verb|0x|
\end{itemize}

\newpage

\subsubsection{Declaring a Pointer}
\begin{itemize}
	\item Just like a variable, a pointer needs to be declared.
	\item You normally declare a pointer to point to a specific \textbf{type} (ex: int*)
	\item This would mean that the \textbf{address} contained in the \textbf{pointer} points to a \textbf{location in the memory} that holds an integer.
\item A \vocab{void pointer} (\verb|void*|) is a pointer that can store the address of any data type. It is often used to represent a generic pointer to raw memory. However, since it has no type information, it cannot be dereferenced directly and must first be cast to a specific pointer type.

	\begin{verbatim}
#include <iostream>
using namespace std;

int main()
{
    int x = 10;
    double y = 3.14;

    void* ptr;

    ptr = &x;   // void pointer holds address of int
    cout << *(static_cast<int*>(ptr)) << endl;

    ptr = &y;   // now holds address of double
    cout << *(static_cast<double*>(ptr)) << endl;

    return 0;
}
	\end{verbatim}


	\pbreak

\item As in the case with most variables, unless you initialize a pointer it will contain a random value.
\item To avoid accessing a \textbf{random (uninitialized) memory address}, pointers should be initialized to \textbf{NULL} (C style) or \textbf{nullptr} (C++11 and later). This ensures the pointer is explicitly set to point to nothing until it is assigned a valid address.

	\item \textbf{nullptr} is preferred in modern C++ because it is type-safe, whereas \textbf{NULL} is typically just defined as 0 and may cause ambiguity in overloaded function calls.
	
	\item Thus, declaring a pointer to an integer would be

		\begin{verbatim}
int *pointsToInt = NULL; // NULL lol
		\end{verbatim}

\end{itemize}

\newpage

\subsubsection{Determing the Address of a Variable Using the Reference Operator (\&)}
\begin{itemize}
	\item if \verb|varName| is a variable, \verb|&varName| gives the \textbf{address in memory} where its value is placed.

	\item So, if you have declared an integer \verb|int age = 30;|, \verb|&age| would be the \textbf{address in memory} where the value 30 is placed.

	\item \fbox{\textbf{Listing 8.1} Addresses} VERY IMPORTANT
\end{itemize}

\newpage

\subsubsection{Access Pointer Data Using the Dereference Operator (*)}

\begin{itemize}
	\item The dereference operator (\texttt{*}) is used to access the value stored at the memory address held by a pointer.

	\item Instead of working with the address itself, dereferencing allows you to work with the \textbf{actual data stored at that address}.

	\item Example:
	\begin{verbatim}
int x = 10;
int *p = &x;

cout << p << endl;   // prints address of x
cout << *p << endl;  // prints value of x (10)
	\end{verbatim}

	\item The expression \verb|*p| means ``go to the memory location stored in p and access its value.''

	\item If the value at that address is changed through the pointer, the original variable is also modified.

	\item Example:
	\begin{verbatim}
int x = 10;
int *p = &x;

*p = 20;   // modifies x through the pointer

cout << x << endl;   // outputs 20
	\end{verbatim}

	\item Dereferencing only works on valid (non-null) pointers. Using \verb|*| on an uninitialized or null pointer leads to undefined behavior.

		\pbreak

	\item For example, just look at \textbf{\fbox{Listing 8.1}}

	\pbreak

\item When the \vocab{dereference operator (*)} is used, the program uses the \textbf{address stored in the pointer} to access the \textbf{object located at that memory address}. The amount of memory accessed depends on the \textbf{data type of the pointer} (e.g., \verb|int|, \verb|double|), as determined by \verb|sizeof(type)| on the given system.


\item The \textbf{dereference operator} is also called the \vocab{indirection operator} (never heard that in my life)

\item Basically, the \textbf{memory address} stored in a \textbf{pointer} is very important and can cause issues if the address is invalid.

\end{itemize}

\newpage

\subsubsection{What is sizeof() a Pointer?}
\begin{itemize}
	\item The \textbf{pointer} is just another \textbf{variable} that contains a \textbf{memory address}.
	\item Hence, \underline{\textbf{irrespsecitve}} of the \textbf{type} that is being pointed to, the content of a \textbf{pointer} is an \textbf{address} --- a number.

	\item The \textbf{length of an address}, that is the \textbf{number of bytes} required to store it, is a \textbf{constant} for a given system.

	\item The \verb|sizeof()| a \textbf{pointer} is hence \underline{\textbf{dependent}} on the \textbf{compiler} and the \textbf{operating system} the program has been compiled for and is \underline{\textbf{not dependent}} on the nature of the \textbf{data} being pointed to.

	\item \fbox{\textbf{Listing 8.6}}

	\item The output of the \verb|sizeof(int*)| would be 
		\begin{itemize}
			\item \textbf{4 bytes} on a \textbf{32-bit compiler} and 
			\item \textbf{8 bytes} on a \textbf{64-bit compiler}. 
		\end{itemize}
\end{itemize}

\newpage
\subsection{Dynamic Memory Allocation}

\vocab{Dynamic memory allocation} allows a program to \textbf{request memory at runtime} from the \textbf{heap} instead of the \textbf{stack}.\\
This is useful when the \textbf{required memory size} is \textbf{not known} at \textbf{compile time}.\\ 

The \verb|sizeof()| operator is very useful here.

\pbreak

\subsubsection{Using new and delete or delete[]}

\begin{itemize}
	\item In C++, dynamic memory is allocated using the \textbf{new} operator and released using \textbf{delete}.
	\item For arrays, \textbf{new[]} and \textbf{delete[]} must be used.

	\item Example:
	\begin{verbatim}
int *p = new int;   // allocate single integer
*p = 10;

cout << *p << endl;

delete p;           // free memory
	\end{verbatim}

	\item Example (arrays):
	\begin{verbatim}
int *arr = new int[5];

for (int i = 0; i < 5; i++)
{
    arr[i] = i * 10;
}

for (int i = 0; i < 5; i++)
{
    cout << arr[i] << " ";
}

delete[] arr;   // free array memory
	\end{verbatim}

	\item \textbf{Important:} Using \textbf{delete} on memory allocated with \textbf{new[]} (or vice versa) leads to undefined behavior.
\end{itemize}

\newpage
\subsubsection{Using malloc, realloc, etc and free}

\begin{itemize}
	\item In C (also available in C++), dynamic memory is managed using \textbf{malloc}, \textbf{calloc}, \textbf{realloc}, and \textbf{free} from the \verb|<cstdlib>| library.

	\item \textbf{malloc}: Allocates a block of memory (uninitialized).
	\begin{verbatim}
int *p = (int*) malloc(5 * sizeof(int));

for (int i = 0; i < 5; i++)
{
    p[i] = i + 1;
}

free(p);
	\end{verbatim}

	\item \textbf{calloc}: Allocates memory and initializes it to zero.
	\begin{verbatim}
int *p = (int*) calloc(5, sizeof(int));
	\end{verbatim}

	\item \textbf{realloc}: Resizes previously allocated memory block.
	\begin{verbatim}
p = (int*) realloc(p, 10 * sizeof(int));
	\end{verbatim}

	\item \textbf{free}: Releases memory allocated by malloc/calloc/realloc.
	\begin{verbatim}
free(p);
	\end{verbatim}

	\item \textbf{Important difference}: Unlike \textbf{new/delete}, these functions do not call constructors or destructors in C++ and should be used carefully in modern C++ code.
\end{itemize}

\newpage

Important side note:

\begin{itemize}
	\item In C-style dynamic memory allocation, functions like \verb|malloc|, \verb|calloc|, and \verb|realloc| return a \vocab{void pointer (void*)}, which is a generic pointer type that can hold the address of any data type.

	\item A \verb|void*| is used because these functions do not know in advance what type of data the allocated memory will store. Instead, they simply return a raw memory address from the heap.

	\item Example:
	\begin{verbatim}
	int *p = (int*) malloc(5 * sizeof(int));
	\end{verbatim}

	\item The cast \verb|(int*)| is required in C++ because a \verb|void*| cannot be automatically converted into a typed pointer without an explicit conversion.

	\item \textbf{Why malloc returns void*:}
	\begin{itemize}
		\item It allows a single allocation function to be \textbf{generic} and work with any data type.
		\item It represents raw, untyped memory returned from the heap.
		\item The programmer decides how to interpret that memory by casting it to the appropriate pointer type.
	\end{itemize}

	\item Example showing interpretation:
	\begin{verbatim}
	void* ptr = malloc(4 * sizeof(int));

	int* iptr = (int*) ptr;
	iptr[0] = 10;
	\end{verbatim}

	\item \textbf{Important:} Unlike \textbf{new}, these functions do not perform type construction or initialization—they only allocate raw memory.

	\pbreak

	\item In C++, it is recommended to use \verb|static_cast| instead of C-style casts when converting a \verb|void*| to a typed pointer.

\item Example:
\begin{verbatim}
void* ptr = malloc(4 * sizeof(int));

int* iptr = static_cast<int*>(ptr);

// or really just

int* cppiptr = new int;

iptr[0] = 10;
\end{verbatim}

\item The expression \verb|(int*) ptr| is a C-style cast, while \verb|static_cast<int*>(ptr)| is the safer and more explicit C++-style cast.

\item \textbf{Note:} The cast is required in C++ because a \verb|void*| cannot be implicitly converted to another pointer type.


\end{itemize}

\newpage

Additional important side notes:
\begin{itemize}
\item The \textbf{new operator} in C++ can also be \textbf{overloaded} for custom memory management in user-defined classes. This allows a class to define how memory is allocated when \verb|new| is used.

\item Example:
\begin{verbatim}
class A {
public:
    void* operator new(size_t size);
};

void* A::operator new(size_t size)
{
    cout << "Custom new called" << endl;
    return malloc(size);
}
\end{verbatim}

\item In this example, the class \verb|A| defines its own version of the \verb|new| operator. The parameter \verb|size_t size| represents the number of bytes required to allocate an object of type \verb|A|.

\item \textbf{Note:} Although the operator returns \verb|void*| internally in the declaration above, in practice it is typically expected to return a pointer of type \verb|A*| after allocation.
\end{itemize}

\newpage

\begin{itemize}
\item \fbox{\textbf{Listing 8.7} new and delete with operator *}
\item \textbf{Operator \textit{delete}} canot be invoked on any address containing a pointer, rather only those that have been returned by the \textbf{operator \textit{new}} and \textbf{only} those that have not already been released by a \textit{delete}.

\item \fbox{\textbf{Listing 8.8}}
\end{itemize}

\newpage

Side note:
\subsubsection{Operators vs Functions vs constexpr}

\begin{itemize}
	\item In C++, computation can be expressed using \textbf{operators}, \textbf{functions}, and \textbf{constexpr expressions}. Each serves a different purpose in how code is written and evaluated.

	\item \textbf{Operators}: Special symbols that perform built-in operations on operands. They are concise and have fixed precedence and associativity rules.
	\begin{verbatim}
int x = 5 + 3;   // + is an operator
	\end{verbatim}

	\item \textbf{Functions}: Named blocks of reusable code that perform a task. They are called explicitly and can contain complex logic.
	\begin{verbatim}
int add(int a, int b)
{
    return a + b;
}

int x = add(5, 3);
	\end{verbatim}

	\item \textbf{constexpr}: A specifier that tells the compiler that a value or function can be evaluated at \textbf{compile time} if given constant inputs. This improves performance and enables compile-time computation.

	\item Example of constexpr variable:
	\begin{verbatim}
constexpr int x = 5 + 3;   // evaluated at compile time
	\end{verbatim}

	\item Example of constexpr function:
	\begin{verbatim}
constexpr int square(int x)
{
    return x * x;
}

int y = square(4);   // may be computed at compile time
	\end{verbatim}

	\item \textbf{Key differences}:
	\begin{itemize}
		\item Operators: built-in symbolic operations with fixed rules
		\item Functions: runtime or compile-time logic (depending on usage)
		\item constexpr: guarantees compile-time evaluation when possible
	\end{itemize}

	\item \textbf{Summary}:
	\begin{itemize}
		\item Operators → concise symbolic computation
		\item Functions → reusable procedural logic
		\item constexpr → compile-time computable values/functions
	\end{itemize}

	\pbreak

	\item Operators and functions are handled differently by the compiler. Function calls use an explicit argument list in parentheses, while operators are part of the C++ language grammar and are parsed as expressions rather than function calls.

\item For example, a function call is parsed as:
\begin{verbatim}
add(5, 3);
\end{verbatim}

\item In contrast, an operator expression is parsed using syntax rules defined by the language:
\begin{verbatim}
5 + 3
\end{verbatim}

\item This difference comes from the compiler's parser: function calls follow a "function-call grammar rule", while operators follow "expression grammar rules" with precedence and associativity.

\item This is why operators such as \verb|+|, \verb|-|, and \verb|new| do not use function-style parentheses for arguments, even though they conceptually operate on operands.

\item In summary, functions are explicitly invoked using argument lists, while operators are built into the language syntax and are interpreted as part of expressions during parsing.

	\pbreak

	\item When the compiler parses an expression, it first performs lexical and syntactic analysis. If it encounters an \textbf{identifier}, it determines whether it refers to a variable, function, or other named entity. If it encounters an \textbf{operator}, it applies the language grammar rules for expressions, using operator precedence and associativity to determine how operands are grouped and evaluated.

\item For example, in a function call like:
\begin{verbatim}
add(5, 3);
\end{verbatim}
the compiler treats \verb|add| as an identifier referring to a function and processes the parentheses as a function-call argument list.

\item In contrast, in an expression like:
\begin{verbatim}
5 + 3 * 2
\end{verbatim}
the compiler recognizes \verb|+| and \verb|*| as operators and constructs an expression tree based on precedence rules, evaluating multiplication before addition.

\pbreak

\item When the compiler encounters a custom operator, it does not introduce new syntax rules. The expression is first parsed using the same grammar rules as built-in operators. For example, an expression like \verb|a + b| is always parsed as a binary operator expression.

\item During semantic analysis, the compiler then performs \textbf{operator overload resolution}. It checks whether a valid function named \verb|operator+| exists for the given operand types, either as a member function or a non-member function.

\item For example, the expression:
\begin{verbatim}
a + b
\end{verbatim}
may be resolved internally as:
\begin{verbatim}
a.operator+(b)
\end{verbatim}

\item If a matching overload is found, the compiler replaces the operator expression with the corresponding function call. If no valid overload exists, a compilation error occurs.

\item In summary, custom operators are not new language constructs; they are existing operators that are bound to user-defined functions during the overload resolution phase of compilation.

	\pbreak

\item In C++, operators have a fixed number of operands (called \textbf{arity}) determined by the language grammar. Operators cannot accept an arbitrary number of arguments like functions.

\item \textbf{Unary operators} operate on one operand:
\begin{verbatim}
!a
++a
-a
\end{verbatim}

\item \textbf{Binary operators} operate on two operands:
\begin{verbatim}
a + b
a < b
a == b
\end{verbatim}

\item \textbf{Ternary operator} is the only operator with three operands:
\begin{verbatim}
condition ? expr1 : expr2
\end{verbatim}

\item The compiler does not allow changing the arity of an operator. For example, \verb|a + b + c| is not a three-operand addition operator, but is instead parsed as:
\begin{verbatim}
(a + b) + c
\end{verbatim}

\item During compilation, the expression is first parsed using fixed grammar rules (unary, binary, or ternary forms), and then overload resolution determines whether a matching \verb|operator+|, \verb|operator!|, etc., function exists for the given operand types.

\item In summary, operator argument count is fixed by the language, and the compiler handles operators by first parsing expression structure and then binding them to appropriate overloaded functions if available.
\end{itemize}

\newpage
\subsubsection{Operator Overloading Examples (Unary, Binary, and Ternary-like)}

\begin{itemize}

\item \textbf{Unary operator example (prefix ++)}:
\begin{verbatim}
#include <iostream>
using namespace std;

class A {
    int x;

public:
    A(int val) : x(val) {}

    // Unary prefix ++
    A operator++()
    {
        ++x;
        return *this;
    }

    void show()
    {
        cout << x << endl;
    }
};

int main()
{
    A obj(5);
    ++obj;
    obj.show();   // 6
}
\end{verbatim}

\item \textbf{Binary operator example (+)}:
\begin{verbatim}
#include <iostream>
using namespace std;

class A {
    int x;

public:
    A(int val) : x(val) {}

    // Binary +
    A operator+(const A& other)
    {
        return A(x + other.x);
    }

    void show()
    {
        cout << x << endl;
    }
};

int main()
{
    A a(5), b(10);
    A c = a + b;

    c.show();   // 15
}
\end{verbatim}

\item \textbf{Ternary operator (?:) note}:
\begin{itemize}
    \item The ternary operator \verb|?:| cannot be overloaded in C++.
    \item It is built into the language grammar and cannot be redefined.
\end{itemize}

\item \textbf{Ternary-like behavior using a function}:
\begin{verbatim}
#include <iostream>
using namespace std;

class A {
    int x;

public:
    A(int val) : x(val) {}

    void show()
    {
        cout << x << endl;
    }
};

A choose(bool cond, A a, A b)
{
    return cond ? a : b;
}

int main()
{
    A a(5), b(10);

    A c = choose(true, a, b);
    c.show();   // 5
}
\end{verbatim}


\pbreak 

\item In an expression like \verb|a + b|, the operator \verb|+| is not a standalone operation. If \verb|operator+| is defined as a member function inside class \verb|A|, the expression is interpreted as a function call on the left-hand object.

\item For example:
\begin{verbatim}
A c = a + b;
\end{verbatim}

\item The compiler rewrites this as:
\begin{verbatim}
A c = a.operator+(b);
\end{verbatim}

\item This means the operator "belongs" to the left operand (\verb|a|), and the right operand (\verb|b|) is passed as an argument.

\item In general:
\begin{itemize}
    \item \verb|a + b| $\rightarrow$ \verb|a.operator+(b)|
    \item \verb|b + a| $\rightarrow$ \verb|b.operator+(a)|
\end{itemize}

\item This is why operator order matters when using member operator overloading. If symmetry is required, a non-member (friend) operator overload is often used instead.

\item \fbox{\textbf{To create an overloaded operator function, you write \textit{operator} followed by the operator symbol.}}
\end{itemize}


\newpage

\subsubsection{Incrementing (++) and Decrementing ($--$ on Pointers)}
\begin{itemize}
	\item A \textbf{pointer} contains a \textbf{memory address}, ex: \verb|0x002EFB34| --- the \textbf{address} where, say, an integer is placed.

	\item The \textbf{integer} is \textbf{4 bytes} long on our system, so it \textbf{occupies} \textbf{4 places in memory} from:
		\begin{itemize}
			\item \verb|0x002EFB34| to \verb|0x002EFB37| inclusive
		\end{itemize}
	
	\item \textbf{Incrementing} the \textbf{pointer} using the \textbf{operator} (++) would \underline{\textbf{not}} result in the \textbf{pointer} pointing to \verb|0x002EFB35| (notice the 4 to the 5), for pointing to the middle of the integer would be pointless.

	\pbreak

\item An \textbf{increment} or \textbf{decrement} operation on a \textbf{pointer} is interpreted by the \textbf{compiler} as your need to point to the \underline{\textbf{next value}} in the \textbf{block of memory}, assuming it would be of the same \textbf{type}, and \underline{\textbf{not}} to the \textbf{next byte} (unless the value type is \textbf{1 byte} large, like a \verb|char| for instance).\\

	\textbf{Example:}

Suppose we have the following code:

\begin{verbatim}
int arr[] = {10, 20, 30, 40};
int *p = arr;
\end{verbatim}

Assume an \verb|int| occupies \textbf{4 bytes} in memory. The array could be stored as follows:

\[
\begin{array}{|c|c|}
\hline
\textbf{Address} & \textbf{Value} \\
\hline
1000 & 10 \\
1004 & 20 \\
1008 & 30 \\
1012 & 40 \\
\hline
\end{array}
\]

Initially, \verb|p| points to address \verb|1000|, so:

\begin{verbatim}
*p    // 10
\end{verbatim}

After incrementing the pointer:

\begin{verbatim}
p++;
\end{verbatim}

the compiler knows that \verb|p| is an \verb|int *|, so it advances the pointer by
\verb|sizeof(int)| (4 bytes), not by 1 byte. The pointer now points to address
\verb|1004|:

\[
\begin{array}{|c|c|}
\hline
\textbf{Address} & \textbf{Value} \\
\hline
1000 & 10 \\
\mathbf{1004 \leftarrow p} & \mathbf{20} \\
1008 & 30 \\
1012 & 40 \\
\hline
\end{array}
\]

Thus,

\begin{verbatim}
*p    // 20
\end{verbatim}

Incrementing the pointer again moves it to address \verb|1008|, where:

\begin{verbatim}
*p    // 30
\end{verbatim}

\noindent
In general, for a pointer of type \verb|T *|, the operation \verb|p++| advances
the pointer by \verb|sizeof(T)| bytes, making it point to the \textbf{next object}
of type \verb|T| in memory rather than simply the next byte.

\pbreak

\item So, \textbf{incrementing} a \textbf{pointer} such as 
	\begin{verbatim}
int age = 30;
int* pointsToInt = &age;
	\end{verbatim}
	results it being \textbf{incremented} by \textbf{4 bytes}, which is the \verb|sizeof(int)|.

\item \textbf{\fbox{It increments based on the \underline{size of the type} it points to, \underline{not} the size of the pointer itself.}}

\item Using ++ on this \textbf{pointer} is telling the compiler that you want it to point to the \textbf{next consecutive integer}.

	Hence, after incrementing, the pointer would then point to \verb|0x002EFB38| (from 4 to 8 in hex)

\item Similarly, adding 2 to this pointer would result in it moving 2 \textbf{integers} ahead, that is \textbf{8 bytes} (in this case).

\item \textbf{Decrementing} pointers using $--$ gives the same effect but you reduce by the sizeof what the pointer is pointing to.
\end{itemize}

\newpage

\subsubsection{Using the const keyword on Pointers}
\begin{itemize}
	\item Lesson 3 we learned that declaring a \textbf{variable} as a \verb|const| ensures that value of the variable is \textbf{fixed} as the initialization value for the \textbf{life} of the variable.

	\item The value of a \textbf{const-variable} \underline{cannot be changed}, and therefore it \underline{cannot be used} as an \fbox{\textbf{l-value}}.
		\pbreak

	\item \textbf{Pointers} are \textbf{variable} too, and hence the \verb|const| keyword can be used here.
	\item \underline{\textbf{However}}, \textbf{pointers} are a \textbf{special kind of variable} as they contain a \textbf{memory address} and are used to modify \textbf{memory} at that address.

	\item Thus, when it comes to \textbf{pointers} and \textbf{constants}, you have the following combinations:

		\begin{enumerate}
\item \boxednote{
		The \textbf{address} contained in the \textbf{pointer} is \underline{\textbf{constant}} and \underline{\textbf{cannot be changed}}, yet the \textbf{data} at that address \underline{\textbf{can be changed}}.
}
			\begin{verbatim}
int daysInMonth = 30;
int* const pDaysInMonth = &daysInMonth;
*pDaysInMonth = 31; // OK! Data pointed to can be changed
int daysInLunarMonth = 28;
pDaysInMonth = &daysInLunarMonth; // NOT OK! Cannot change address!
			\end{verbatim}

	Note here also, the \textbf{const} keyword is after the int*!!!\\ 

	\begin{verbatim}
const int *p;
^^^^^^^^^
data is const (the int being pointed to is const)

int *const p;
     ^^^^^
pointer is const (the pointer variable p cannot be changed)
	\end{verbatim}

	Also

\begin{verbatim}
Valid:

const int x = 5;
int const x = 5;

Invalid:

int x const = 5;   // not allowed in C/C++

-----

const int x -> x is a const int
int const x -> same thing (just different order)
int x const -> not legal grammar in C/C++
\end{verbatim}

	\pbreak

\item \boxednote{
		\textbf{Data} pointed to is \textbf{constant} and \textbf{cannot be changed}, yet the \textbf{address} contained in the pointer \textbf{can be changed} --- that is, the \textbf{pointer} can also point elsewhere: 
	}
\begin{verbatim}
int hoursInDay = 24;
const int* pointsToInt = &hoursInDay; // OK! address can be changed but not the data.
int monthsInYear = 12;
pointsToInt = &monthsInYear; // OK again for same reason
*pointsToInt = 13; // NOT OK! Cannot change data being pointed to

int* newPointer = pointsToInt; // NOT OK! Cannot assign const to non-const 
// cause this would entail you being able to change data.
// const int *p  means:
// "I promise I will not modify the int through this pointer"
// int* newPointer = pointsToInt; violates that contract.
\end{verbatim}

\pbreak

\item \boxednote{
		Both the \textbf{address} conatined in the pointer and the \textbf{value} being pointed to are \textbf{constant} and \underline{cannot be changed} (most restrictive variant):
	}

\begin{verbatim}
int hoursInDay = 24;
const int* const pHoursInDay = &hoursInDay;

*pHoursInDay = 25; // NOT OK! Cannot change data being pointed to.

int daysInMonth = 30;
pHoursInDay = &daysInMonth; // NOT OK! Cannot change address.
\end{verbatim}
	\end{enumerate}
\end{itemize}

\pbreak

\begin{itemize}
	\item These different forms of \textbf{const} is particularly useful when passing \textbf{pointers} to \textbf{functions}.
	\item Function parameters need to be declared to support the highest possible (restrictive) level of \textbf{const-ness}, to ensure that the function does \textbf{not} modify the pointed value when it is not supposed to.
\end{itemize}

\newpage
\subsubsection{Passing by ref C vs C++}

\begin{verbatim}
// C++ only (reference)
void foo(int& v) {
    v = 5;
}

int main() {
    int v = 0;
    foo(v);
}

// C/C++ compatible (pointer to int)
void foo(int* v) {
    *v = 5;
}

int main() {
    int v = 0;
    foo(&v);
}
\end{verbatim}

\newpage
\subsubsection{Clarification for pointer declarations}
The confusing part is that the \verb|*| is \textbf{part of the declaration}, not part of the type that is being made \verb|const|.

\subsubsection*{Step 1: A Regular Variable}

A normal integer declaration is:

\begin{verbatim}
int x;
\end{verbatim}

If we make the variable constant:

\begin{verbatim}
const int x;
// or equivalently
int const x;
\end{verbatim}

the \verb|int| itself becomes \textbf{constant}.

\subsubsection*{Step 2: Introduce a Pointer}

Consider the declaration:

\begin{verbatim}
int *p;
\end{verbatim}

This is read as:

\begin{quote}
``\verb|p| is a pointer to an \verb|int|.''
\end{quote}

Notice that there are now \textbf{two different things (or both to be the most restrictive!)} that could be made constant:

\begin{enumerate}
    \item the \textbf{integer} being pointed to,
    \item the \textbf{pointer} itself.
\end{enumerate}

\subsubsection*{Case 1: Make the Integer Constant}

Starting from:

\begin{verbatim}
int *p;
\end{verbatim}

we can make the integer constant:

\begin{verbatim}
const int *p;
\end{verbatim}

The \verb|const| applies to the \verb|int|, not to the pointer.

You can think of it conceptually as:

\begin{verbatim}
(const int) *p
\end{verbatim}

This means:

\begin{itemize}
    \item the integer cannot be modified through \verb|p|,
    \item the pointer itself may still point somewhere else.
\end{itemize}

Conceptually:

\begin{verbatim}
p
|
v
+-----+
| 42  |   <- const int
+-----+
\end{verbatim}

Therefore,

\begin{verbatim}
*p = 10;        // NOT OK
p = &otherInt;  // OK
\end{verbatim}

\subsubsection*{Case 2: Make the Pointer Constant}

Again start with:

\begin{verbatim}
int *p;
\end{verbatim}

This time make the pointer constant:

\begin{verbatim}
int *const p;
\end{verbatim}

Here, the \verb|const| appears \emph{after} the \verb|*|, so it applies to the pointer itself.

Conceptually:

\begin{verbatim}
(int *) const p
\end{verbatim}

This means:

\begin{itemize}
    \item the pointer cannot be changed,
    \item the integer being pointed to may still be modified.
\end{itemize}

Conceptually:

\begin{verbatim}
const p
|
v
+-----+
| 42  |
+-----+
\end{verbatim}

Therefore,

\begin{verbatim}
*p = 10;        // OK
p = &otherInt;  // NOT OK
\end{verbatim}

\subsubsection*{Why Doesn't \texttt{const int *} Mean ``Constant Pointer''?}

This is because, in C/C++, the \verb|*| belongs to the \textbf{declarator} (the variable being declared), not to the base type.

For example,

\begin{verbatim}
int *p, q;
\end{verbatim}

does \textbf{not} declare two pointers.

Instead, it declares:

\begin{verbatim}
int *p;   // pointer
int q;    // ordinary int
\end{verbatim}

If the \verb|*| were part of the type, then both \verb|p| and \verb|q| would be pointers, but they are not.

Therefore, the compiler interprets

\begin{verbatim}
const int *p;
\end{verbatim}

as

\begin{verbatim}
(const int) *p
\end{verbatim}

rather than

\begin{verbatim}
const (int *)
\end{verbatim}

\subsubsection*{Summary}

There are two separate objects involved in a pointer declaration:

\begin{enumerate}
    \item the object being pointed to,
    \item the pointer itself.
\end{enumerate}

The position of \verb|const| determines which one is constant:

\[
\begin{array}{|l|l|}
\hline
\textbf{Declaration} & \textbf{Meaning} \\
\hline
\verb|int *p| & \text{pointer to an int} \\
\verb|const int *p| & \text{pointer to a constant int} \\
\verb|int *const p| & \text{constant pointer to an int} \\
\verb|const int *const p| & \text{constant pointer to a constant int} \\
\hline
\end{array}
\]

A useful way to read declarations is:

\begin{itemize}
    \item \verb|const int *p| \quad $\rightarrow$ \quad ``\verb|p| is a pointer to a \textbf{constant int}.''
    \item \verb|int *const p| \quad $\rightarrow$ \quad ``\verb|p| is a \textbf{constant pointer} to an int.''
\end{itemize}

Once you recognize that there are \textbf{two different things} that can be made constant (the pointee and the pointer), the placement of \verb|const| becomes much more intuitive.

\pbreak

\begin{itemize}
    \item \textbf{Add const (always allowed):}
    \begin{verbatim}
int *p
const int *p      (OK: add const to data)
int *const p      (OK: add const to pointer)
const int *const p (OK: add const to both)
    \end{verbatim}

    \item \textbf{Remove const (NOT allowed implicitly):}
    \begin{verbatim}
const int *p  -> int *p      (NOT OK)
int *const p  -> int *p      (NOT OK)
const int x   -> int x       (NOT OK)
    \end{verbatim}

    \item \textbf{Safe conversions (const-preserving):}
    \begin{verbatim}
int *p        -> const int *p   (OK)
int x         -> const int x    (OK)
int *p        -> int *const p   (OK only when initializing)
    \end{verbatim}

    \item \textbf{Pointer rule intuition:}
    \begin{verbatim}
You can always promise "I won't modify",
but you cannot take back that promise implicitly.
    \end{verbatim}
\end{itemize}

\newpage

\subsubsection{Passing Pointers to Functions}
\begin{itemize}
	\item Pointers are an effective way to \textbf{pass memory space} that contains relevant data for functions to work on.

	\item The \textbf{memory space} shared can also return the result of an \textbf{operation}.
	\item Given the previous section, you should modify pointers as parameters to be const or not if you want to restrict the callee modifiying the pointer or the value pointed to or both.

	\item \fbox{\textbf{Listing 8.10} Use const in pointer function parameters}
\end{itemize}

\newpage

\subsubsection{Pointers in Static vs Dynamic Memory}

\begin{itemize}

\item \textbf{Pointers are not tied only to dynamic memory.}
Pointers can store the address of variables in \textbf{stack (static)} memory or \textbf{heap (dynamic)} memory.

\item For example, static arrays are just a single pointer underneath and we go to each element via adding the sizeof(element\_type) to the current pointer and increment or decrement the address by that size in bytes (cause \verb|sizeof| returns bytes)

\item \textbf{Static (stack) memory example:}
\begin{verbatim}
int x = 10;
int *p = &x;
\end{verbatim}

\begin{itemize}
    \item \verb|x| is stored in stack memory
    \item \verb|p| simply stores its address
    \item no dynamic allocation is involved
\end{itemize}

\item \textbf{Array example (contiguous static memory):}
\begin{verbatim}
int arr[] = {1, 2, 3};
int *p = arr;
\end{verbatim}

\begin{itemize}
    \item arrays decay into pointers in expressions
    \item pointer arithmetic allows traversal of elements
    \item \verb|p++| moves to the next array element, not next byte
\end{itemize}

\item \textbf{Dynamic (heap) memory example:}
\begin{verbatim}
int *p = malloc(sizeof(int));
*p = 42;

// or

int *p = new int;
*p = 42;
\end{verbatim}

\begin{itemize}
    \item memory is allocated on the heap
    \item pointer is required to access and manage that memory
    \item must be manually freed in C using \verb|free(p)| or \verb|delete ptr| or \verb|delete[] ptr| in C++
\end{itemize}

\item \textbf{Function parameter example (pass by address):}
\begin{verbatim}
void foo(int *p) {
    *p = 100;
}

int main() {
    int x = 5;
    foo(&x);
}
\end{verbatim}

\begin{itemize}
    \item allows modification of original variable
    \item avoids copying large data structures
\end{itemize}

\item \textbf{Key insight:}
Pointers are a \textbf{general mechanism for storing memory addresses}, not a feature exclusive to dynamic memory.

\item \textbf{Exam takeaway:}
\begin{itemize}
    \item stack variables can have pointers
    \item heap variables require pointers
    \item arrays behave like pointers in most expressions
    \item pointer arithmetic works based on data type size, not byte-by-byte movement
\end{itemize}

\end{itemize}

\newpage

\subsection{Common Programming Mistakes when using Pointers}

\begin{itemize}

\item \textbf{Uninitialized pointers (wild pointers)}
\begin{verbatim}
int *p;
*p = 10;   // undefined behavior
\end{verbatim}

\begin{itemize}
    \item Pointer is not initialized to valid memory.
    \item Must always assign a valid address before dereferencing.
\end{itemize}

\item \textbf{Dangling pointers (use-after-free)}
\begin{verbatim}
int *p = malloc(sizeof(int));
free(p);
*p = 5;   // undefined behavior
\end{verbatim}

\begin{itemize}
    \item Memory has been freed but pointer still holds the old address.
    \item Accessing it leads to undefined behavior.
\end{itemize}

\item \textbf{Memory leaks}
\begin{verbatim}
int *p = malloc(sizeof(int));
p = NULL;   // lost reference to allocated memory
\end{verbatim}

\begin{itemize}
    \item Original allocated memory is no longer reachable.
    \item This causes memory that cannot be freed (leak).
\end{itemize}

\item \textbf{Null pointer dereference}
\begin{verbatim}
int *p = NULL;
*p = 10;   // undefined behavior
\end{verbatim}

\begin{itemize}
    \item NULL does not point to valid memory.
    \item Dereferencing NULL leads to crashes or undefined behavior.
\end{itemize}

\item \textbf{Incorrect pointer types}
\begin{verbatim}
int x = 10;
float *p = &x;   // type mismatch
\end{verbatim}

\begin{itemize}
    \item Pointer type must match the type of the object.
    \item Otherwise memory is interpreted incorrectly.
\end{itemize}

\item \textbf{Confusing pointer assignment vs dereference}
\begin{verbatim}
int x = 5;
int *p = &x;

p = 10;   // incorrect: changes pointer, not value
*p = 10;  // correct: changes value at address
\end{verbatim}

\begin{itemize}
    \item \verb|p| stores an address.
    \item \verb|*p| accesses the value at that address.
\end{itemize}

\item \textbf{Out-of-bounds pointer access}
\begin{verbatim}
int arr[3] = {1, 2, 3};
int *p = arr;

p[5] = 10;   // undefined behavior
\end{verbatim}

\begin{itemize}
    \item Pointer arithmetic must stay within allocated memory.
    \item Accessing outside bounds is undefined behavior.
\end{itemize}

\end{itemize}

\newpage

\subsection{References}

\begin{itemize}
    \item A \vocab{reference} is an \textbf{alias} for an existing variable.
    \item Once a reference is created, it \textbf{must be initialized immediately}.
    \item A reference does not store a separate copy of the data; it refers to the \textbf{same memory location} as the original variable.
    \item Any operation on the reference is performed on the \textbf{original variable}.

    \pbreak

    \item References are declared using the \textbf{reference operator (\&)}:
    \begin{verbatim}
VarType original = value;
VarType& ref = original;
    \end{verbatim}

    \item After binding, the reference becomes an alternative name for the same variable:
    \begin{verbatim}
ref = 10;   // changes original as well
    \end{verbatim}

    \pbreak

    \item \textbf{Important properties of references:}
    \begin{itemize}
        \item Must be initialized at declaration
        \item Cannot be null (unlike pointers)
        \item Cannot be reseated to refer to another variable
        \item Acts as an automatic dereference (no need for *)
    \end{itemize}

    \pbreak

    \item \textbf{Example:}
    \begin{verbatim}
int x = 5;
int& r = x;

r = 10;   // x is now 10
    \end{verbatim}

    \item \textbf{Key intuition:}
    A reference behaves like the original variable, just under a different name.
\end{itemize}

\newpage

\subsubsection{References as Function Parameters}

\begin{itemize}
    \item References are commonly used as \textbf{function parameters} in C++ to allow functions to modify the original argument.

    \item When a parameter is declared as a reference, the function receives an \textbf{alias} to the original variable, not a copy.

    \pbreak

    \item \textbf{Pass-by-value vs Pass-by-reference:}

    \begin{verbatim}
// Pass-by-value (copy)
void foo(int x) {
    x = 10;   // does NOT affect original
}

// Pass-by-reference (no copy)
void bar(int& x) {
    x = 10;   // modifies original
}
    \end{verbatim}

    \item In pass-by-value:
    \begin{itemize}
        \item A copy of the argument is created
        \item Changes inside the function do not affect the original variable
    \end{itemize}

    \item In pass-by-reference:
    \begin{itemize}
        \item No copy is made
        \item The parameter becomes an alias for the original variable
        \item Changes inside the function affect the caller's variable
    \end{itemize}

    \pbreak

    \item \textbf{Example:}
    \begin{verbatim}
void increment(int& x) {
    x++;
}

int main() {
    int a = 5;
    increment(a);
    // a is now 6
}
    \end{verbatim}

    \pbreak

    \item \textbf{Why use references in parameters?}
    \begin{itemize}
        \item Avoid copying large objects (performance improvement)
        \item Allow functions to modify input arguments
        \item Cleaner syntax compared to pointers (no * or \&)
    \end{itemize}

    \pbreak

    \item \textbf{Important rule:}
    \begin{itemize}
        \item A reference parameter behaves exactly like the original variable inside the function.
        \item It is just another name for the same memory location.
    \end{itemize}
\end{itemize}

\newpage

\subsubsection{References vs Pointers as Function Arguments}

\begin{itemize}

\item In C++, there are two common ways to allow a function to modify an argument:
\begin{itemize}
    \item \textbf{References (C++)}
    \item \textbf{Pointers (C/C++)}
\end{itemize}

\pbreak

\item \textbf{Pointer-based arguments (C/C++ compatible):}
\begin{verbatim}
void foo(int* p) {
    *p = 10;
}

int main() {
    int x = 5;
    foo(&x);
}
\end{verbatim}

\begin{itemize}
    \item The function receives the \textbf{address} of the variable.
    \item You must explicitly pass \verb|&x| (address-of operator).
    \item Inside the function, you must dereference using \verb|*p|.
\end{itemize}

\pbreak

\item \textbf{Reference-based arguments (C++ only):}
\begin{verbatim}
void foo(int& x) {
    x = 10;
}

int main() {
    int x = 5;
    foo(x);
}
\end{verbatim}

\begin{itemize}
    \item The function receives an \textbf{alias} to the original variable.
    \item No need to use \verb|&| or \verb|*| at the call site.
    \item The syntax is cleaner and behaves like the original variable.
\end{itemize}

\pbreak

\item \textbf{Key difference in calling style:}

\begin{verbatim}
Pointer version:     foo(&x);
Reference version:   foo(x);
\end{verbatim}

\pbreak

\item \textbf{Conceptual difference:}
\begin{itemize}
    \item Pointer: function receives an \textbf{address} and manually dereferences it.
    \item Reference: function receives an \textbf{alias} automatically (implicit dereference).
\end{itemize}

\pbreak

\item \textbf{Important note:}
\begin{itemize}
    \item You do NOT write \verb|foo(&arg)| when the function expects a reference.
    \item \verb|&| is only used when the function expects a pointer.
\end{itemize}

\end{itemize}

\newpage

\begin{itemize}
	\item \fbox{\textbf{Listing 8.17}}
	\item \fbox{\textbf{Listing 8.18}}
	\item \fbox{\textbf{Listing 8.19}}
\end{itemize}

\pbreak

\subsubsection{Using const on References}
\begin{itemize}
	\item You might need to have \textbf{references} that are \underline{not allowed} to change the \textbf{value}of the original variable being \textbf{aliased}.

	\item Using \textbf{const} when declaring such references is the way to achieve this:

		\begin{verbatim}
int original = 30;
const int& constRef = original; // think of (const belongs to const (int) &, so can't change reference value)
constRef = 40; // NOT OK! constRef can't change value in original
int& ref2 = constRef; // NOT OK! ref2 is not const
const int& constRef2 = constRef; // OK
		\end{verbatim}
\end{itemize}

\newpage

\subsubsection{Extended: Using const with References}

\begin{itemize}
    \item A \textbf{const reference} prevents modification of the value being referenced.

    \begin{verbatim}
int original = 30;
const int& ref = original;

ref = 40;   // NOT OK: cannot modify through const reference
    \end{verbatim}

    \item A const reference still refers to the \textbf{same object}, but only allows read-only access.

    \item You can bind const references to const or non-const variables:

    \begin{verbatim}
int x = 10;
const int& r1 = x;     // OK
const int& r2 = r1;    // OK
    \end{verbatim}

    \item You cannot modify the original through a const reference:

    \begin{verbatim}
r1 = 20;   // NOT OK
    \end{verbatim}

    \item \textbf{Important rule:}
    \begin{itemize}
        \item References in C++ are already fixed (cannot be reseated)
        \item \verb|const| only controls whether the value can be modified
        \item \verb|int& const| is meaningless in practice
    \end{itemize}

	\begin{verbatim}
int& -> fixed alias (cannot change what it refers to)
const int& -> fixed alias + cannot modify value
int& const -> conceptually meaningless / not used
	\end{verbatim}
\end{itemize}

\newpage
\subsubsection{Definition of Reference in C++ and why it's not valid in C}

In C++, a \textit{reference} is a language-level feature that creates an alias for an existing variable. Unlike pointers, references are not separate objects that store memory addresses; instead, they directly refer to the original variable.

\begin{itemize}
    \item A reference must be initialized when it is declared.
    \item A reference cannot be reseated to refer to another variable after initialization.
    \item A reference behaves as an alias of the original variable.
    \item C does not support this language feature; it only provides pointers.
\end{itemize}

In C++, references are declared using the \texttt{\&} symbol:

\begin{verbatim}
int a = 10;
int& ref = a;   // ref is a reference to a
\end{verbatim}

This allows \texttt{ref} to be used exactly like \texttt{a}:

\begin{verbatim}
ref = 20;   // modifies a directly
\end{verbatim}

In contrast, C does not have references as a built-in type system feature. Instead, it uses pointers to achieve similar indirect access to variables:

\begin{verbatim}
int a = 10;
int *p = &a;   // pointer to a
\end{verbatim}

However, pointers differ fundamentally from C++ references:

\begin{itemize}
    \item Pointers can be reassigned to point to different variables.
    \item Pointers require explicit dereferencing using \texttt{*}.
    \item Pointers store memory addresses explicitly, whereas references are aliases.
\end{itemize}

For example:

\begin{verbatim}
int a = 10;
int b = 20;

int* p = &a;
p = &b;   // valid: pointer reassignment

// In C++, a reference cannot be "re-pointed":
int& r = a;
// r = b;  // this assigns the value of b to a, not rebinding r
\end{verbatim}

Thus, the key difference is that C++ provides true reference semantics as part of the language, while C only provides pointer-based indirect access, which is more flexible but less abstract.

\newpage
\subsubsection{QA}
\begin{itemize}
\item \textbf{Q: Why dynamically allocated when you can do with static arrays where you don't need to worry about deallocation?}

	A: Static arrays have \textbf{fixed size} and will neither scale upward if your application needs more \textbf{memory} nor will they optimize your application needs less. Dynamic memory fixes that. \\ 

\item \textbf{Q: I have 2 pointers:}
	\begin{verbatim}
int* pointToAnInt = new int;
int* pCopy = pointToAnInt;
	\end{verbatim}
	\textbf{Am I not better off calling 'delete' using both to ensure memory is gone?}

	A: No. You are allowed to invoke \verb|delete| \underline{only once} on the \textbf{address} returned by \verb|new|.

	Also, you would ideally \textbf{avoid} having two pointers pointing to the \textbf{same address} b/c performing \verb|delete| on any one would \textbf{invalidate} the other. 

	Your program should also not be written in a way that you have any uncertainity about the validity of \textbf{pointers} used.\\

\item \textbf{Q: When should i use 'new(nothrow)'?}

	A: If you don't want to handle the \textbf{exception} \verb|std::bad_alloc|, you use the \verb|nothrow| version of \textbf{operator} \verb|new| that returns the \textbf{NULL} if the requested allocation \textbf{fails}.\\ 

\item \textbf{Q: I can call a function to calculate the area using the following two methods:}

\begin{verbatim}
void CalculateArea(const double* const ptrRadius, double* const ptrArea);
void CalculateArea(const double& radius, double& area);
\end{verbatim}

\textbf{Which variant should I prefer?}

A: In general, prefer the \textbf{reference version} (C++):

\begin{verbatim}
void CalculateArea(const double& radius, double& area);
\end{verbatim}

\textbf{Reason:}
\begin{itemize}
    \item References provide \textbf{cleaner syntax} (no \*, no \& at call site).
    \item They guarantee a \textbf{valid object} must be passed (cannot be null).
    \item They make the function easier to read and use.
\end{itemize}

\pbreak

\textbf{Important clarification:}
\begin{itemize}
    \item It is not that pointers are always "invalid".
    \item Pointers can be valid and safe if used correctly.
    \item However, pointers can also be \textbf{null or uninitialized}, which introduces extra risk.
\end{itemize}

\pbreak

\textbf{Comparison:}

\begin{itemize}
    \item Pointer version:
    \begin{itemize}
        \item Can be null or invalid if misused
        \item Requires explicit dereferencing (\verb|*|)
        \item More verbose at call site (\verb|&x|)
    \end{itemize}

    \item Reference version:
    \begin{itemize}
        \item Must always refer to a valid object
        \item Cleaner syntax (acts like a normal variable)
        \item Safer by default in most use cases
    \end{itemize}
\end{itemize}

\pbreak

\textbf{C vs C++ note:}
\begin{itemize}
    \item C does not support references.
    \item Therefore, C must use pointers.
    \item C++ supports both pointers and references.\\
\end{itemize}


\item \textbf{Q: I have a pointer}
	\begin{verbatim}
int number = 30;
const int* pointToAnInt = &number;
	\end{verbatim}

\textbf{I understand that I cannot change the value of 'number' using the pointer 'pointToAnInt' due to the 'const' declaration. Can I assign 'pointToAnInt' to a non-const pointer and then use it to manipulate the value contained in integer number?}

A: No, you cannot change the \textbf{const-correctness} of the pointer:
\begin{verbatim}
int* pAnother = pointToAnInt; // cannot assign pointer to const to a non-const
\end{verbatim}

\item \textbf{Q: What is the difference btw these 2 declarations?}
\begin{verbatim}
int myNums[100];
int* myArrays[100];
\end{verbatim}

A: \verb|myNums| is an array of integers --- that is, \verb|myNums| is a \textbf{pointer} to a \textbf{memory location} that holds 100 integers, points to the \textbf{first index} at index 0.

It is the static alternative to the following:

\begin{verbatim}
int* mNums = new int[100]; // dyanmically allocated array
delete[] myNums;
\end{verbatim}

\verb|myArrays|, on the other hand, is an array of 100 \textbf{pointers}, each pointer being capable of pointing to an integer or an array of integers.
\end{itemize}

\newpage

\subsubsection{Quiz}
\begin{enumerate}
	\item Why can't you assign a \verb|const| reference to a \verb|non-const| reference?

	A: You cannot bind a \verb|const| reference to a \verb|non-const| reference because it would allow the possibility of modifying a value that was originally declared as constant. This would violate C++'s type safety rules by potentially stripping away const-correctness. In general, C++ prevents dropping \verb|const| qualifiers through reference assignment to ensure that immutable objects remain protected.\\

	\item Are \verb|new| and \verb|delete| functions?

	A: No, \verb|new| and \verb|delete| are not functions. They are C++ language operators. Unlike functions, they are built into the language and can be overloaded, but they do not behave like normal function calls.\\

	\item What is the nature of value contained in a \textbf{pointer variable}?

	A: A pointer variable stores a memory address as its value. This address refers to the location of another variable or object in memory.\\

	\item What \textbf{operator} would you use to access the data pointed by a pointer?

	A: The dereference operator \verb|*| is used to access the value stored at the memory address pointed to by a pointer.\\
\end{enumerate}

\newbp

\section{Classes and Objects}
\begin{itemize}
	\item Thus far, we've learned \vocab{procedural programming} and haven't dived into \vocab{object-oriented programming}. We now will.
\end{itemize}

\pbreak

\subsection{Classes}
\begin{itemize}
	\item You declare a \vocab{class} by using the keyword \verb|class| followed by the \textbf{name} of the class, followed by a \textbf{statement block} \verb|{...}|
\item A \textbf{declaration} of a \textbf{class} tells the \textbf{compiler} about the class and its properties.
	\pbreak

\item A class consists of its \textbf{fields} (data members) and \vocab{methods} (member functions).

\item \vocab{Encapsulation} is the ability to logically \textbf{group data} and \textbf{functions} and we do that in our class in object-oriented programming.

\end{itemize}

\pbreak

\subsubsection{An Object as an Instance of a Class}
\begin{itemize}
	\item A \vocab{class} is the \textbf{blueprint} and declaring a class alone has no interesting effect.

	The \textbf{real-world} avator of a class at program execution time is an \vocab{object}. 

\item To use this blueprint you create an \textbf{instance} of the class, i.e. an \vocab{object}.

\item You use this \textbf{object} to access its \textbf{attributes} and \textbf{member methods}. (or just attributes/fields and methods/functions)

\pbreak

\item Creating an \textbf{object} of type \verb|class Human| is similar to creating an instance of another type, say \verb|double|:

\begin{verbatim}
double pi = 3.14;
Human firstMan;		// statically allocated object
\end{verbatim}

\item Alternatively, you can \textbf{dynamically} create an \textbf{instance} of \verb|class Human| using \verb|new| as you would for another type say \verb|int|:

\begin{verbatim}
int* pointToAnInt = new int;		// an integer allocated dynamically
delete pointsToAnInt;

Huamn* firstHuman = new Human();	// dynamically allocated Human
delte firstHuman;
\end{verbatim}

\end{itemize}


\newpage
\subsubsection{Accessing Members Using the Dot Operator (\texttt{.})}

The dot operator (\texttt{.}) is used to access the data members and member functions of an object. It is applicable when working directly with an instance of a class rather than a pointer to that instance.

\begin{itemize}
    \item Used with objects (instances) of a class.
    \item Accesses both data members and member functions.
    \item The object itself owns the members being accessed.
\end{itemize}

For example:

\begin{verbatim}
class Student {
public:
    std::string name;
    int age;

    void display() {
        std::cout << name << " is "
                  << age << " years old.\n";
    }
};

Student s;

s.name = "Alice";
s.age = 20;
s.display();
\end{verbatim}

In this example, the object \texttt{s} accesses its data members and member function using the dot operator.

\newpage
\subsubsection{Accessing Members Using the Pointer Operator (\texttt{->})}

The pointer operator (\texttt{->}) is used to access the members of an object through a pointer. It combines pointer dereferencing and member access into a single operator.

\begin{itemize}
    \item Used with pointers to objects.
    \item Accesses both data members and member functions.
    \item Equivalent to dereferencing the pointer and then using the dot operator:
    \begin{verbatim}
(*pointer).member
    \end{verbatim}
\end{itemize}

For example:

\begin{verbatim}
Student* p = new Student;

p->name = "Bob";
p->age = 22;
p->display();

delete p;
\end{verbatim}

The above is equivalent to:

\begin{verbatim}
(*p).name = "Bob";
(*p).age = 22;
(*p).display();
\end{verbatim}

The pointer operator is generally preferred because it is shorter, easier to read, and less prone to errors than explicitly dereferencing the pointer before using the dot operator.

\newpage

\begin{itemize}
	\item \fbox{\textbf{Listing 9.1} A compile-worthy class Human}
\end{itemize}

\newpage
\subsubsection{Class vs Struct}

In C++, both \verb|class| and \verb|struct| can contain data members, member functions, constructors, destructors, and other class features. The primary difference between them is their default access level.

\begin{itemize}
    \item Both \verb|class| and \verb|struct| can contain data members and member functions.
    \item Both support constructors, destructors, inheritance, and polymorphism.
    \item The default access specifier for a \verb|class| is \verb|private|.
    \item The default access specifier for a \verb|struct| is \verb|public|.
    \item The default inheritance mode for a \verb|class| is \verb|private| inheritance.
    \item The default inheritance mode for a \verb|struct| is \verb|public| inheritance.
    \item By convention, \verb|struct| is typically used for simple data structures, while \verb|class| is used for objects that encapsulate both data and behavior.
\end{itemize}

For example, the following class has private members by default:

\begin{verbatim}
class Student {
    int age;    // private by default
};
\end{verbatim}

To make members accessible, the \verb|public| access specifier must be used:

\begin{verbatim}
class Student {
public:
    std::string name;
    int age;
};
\end{verbatim}

In contrast, a structure has public members by default:

\begin{verbatim}
struct Student {
    std::string name;   // public by default
    int age;
};
\end{verbatim}

Thus, although \verb|class| and \verb|struct| are nearly identical in functionality, they differ primarily in their default access and inheritance rules, as well as their typical usage conventions.

\newpage

\subsubsection{Different Ways to Declare a \texttt{struct}}

\begin{enumerate}
    \item \textbf{Define the structure and declare an object immediately:}

\begin{verbatim}
struct Student {
    std::string name;
    int age;
} student1;
\end{verbatim}

    Here, the structure type \verb|Student| is defined, and an object named
    \verb|student1| is created immediately.

    \item \textbf{Define the structure first, then declare objects later (most common):}

\begin{verbatim}
struct Student {
    std::string name;
    int age;
};

Student student1;
Student student2;
\end{verbatim}

    This is the preferred style because the type can be reused to create
    multiple objects.

    \item \textbf{Declare multiple objects when defining the structure:}

\begin{verbatim}
struct Student {
    std::string name;
    int age;
} student1, student2, student3;
\end{verbatim}

    Multiple objects can be declared immediately after the structure
    definition.

    \item \textbf{Declare a pointer to a structure:}

\begin{verbatim}
struct Student {
    std::string name;
    int age;
};

Student* ptr = new Student;
\end{verbatim}

    The pointer \verb|ptr| stores the address of a dynamically allocated
    \verb|Student| object.

    \item \textbf{Anonymous structure (rarely used):}

\begin{verbatim}
struct {
    std::string name;
    int age;
} student1;
\end{verbatim}

    This creates an object named \verb|student1| whose type has no name.
    Since the structure is anonymous, additional objects of the same type
    cannot be declared later.
\end{enumerate}

\newpage

\paragraph{Note on \texttt{typedef struct}}

In the C programming language, structures are commonly declared using
\verb|typedef| so that the \verb|struct| keyword does not need to be written
when declaring variables.

For example:

\begin{verbatim}
typedef struct Student {
    char name[50];
    int age;
} Student;

Student s1;
\end{verbatim}

It is also common to define an anonymous structure using \verb|typedef|:

\begin{verbatim}
typedef struct {
    char name[50];
    int age;
} Student;

Student s1;
\end{verbatim}

In C++, however, this is unnecessary because the structure name is
automatically introduced as a type name:

\begin{verbatim}
struct Student {
    std::string name;
    int age;
};

Student s1;
\end{verbatim}

Although \verb|typedef struct| is valid C++, it is considered redundant and is
rarely used in modern C++ code.

\newpage
\subsubsection{Public, Private, Protected, Friend in C++}

Access specifiers in C++ control the visibility and accessibility of class members. They are essential for encapsulation and data hiding.

\begin{itemize}
    \item \textbf{public}
    \begin{itemize}
        \item Members declared as \verb|public| are accessible from anywhere in the program.
        \item They form the interface of the class.
    \end{itemize}

    \item \textbf{private}
    \begin{itemize}
        \item Members declared as \verb|private| are only accessible within the same class.
        \item They are used to hide implementation details.
    \end{itemize}

    \item \textbf{protected}
    \begin{itemize}
        \item Members declared as \verb|protected| are accessible within the class and its derived (child) classes.
        \item They are not accessible from outside the class hierarchy.
    \end{itemize}

    \item \textbf{friend}
    \begin{itemize}
        \item The \verb|friend| keyword allows a non-member function or another class to access private and protected members of a class.
        \item Friendship is granted explicitly and is not inherited or transitive.
    \end{itemize}
\end{itemize}

\paragraph{Example:}

\begin{verbatim}
#include <iostream>
using namespace std;

class Base {
private:
    int privateValue = 10;

protected:
    int protectedValue = 20;

public:
    int publicValue = 30;

    void show() {
        cout << privateValue << " "
             << protectedValue << " "
             << publicValue << endl;
    }

    friend void friendFunction(Base obj);
};

void friendFunction(Base obj) {
    // Can access private and protected members
    cout << obj.privateValue << endl;
    cout << obj.protectedValue << endl;
    cout << obj.publicValue << endl;
}

class Derived : public Base {
public:
    void access() {
        // cout << privateValue;   // not accessible
        cout << protectedValue << endl; // accessible
        cout << publicValue << endl;    // accessible
    }
};

int main() {
    Base b;
    b.publicValue = 100;

    b.show();
    friendFunction(b);

    Derived d;
    d.access();

    return 0;
}
\end{verbatim}

\paragraph{Key Summary:}

\begin{itemize}
    \item \verb|public| → accessible everywhere
    \item \verb|private| → accessible only inside the class
    \item \verb|protected| → accessible in class and derived classes
    \item \verb|friend| → grants special access to external functions/classes
\end{itemize}

\newpage
\subsubsection{This Keyword}

The \verb|this| keyword in C++ is an implicit pointer available inside all non-static member functions. It points to the object that invoked the member function.

\begin{itemize}
    \item \verb|this| is a pointer to the current object.
    \item It is automatically passed to all non-static member functions.
    \item It is used to access members of the current object explicitly.
    \item It helps resolve ambiguity when parameter names shadow member variables.
    \item It cannot be used in static member functions.
\end{itemize}

\paragraph{Example:}

\begin{verbatim}
#include <iostream>
using namespace std;

class Student {
private:
    int age;

public:
    void setAge(int age) {
        // 'this->age' refers to the member variable
        // 'age' refers to the parameter
        this->age = age;
    }

    void display() {
        cout << "Age: " << this->age << endl;
    }
};

int main() {
    Student s;
    s.setAge(20);
    s.display();

    return 0;
}
\end{verbatim}

\paragraph{Key Points:}

\begin{itemize}
    \item \verb|this| is an implicit pointer of type \verb|ClassName* const|.
    \item It always points to the current object instance.
    \item It is commonly used to disambiguate member variables from parameters.
    \item It improves readability in method chaining scenarios.
\end{itemize}

\newpage
\subsubsection{Static Variables in Classes}

A \verb|static| data member in a class is shared among all objects of that class rather than each object having its own separate copy.

\begin{itemize}
    \item A \verb|static| member variable belongs to the class, not individual objects.
    \item Only one copy of the variable exists, regardless of how many objects are created.
    \item It is initialized outside the class definition.
    \item It can be accessed using the class name or through an object.
    \item It retains its value throughout the lifetime of the program.
\end{itemize}

\paragraph{Example:}

\begin{verbatim}
#include <iostream>
using namespace std;

class Counter {
public:
    static int count;

    Counter() {
        count++;
    }

    void display() {
        cout << "Count: " << count << endl;
    }
};

// Definition and initialization of static member
int Counter::count = 0;

int main() {
    Counter c1;
    Counter c2;
    Counter c3;

    c1.display();
    c2.display();
    c3.display();

    cout << "Final count: " << Counter::count << endl;

    return 0;
}
\end{verbatim}

\paragraph{Key Points:}

\begin{itemize}
    \item Static members are shared across all instances of a class.
    \item They must be defined outside the class using scope resolution operator \verb|::|.
    \item They are useful for tracking shared state (e.g., object count, configuration values).
\end{itemize}


\newpage

\subsection{Constructors}
\begin{itemize}
	\item A \vocab{constructor} is a \textbf{special function} (or method) invoked durinkg the \textbf{instantiation} of a \textbf{class} to construct an \textbf{object}.
\end{itemize}

\pbreak

\subsubsection{Declaring and Implementing a Constructor}
\begin{itemize}
	\item A \textbf{constructor} is a \textbf{special function} that takes the name of the class and returns \textbf{no value}.
	\item So, \verb|class Human| would have a constructor that is declared like this:
		\begin{verbatim}
class Human
{
public:
	Human();	// declaration of a constructor
}
		\end{verbatim}

	\pbreak

\item This constructor can be implemented either \textbf{inline} within the class or \textbf{externally} outside the class declaratiopn.

\item An \vocab{implementation} (also called \vocab{definition}) inside the class looks like this:

\begin{verbatim}
class Human
{
public:
	Human()
	{
		// constructor code here
	}
};
\end{verbatim}

\item A variant enabling you to define the constructor \textbf{outside} the class's declaration looks like this:

\begin{verbatim}
class Human
{
public:
	Human();	// constructor declaration
};

// constructor implementation (definition)
Human::Human()
{
	// constructor code here
}
\end{verbatim}


\item the \verb|::| is called the \vocab{scope resolution operator}.

	For example, \verb|Human::dateOfBirth| is referring to the \textbf{variable} \verb|dateOfBirth| declared \underline{\textbf{within the scope}} of class \verb|Human|.

\item \verb|::dateOfBirth| would refer to the \verb|dateOfBirth| in the \underline{\textbf{global scope}}.
\end{itemize}

\newpage

\subsubsection{When and How to Use Constructors}
\begin{itemize}
	\item An \textbf{constructor} is always invoked during \textbf{object creation}, when an \textbf{instance} of a class is constructed.
	\item This is a perfect time to initialize the class members of your instance.

	\item \fbox{\textbf{Listing 9.3} Using Constructors}
\end{itemize}

\pbreak

\subsubsection{Overloading Constructors}
\begin{itemize}
	\item Constructors can be \textbf{overloaded} (mean multiple constructors with different parameters) just like functions.
	\item Example:
	\begin{verbatim}
class Human
{
public:
	Human() {...} // default constructor code

	Human(string humansName) {...} // overloaded constructor code 
};
	\end{verbatim}

\item \fbox{\textbf{Listing 9.4} Multiple Constructors}

\item You can choose to \textbf{not} implement the \vocab{default constructor} to enforce \textbf{object instantiation} with certain minimal parameters as explained in the next section.
\end{itemize}

\pbreak

\subsubsection{Class Without a Default Constructor}
\begin{itemize}
\item \fbox{\textbf{Listing 9.4} Class with overload constructor and No Default Constructor}

\item You also need not to declare any constructor since the default one is always provided implicitly
\end{itemize}

\newpage
\subsubsection{Constructor Parameters with Default Values}
\begin{itemize}
	\item Constructors can have \textbf{default values} like functions, allowing objects to be created with fewer arguments while still initializing members safely.

	\begin{verbatim}
#include <iostream>
#include <string>

using namespace std;

class Student {
private:
    string name;
    int age;
    double gpa;

public:
    // Constructor with default parameter values
    Student(string n = "Unknown", int a = 0, double g = 0.0) {
        name = n;
        age = a;
        gpa = g;
    }

    void display() {
        cout << "Name: " << name << endl;
        cout << "Age: " << age << endl;
        cout << "GPA: " << gpa << endl;
    }
};

int main() {
    Student s1;                      // uses all defaults
    Student s2("Alice");            // only name provided
    Student s3("Bob", 22);         // name + age
    Student s4("Charlie", 21, 3.9);// full initialization

    s1.display();
    cout << "-----------------" << endl;

    s2.display();
    cout << "-----------------" << endl;

    s3.display();
    cout << "-----------------" << endl;

    s4.display();

    return 0;
}
	\end{verbatim}
\end{itemize}

\newpage
\subsubsection{Constructors with Initialization Lists}

An initialization list initializes a class's data members before the constructor body is executed. This is the preferred method of initializing member variables in modern C++.

\begin{itemize}
    \item An initialization list appears after the constructor parameter list and before the constructor body.
    \item Data members are initialized directly rather than first being default-constructed and then assigned values.
    \item Initialization lists are generally more efficient than assigning values inside the constructor body.
    \item They are \textbf{required} when initializing:
    \begin{itemize}
        \item \verb|const| data members.
        \item Reference (\verb|&|) data members.
        \item Class members that do not have a default constructor.
    \end{itemize}
    \item Data members are initialized in the order they are declared in the class, \textbf{not} in the order they appear in the initialization list.
\end{itemize}

\paragraph{Syntax:}

\begin{verbatim}
ClassName(type parameter)
    : member(parameter)
{
    // Constructor body
}
\end{verbatim}

\paragraph{Example:}

\begin{verbatim}
#include <iostream>
#include <string>

using namespace std;

class Student {
private:
    string name;
    int age;
    double gpa;

public:
    Student(string n, int a, double g)
        : name(n), age(a), gpa(g)
    {
    }

    void display() {
        cout << "Name: " << name << endl;
        cout << "Age: " << age << endl;
        cout << "GPA: " << gpa << endl;
    }
};

int main() {
    Student s1("Alice", 20, 3.8);
    s1.display();

    return 0;
}
\end{verbatim}

\paragraph{Example Requiring an Initialization List:}

\begin{verbatim}
#include <iostream>

using namespace std;

class Example {
private:
    const int id;
    int& value;

public:
    Example(int i, int& v)
        : id(i), value(v)
    {
    }

    void display() {
        cout << "ID: " << id << endl;
        cout << "Value: " << value << endl;
    }
};

int main() {
    int x = 100;

    Example obj(1, x);
    obj.display();

    return 0;
}
\end{verbatim}

In the second example, the initialization list is mandatory because both
\verb|id| (a \verb|const| member) and \verb|value| (a reference member) must
be initialized when the object is created. They cannot be assigned values
inside the constructor body.

\newpage
\subsection{Destructor}
\begin{itemize}
	\item A \vocab{destructor}, like a \textbf{constructor}, is a \textbf{special function}.
	\item A \textbf{constructor} is invoked at \textbf{object instantiation}, and a \fbox{\textbf{destructor}} is \underline{automatically} invoked when an \textbf{object} is \textbf{destroyed}.
\end{itemize}

\pbreak

\subsubsection{Declaring and Implementing a Destructor}
\begin{itemize}
	\item The \textbf{destructor} looks like a \textbf{function} that takes the name of the \textbf{class}, yet has a \textbf{tilde (\textasciitilde{})}.

	\item Example:
	\begin{verbatim}
class Human 
{
public:
	~Huamn();	// declaration of a destructor
};
	\end{verbatim}

	\pbreak

\item The \textbf{destructor} can be implemented \textbf{inline} in the class or \textbf{externally} outside the class declaration.

	\fbox{internally}:
\begin{verbatim}
class Human
{
public:	
	~Human()
	{
		// destructor code here
	}
};

\end{verbatim}

\fbox{externally}:
\begin{verbatim}
class Human
{
public:
	~Human();	// destructor declaration
};

// destructor defintion (implementation)
Human::~Human()
{
	// destructor code here
}

\end{verbatim}

\end{itemize}

\newpage

\subsubsection{When and How to Use a Destructor}
\begin{itemize}
	\item A \textbf{destructor} is always invoked when an \textbf{object} of a class is \textbf{destroyed} when it \textbf{goes out of scope} or is \textbf{deleted} via \verb|delete| or \verb|delete[]| or free, etc.

	\item This property makes a \textbf{destructor} the ideal place to reset variables and \textbf{release dynamically allocated memory} and other resources.	 

\begin{itemize}
    \item \verb|std::string| automatically manages dynamic memory for its character data, relieving the programmer from manually allocating and deallocating memory.
    \item The \verb|std::string| object itself is stored wherever it is created (e.g., on the stack for local variables or on the heap if dynamically allocated).
    \item The character data managed by \verb|std::string| may be stored within the object itself (using Small String Optimization) or in dynamically allocated heap memory, depending on the implementation and the length of the string.

	\item But internally it would use constructors and destructors and operators.
\end{itemize}

\pbreak
\item \fbox{\textbf{Listing 9.7} Destructor}

\item Smart pointers (Lesson 26) is where destructors play a critical role in working with pointers in a smarter way.

\item \textbf{Destructors} \underline{cannot be overloaded}.

\item A class can only have \textbf{one destructor}, the compiler creates and invokes a dummy destructor, that is, an empty one (that does no cleanup of dynamically allocated memory)

\end{itemize}

\newpage

\subsection{Copy Constructors}

\begin{itemize}
	\item Argument passing to functions copy the argument into the function stack frame whenevery the function is called and of course if it has arguments passed to it.

	\item This rule applies to \textbf{objects}, that is, instances of classes as well.
\end{itemize}

\pbreak

\subsubsection{Shallow Copying and Associated Problems}
\begin{itemize}
	\item Classes such as \verb|MyString| in \textbf{Listing 9.7} had a \textbf{pointer member} \verb|buffer| that points to \textbf{dynamically allocated memory}, allocated in the constructor using \verb|new| and then deallocated in the destructor using \verb|delete[]|.

\item When an \textbf{object} of this class is \textbf{copied}, the \textbf{pointer member} is \underline{\textbf{copied}}, but not the \textbf{pointer memory}, resulting in 2 objects pointing to the \textbf{same dynamically allocated buffer in memory}.

	So,

	\begin{itemize}
    \item Classes such as \verb|MyString| in \textbf{Listing 9.7} contain a pointer member (\verb|buffer|) that points to dynamically allocated memory. The memory is allocated in the constructor using \verb|new| and released in the destructor using \verb|delete[]|.

    \item When an object is copied, the compiler performs a \textbf{shallow copy} by default. This copies the \textbf{pointer value} (memory address), but \textbf{does not allocate or copy the dynamically allocated memory} it points to.

    \item As a result, both objects' pointer members refer to the \textbf{same block of dynamically allocated memory}. If one object modifies or deallocates the memory, it affects the other object, potentially leading to errors such as double deletion.
\end{itemize}

\item Hence, when an \textbf{object} is \textbf{destructed}, \verb|delete[]| deallocates the memory, thereby \textbf{invalidating} the \textbf{pointer copy} held by the other object.

\item Such copies are \vocab{shallow} and are a \textbf{threat} to the stability of the program, as \fbox{\textbf{Listing 9.8} Shallow copy Problem}

\end{itemize}
	\newpage
	Here is another small example:
\begin{verbatim}
#include <iostream>
#include <cstring>

using namespace std;

class MyString {
private:
    char* buffer;

public:
    MyString(const char* str) {
        buffer = new char[strlen(str) + 1];
        strcpy(buffer, str);
    }

    ~MyString() {
        delete[] buffer;
    }

    void display() {
        cout << buffer << endl;
    }
};

int main() {
    MyString s1("Hello");

    // Default copy constructor performs a shallow copy
    MyString s2 = s1;

    s1.display();
    s2.display();

    // When main() ends:
    // ~MyString() is called for s2 -> delete[] buffer;
    // ~MyString() is called for s1 -> delete[] buffer; (same memory!)
    // Result: Undefined behavior (double deletion)

    return 0;
}
\end{verbatim}

	This is a classic example demonstrating why the default copy constructor can be dangerous for classes that own dynamically allocated memory, and why a deep copy (custom copy constructor) is needed.

	\pbreak

	Output for listing 9.8:
	\begin{verbatim}
String buffer in sayHello is 23 characters long
Buffer contains: Hello from String Class
Invoking destructor, clearing up for id: 21
Invoking destructor, clearing up for id: 0
free(): double free detected in tcache 2
Aborted                    (core dumped) ./a.out
	\end{verbatim}

	where i changed the id in the callee to 21 and that shows the passed arg first destructs;

	\begin{itemize}
		\item So, the issue is both objects point to the same memory address, but when they both destructor we get the error of freeing an already deleted/freed data
	\end{itemize}

\newpage
\subsubsection{Ensuring Deep Copy Using a Copy Constructor}
\begin{itemize}
	\item The \textbf{copy constructor} is an \textbf{\underline{overloaded constructor}} that you supply.

	\item It is invoked everytime an \textbf{object} of the class is \textbf{copied}.
	\item So, we shall write a \vocab{copy constructor} to help solve the issues with \textbf{shallowing copying}, which we call \vocab{deep copying}.

		\pbreak

	\item The \textbf{declaration syntax} of a \textbf{copy constructor} for \verb|class MyString| is the following:

		\begin{verbatim}
class MyString
{
	MyString(const MyString& copySource);	// copy constructor
	// the const says "I promise not to modify the object I'm copying from."
};

another reason for the const:
MyString::MyString(const MyString& copySource) {
    // OK: read from copySource
    std::cout << copySource.buffer << std::endl;

    // Error: cannot modify the source object
    // copySource.buffer = nullptr;
}

...

MyString::MyString(const MyString& copySource) 
{
	// copy constructor implementation code
}
		\end{verbatim}

\pbreak

\item Thus, a \vocab{copy constructor} takes an \textbf{object} of the \textbf{same class} by \textbf{reference} as a parameter.
\item This parameter is an \textbf{alias} of the \underline{\textbf{source object}} and is the handle you have in writing your \textbf{custom copy code}.

\item You would use the \textbf{copy constructor} to \textbf{ensure} a \vocab{deep copy} of \textbf{all buffers} in the source.

	Like before you would only just copy the memory address, but here you would \textbf{allocate new memory on the heap} and then \textbf{copy the actual data contents} from the source object into this newly allocated memory, so that \textbf{both} objects \textbf{own independent copies} of the \textbf{data} \underline{rather than} sharing the \textbf{same memory address}.

\item \fbox{\textbf{Listing 9.9} Define a copy constructor to ensure deep copy of dynamically allocated buffers}.

	\pbreak

\item The \textbf{syntax difference} between a \textbf{regular constructor} and a \textbf{copy constructor} is in their \textbf{parameter lists}.

	A regular constructor takes \textbf{raw values} (e.g., integers, strings, pointers), while a copy constructor takes a \textbf{reference to another object of the same class}.

    \begin{verbatim}
    // Regular constructor
    MyString(const char* str);

    // Copy constructor
    MyString(const MyString& other);
    \end{verbatim}

    \item The copy constructor is used to create a new object from an existing object of the same type, typically performing a \textbf{deep copy} of dynamically allocated resources.

	\item Yes, the \textbf{\fbox{copy constructor is invoked \underline{implicitly}}} in several situations where an object is copied. These include:
    \begin{itemize}
        \item When an object is initialized from another object:
        \begin{verbatim}
        MyString s2 = s1;
        \end{verbatim}

        \item When an object is passed to a function by value:
        \begin{verbatim}
        void func(MyString s);
        func(s1);
        \end{verbatim}

        \item When an object is returned by value from a function:
        \begin{verbatim}
        MyString func();
        \end{verbatim}
    \end{itemize}

    \item If no copy constructor is defined, C++ generates a default one that performs a \textbf{shallow copy}, meaning it copies pointer values rather than the data they point to.

	\item A \textbf{copy constructor} can \textbf{only take} \textbf{\underline{one argument}}, which is a \textbf{reference to an object} of the \textbf{same class} (typically \verb|const ClassName&|). 

		This is because its purpose is to create a \textbf{new object as a copy} of an \textbf{existing object}.

    \item If more than one parameter is used, it is no longer considered a copy constructor but instead a regular overloaded constructor.

    \item The standard form is:
    \begin{verbatim}
    ClassName(const ClassName& source);
    \end{verbatim}


\item The compiler does not "guess" or manually pass the object. Instead, when a copy constructor is invoked, the compiler automatically selects it based on the context where a copy is required (e.g., initialization from another object, passing by value, or returning by value).

Internally, the compiler rewrites the code so that the source object is passed as the single argument to the copy constructor. For example:

\begin{verbatim}
MyString s2 = s1;
\end{verbatim}

is conceptually transformed by the compiler into:

\begin{verbatim}
MyString s2(s1);
\end{verbatim}

Here, \verb|s1| is explicitly passed as the argument to the copy constructor. The "implicit" part refers to the fact that the programmer does not explicitly call the constructor; the compiler inserts the call automatically based on the context.

\end{itemize}

\newpage


\begin{tcolorbox}
\begin{itemize}
	\item The \textbf{copy constructor} has ensured a \textbf{deep copy} in cases such as:

	\begin{verbatim}
MyString sayHello("Hello");
UseMyString(sayHello);
	\end{verbatim}

\item However, what if you tried copying via \textbf{assignment}:

\begin{verbatim}
MyString overwrite("who cares? ");
overwrite = sayHello;
\end{verbatim}

This would \textbf{still} be a \underline{\textbf{shallow copy}} because you still haven't yet supplied a \vocab{copy assignment operator=}.

\item In the absence of one, the compiler has supplied a default for you that does a \textbf{shallow copy}.

\item The \vocab{copy assignment operator} is discussed in Lesson 12. Listing 12.8 is an improvied \verb|MyString| that implements the same:

	\begin{verbatim}
MyString::operator= (const MyString& copySource)
{
	// ... copy assignment operator code
}
	\end{verbatim}
\end{itemize}
\end{tcolorbox}

\pbreak

\begin{tcolorbox}
	\begin{itemize}
		\item Using \verb|const| in the \textbf{copy constructor declaration} ensures that the \textbf{copy constructor} does \underline{\textbf{not modifiy}} the \textbf{source object} being referred to.

		\item Additionally, the parameter in the \textbf{copy constructor} is \textbf{passed by reference} as a \underline{necessity}.

		If it weren't a reference, the \textbf{copy constructor} would itself invoke a copy, thus invoking itself again and so on till the system runs out of memory.

	\end{itemize}

	\pbreak

	\begin{itemize}
    \item Using \verb|const| in the copy constructor ensures that the source object cannot be modified during copying. This guarantees that the original object remains unchanged and allows copying from both \verb|const| and non-\verb|const| objects.

    \item The copy constructor parameter must be passed by reference. If it were passed by value, a copy would need to be made to pass the argument into the constructor. This would recursively invoke the copy constructor again, leading to infinite recursion and eventually a stack overflow.

    \item Therefore, the correct and standard form of a copy constructor is:
    \begin{verbatim}
    ClassName(const ClassName& other);
    \end{verbatim}
\end{itemize}
\end{tcolorbox}

\newpage
\begin{tcolorbox}
\begin{itemize}
    \item \textbf{Do} always implement a \textbf{copy constructor} and \textbf{copy assignment operator} when your class contains \textbf{raw pointer members} (e.g. \verb|char*|). This prevents shallow copying and avoids issues such as double deletion and dangling pointers.

    \item \textbf{Do} declare the copy constructor using a \textbf{const reference parameter}:
    \begin{verbatim}
    ClassName(const ClassName& other);
    \end{verbatim}
    This ensures the source object is not modified during copying and allows copying from both const and non-const objects.

    \item \textbf{Do} use the \verb|explicit| keyword for constructors that take a single parameter to prevent unintended implicit conversions. For example:
    \begin{verbatim}
    class MyString {
    public:
        explicit MyString(const char* str);
    };

    MyString s1 = "Hello";   // Not allowed (prevents implicit conversion)
    MyString s2("Hello");    // Allowed (explicit construction)
    \end{verbatim}

    \item \textbf{Do} prefer using standard library types such as \verb|std::string| and smart pointers (e.g. \verb|std::unique_ptr|, \verb|std::shared_ptr|) as class members instead of raw pointers, since they automatically manage memory and implement safe copy/move semantics.

    \item \textbf{Don't} use raw pointers as class members unless absolutely necessary, as they require manual memory management and increase the risk of memory leaks, double deletion, and undefined behavior.
\end{itemize}
\end{tcolorbox}

\newpage

\subsubsection{Explicit keyword}
\begin{itemize}
    \item The \verb|explicit| keyword in C++ is used to \textbf{prevent implicit conversions} that the compiler would otherwise perform automatically using constructors.

    \item Without \verb|explicit|, a constructor that takes a single parameter can be used for implicit type conversion. For example:

\begin{verbatim}
class MyString {
public:
    MyString(const char* str) {
        // constructor
    }
};

MyString s = "Hello";   // implicit conversion allowed
\end{verbatim}

    \item In this case, the compiler automatically converts the string literal into a \verb|MyString| object.

    \item When the constructor is marked as \verb|explicit|, implicit conversions are disallowed:

\begin{verbatim}
class MyString {
public:
    explicit MyString(const char* str) {
        // constructor
    }
};

MyString s = "Hello";   // Error: implicit conversion not allowed
MyString s2("Hello");   // OK: explicit construction
\end{verbatim}

    \item The \verb|explicit| keyword is used to prevent unintended or accidental conversions that may lead to bugs or unclear code.

    \item In summary, \verb|explicit| forces the programmer to clearly state when a conversion or constructor call should occur.

	\item \boxednote{
		\vocab{explicit} is a \textbf{compile-time} rule that \textbf{prevents the compiler} from performing \textbf{automatic (implicit) type conversions} using constructors.}
\end{itemize}

\newpage

\subsubsection{Move Constructors Help Improve Performance}
\begin{itemize}
	\item There are cases where \textbf{objects} are subjected to copy steps automatically, due to the nature of the language and its needs.

	\item Consider:

		\begin{verbatim}
class MyString
{
	// pick implementation from listing 9.9
};

MyString Copy(MyString& source)		// function
{
	MyString copyForReturn(source.GetString());		// create copy
	return copyForReturn;	// return by value invokes copy constructor
}

int main()
{
	MyString sayHello("Hi");
	MyString sayHelloAgain(Copy(sayHello));		// invokes 2x copy constructor
	return 0;
}
		\end{verbatim}

		\pbreak

\item As the comment indicates, in the \textbf{instantiation} of \verb|sayHelloAgain|, the \textbf{copy constructor} is invoked twice, thus a \textbf{deep copy} was performed \textbf{twice} b/c our call to function \verb|Copy(sayHello)| returns a \verb|MyString| \textbf{by value}.

\item However, this value returned is \textbf{very temporary} and is \textbf{not avaiable} outside this expression.

	So, the \textbf{copy constructor} invoked in good faith by the C++ compiler is burden on \textbf{performance}.

	This impact becomes significant if our class were to contain objects of great size.

	\pbreak

\item To \textbf{avoid} this performance bottleneck, versions of C++ starting with C++11 feature a \vocab{move constructor} in addition to a \textbf{copy constructor}.

\item The syntax of a \vocab{move constructor} is

	\begin{verbatim}
// move constructor
MyString(MyString&& moveSource)
{
	if (moveSource != NULL)
	{
		buffer = moveSource.buffer;		// take ownership i.e. 'move'
		moveSource.buffer = NULL;		// set the move source to NULL
	}
}
	\end{verbatim}

\begin{itemize}
    \item The constructor
\begin{verbatim}
MyString(MyString&& moveSource)
\end{verbatim}
is called a \textbf{move constructor}.

    \item Here's what each part means:
    \begin{itemize}
        \item \verb|MyString| --- the constructor's name.
        \item \verb|MyString&&| --- an \textbf{rvalue reference} to another \verb|MyString|.
        \item \verb|moveSource| --- the parameter name.
    \end{itemize}

    \item In other words, this constructor says:
    \begin{quote}
    "This constructor takes ownership of another temporary (or movable) \verb|MyString|."
    \end{quote}

    \item Compare the different constructor signatures:
\begin{verbatim}
MyString(MyString other);      // pass by value
MyString(MyString& other);     // lvalue reference
MyString(MyString&& other);    // rvalue reference
\end{verbatim}

    \item They each accept different kinds of arguments.

    \item Normally, variables are \textbf{lvalues}. For example:
\begin{verbatim}
MyString a("hello");
\end{verbatim}

    Since \verb|a| has a name, it is an lvalue.

    \item A temporary object is an \textbf{rvalue}. For example:
\begin{verbatim}
MyString("hello")
\end{verbatim}

    or
\begin{verbatim}
return MyString("hello");
\end{verbatim}

    These objects are about to disappear, so C++ allows their resources to be
    transferred instead of copied.

    \item Example:
\begin{verbatim}
MyString a("hello");
MyString b(std::move(a));   // calls move constructor
\end{verbatim}

    \verb|std::move(a)| converts \verb|a| into an rvalue, so the move constructor
    is selected:
\begin{verbatim}
MyString(MyString&& moveSource)
\end{verbatim}

    \item Inside the move constructor, the implementation may look like:
\begin{verbatim}
MyString::MyString(MyString&& moveSource)
{
    data = moveSource.data;      // steal pointer
    moveSource.data = nullptr;   // leave source safe
}
\end{verbatim}

    No new memory is allocated. The pointer ownership is simply transferred.

    \item Why not use \verb|MyString&|?

    If the constructor were
\begin{verbatim}
MyString(MyString& other)
\end{verbatim}

    then it could receive a normal object that the caller still intends to use.
    Stealing its resources would unexpectedly leave that object in an unusable state.

    \item The \verb|&&| tells the compiler:
    \begin{quote}
    "This object is temporary (or has been explicitly marked movable), so it is
    safe to transfer its resources."
    \end{quote}

    \item Summary:
    \begin{itemize}
        \item \verb|&| = reference to a normal object (do not steal its resources).
        \item \verb|&&| = reference to a temporary or movable object (safe to steal its resources).
        \item \verb|MyString(MyString&& moveSource)| is the constructor that performs a \textbf{move} instead of a copy.
    \end{itemize}
\end{itemize}

\pbreak

\item When a \textbf{move constructor} is programmed, the compiler automatically opts for the same for "moving" the temporary resource and hence \textbf{avoiding} a \textbf{deep-copy} step.

\item With this \textbf{move consutrctor} implemented, the comment should be appropriately changed to the following:
	\begin{verbatim}
MyString sayHelloAgain(Copy(sayHello));	// invokes 1x copy, 1x move constructors
	\end{verbatim}

\item The \vocab{move constructor} is usually implemented with the \vocab{move assigment operator}, which is dicussed in Lesson 12. Listing 12.11 is a better MyString that implements the \textbf{move constructor} and the \textbf{move assignment operator}. 
\end{itemize}

\newpage

\subsection{Different Uses of Constructors and Destructors}
\begin{itemize}
	\item We've learned now concepts that enable you to create classes that can control how they're \textbf{created, copied, destroyed,} or \textbf{expose data (through \vocab{access specificers}}).
\end{itemize}

\pbreak

\subsubsection{Class That Does Not Permit Copying}
\begin{itemize}
	\item Somtimes we don't want duplicates at all in our program:
	\begin{verbatim}
President ourPresident;
DoSomething(ourPresident);		// duplicate created in passing by value
President clone;
clone = curPresident;			// duplicate via assignment
	\end{verbatim}
\item You need to ensure that certain resources \textbf{cannot} be \textbf{copied} or \textbf{duplicated}.

\item If you \textbf{don't} declare a \textbf{copy constructor}, the C++ compiler inserts a \textbf{default public copy constructor} for you.

\item This ruins your design and threatens your implementation.

\item We have a solution.

\item You would ensure that your \verb|class| \textbf{\underline{cannot be copied}} by declaring a \verb|private| \textbf{copy constructor}.

\item This ensures that the function call \verb|DoSomething(ourPresident)| will cause a compiler failure.

\item To avoid \textbf{assignment}, you declare a \verb|private| \textbf{assignment operator}.

	\pbreak

\item Thus, the solution is the following:

	\begin{verbatim}
class President
{
private:
	President(const President&);				// private copy constructor
	President& operator= (const President&);	// private copy assignment operator

	// ... other attributes
};
	\end{verbatim}

\item There is \textbf{no need for implementation} of the \textbf{private copy constructor} or \textbf{assignment operator}.

\item Just \underline{\textbf{declaring them private}} is adequate and \textbf{sufficient} toward fulfilling your goal of ensuring \textbf{non-copyable objects} of \verb|class President|.
\end{itemize}

\newpage

\subsubsection{Singleton Class That Permits a Single Instance}
\begin{itemize}
\item \verb|class President| earlier is good, but has a \textbf{shortcoming}:\\
	\underline{It \textbf{cannot help} creation of multiple presidents via instantiation of \textbf{multiple objects}:}
\begin{verbatim}
President One, Two, Three;
\end{verbatim}

\item Individually, they \underline{are} \textbf{non-copyable} thanks to the \textbf{private copy constructors}, but what youi ideally need is a \verb|class President| that has \textbf{one}, and \underline{only one}, real-world manifestation --- that is, there is \underline{\textbf{only one object}} and creation of additional ones is prohibited. 

\item Welcome to the concept of \vocab{singleton} that uses \vocab{private constructors}, a \vocab{private assignment operator}, and a \vocab{static instance member} to create this (controversally) powerful pattern.
\end{itemize}

\pbreak

\begin{tcolorbox}
	\begin{itemize}
		\item When the keyword \vocab{static} is used on a \textbf{class's data member}, it ensures that the member is \textbf{\underline{shared}} across \textbf{all instances}

		\item When \vocab{static} is used on a \textbf{local variable} declared within the scope of a function, it ensures that the \textbf{variable retains its value} between \textbf{functions calls}.

		\item When \vocab{static} is used on a \textbf{member function} --- a \textbf{method} --- the method is \textbf{shared} across \textbf{all instances} of the class.
	\end{itemize}
\end{tcolorbox}

\pbreak
\begin{itemize}
	\item \fbox{\textbf{Listing 9.10} Singleton class President}
\end{itemize}

\newpage
\subsection{Class That Prohibts Instantiation on the Stack}
\begin{itemize}
	\item Space on the stack is often limited
	\item If you are writing a database that may contain \textbf{terabytes} of data in its internal structures, you might want to ensure that a client of this \textbf{class} \underline{\textbf{cannot instantiate}} it on the \textbf{stack};

		instead it is \textbf{forced} to create \textbf{instances} \underline{only} on the free store.

	\item The key to ensuring this is \textbf{declaring} the \textbf{destructor} \verb|private|:

		\begin{verbatim}
class MonsterDB
{
private:
	~MonsterDB();		// private destructor
	// ... members that consume a huge amount of data
};
		\end{verbatim}

	\item \textbf{Declaring} a \textbf{private destructor} ensures that one is \textbf{not allowed} to create an \textbf{instance} like this:
	\begin{verbatim}
int main()
{
	MonsterDB myDatabase;	// compile error
	// ... more code
	return 0;
}
	\end{verbatim}

	\item This \textbf{instance}, if successfully constructed, would be on the \textbf{stack}.
	\item \textbf{All objects} on the \textbf{stack} get popped when the \textbf{stack} is unwound and therefore the compiler would need to compile and invoke the \textbf{destructor} \verb|~MonsterDB()| at the \textbf{end} of \verb|main()|.

	\item However, this \textbf{destructor} is \textbf{\underline{private}} and there \underline{\textbf{inaccessible}}, resulting in a compile failure.

		\pbreak

	\item A \textbf{private destructor} would \textbf{not stop} you from instantiating on the \vocab{heap}:

		\begin{verbatim}
int main()
{
	MonsterDB* myDatabase = new MonsterDB();	// no error
	// ... more code
	return 0;
}
		\end{verbatim}

	\item We have a memory leak here: as the \textbf{destructor} is \underline{\textbf{not accessible}} from \textbf{main}, you \underline{\textbf{cannot}} do a \verb|delete|, either.

	\item What \verb|class MonsterDB| needs to support is a \verb|public static| \textbf{member function} that would \textbf{destroy} the \textbf{instance} (a class member would have access to the \textbf{private destructor}.)

	\item \fbox{\textbf{Listing 9.11} Only allowed creation on heap}
\end{itemize}

\newpage

\subsubsection{Using Constructors to Convert Types}
\begin{itemize}
	\item Earlier we've learned \textbf{constructors} can be \textbf{overloaded}, that is, they may take \textbf{one or more parameters}.

	\item This feature is often used to \textbf{convert} one type to another.

	\item Let's see a \verb|class Human| that features an \textbf{overloaded constructor} that accepts an integer.

		\begin{verbatim}
class Human
{
	int age;
public:
	Human(int humansAge) : age(humansAge) {}
};


// function that takea  human as a parameter
void DoSomething(Human person)
{
	cout << "Human sent did something" << endl;
	return;
}
		\end{verbatim}

	\item The \textbf{constructor} allows a \textbf{conversion}:

	\begin{verbatim}
Human kid(10); // convert integer into a Human
DoSomething(kid);
	\end{verbatim}


\myhrule{1em}

\item I think the textbook messed here/is misleading so here is update:
\begin{verbatim}
No conversion (explicit construction)
Human kid(10);
DoSomething(kid);
\end{verbatim}

\begin{itemize}
    \item Yes. If you write:

\begin{verbatim}
Human kid(10);
DoSomething(10);
\end{verbatim}

    the \verb|Human| object named \verb|kid| is \textbf{not} used.

    \item Instead, the compiler creates a \textbf{completely separate temporary} \verb|Human|
    object for the function call.

    \item Conceptually, it is as if you had written:

\begin{verbatim}
Human kid(10);      // Your object

Human temp(10);     // Compiler-created temporary
DoSomething(temp);
\end{verbatim}

    \item Therefore, there are \textbf{two different} \verb|Human| objects:
    \begin{itemize}
        \item \verb|kid| --- the object that you explicitly created.
        \item \verb|temp| --- a temporary object automatically created by the compiler.
    \end{itemize}

    \item If, instead, you write:

\begin{verbatim}
Human kid(10);
DoSomething(kid);
\end{verbatim}

    then there is \textbf{no type conversion}. The function receives the object
    \verb|kid| (or a copy of it, since the parameter is passed by value).

    \item Compare the two cases.

    \begin{itemize}
        \item \textbf{Case 1: No conversion}

\begin{verbatim}
Human kid(10);
DoSomething(kid);
\end{verbatim}

        The object flow is:

\begin{verbatim}
kid  ------------------> DoSomething
\end{verbatim}

        \item \textbf{Case 2: Implicit conversion}

\begin{verbatim}
DoSomething(10);
\end{verbatim}

        The object flow is:

\begin{verbatim}
10
 |
 v
Human(10)   <-- temporary object created by the compiler
 |
 v
DoSomething
\end{verbatim}
    \end{itemize}

    \item Therefore, the temporary object is \textbf{not} the same object that you
    previously created. It is a brand new object that exists only for the duration
    of the function call and is destroyed immediately afterward.

    \item This automatic creation of a temporary object is what allows a
    \textbf{single-argument constructor} to perform an \textbf{implicit type conversion}.
\end{itemize}

\item We fix this via the keyword \vocab{explicit}
\end{itemize}

\myhrule{1em}

\pbreak

\begin{tcolorbox}
	\begin{itemize}
		\item Such \textbf{converting constructors} allow \vocab{implicit conversions}.

		\begin{verbatim}
Human anotherKid = 11; // int converted to Human
DoSomething(10);		// 10 converted to Human!
		\end{verbatim}

	\item We declared \verb|DoSomething(Human person)| as a \textbf{function} that accepts a parameter of \textbf{type} \verb|Human| and \underline{does not accept} an \verb|int|!

	\item So why did the line work?

	\item The compiler knows that \verb|class Human| \textbf{supports} a \textbf{constructor} that \underline{accepts an integer} and performed \textbf{implicit conversion}!

		It created an \textbf{object} of type \verb|Human| using the integer you supplied and sent it as an argument to the function.

		\pbreak

	\item To \textbf{void implicit conversions}, use keyword \vocab{explicit} at the time of \textbf{declaring the constructor}.


		\begin{verbatim}
class Human
{
	int age;
public:
	explicit Human(int humansAge) : age(humansAge) {}
};
\end{verbatim}

\item Using \verb|explicit| is \textbf{not} a prerequisite but in many cases a \textbf{good programming practice}.

	\item \fbox{\textbf{Listing 9.12} Using explicit to block implicit conversions}
	\end{itemize}
\end{tcolorbox}

\begin{tcolorbox}
	\begin{itemize}
		\item The problem of \textbf{implicit conversion} and \textbf{avoiding} them using keyword \vocab{explicit} applies to \underline{\textbf{operators}} too.
	\end{itemize}
\end{tcolorbox}

\newpage

\subsection{sizeof() a Class}
\begin{itemize}
    \item The \verb|sizeof| operator returns the number of \textbf{bytes} occupied by an object or a type in memory.

    \item When applied to a class, \verb|sizeof| returns the total amount of memory required to store all of its \textbf{data members}, including any \textbf{padding bytes} that the compiler inserts for memory alignment.

    \item The size of a class \textbf{does not} include the code for its member functions. Member functions are stored only once in the program's executable, regardless of how many objects of the class are created.

    \item A class may contain:
    \begin{itemize}
        \item Fundamental data types (such as \verb|int|, \verb|char|, and \verb|double|).
        \item Objects of other classes.
        \item Structures (\verb|struct|).
        \item Arrays and pointers.
    \end{itemize}

    \item When a class contains another class or a structure as a data member, the
    size of the containing class includes the size of those objects as well.

    \item Consider the following example:

\begin{verbatim}
struct Address
{
    int houseNumber;
    char streetInitial;
};

class Birthday
{
public:
    int day;
    int month;
    int year;
};

class Student
{
private:
    int id;             // Fundamental type
    double gpa;         // Fundamental type
    char grade;         // Fundamental type

    Address address;    // Structure
    Birthday birth;     // Another class
};
\end{verbatim}

    \item The size of a \verb|Student| object is conceptually the sum of the sizes
    of all its data members:

\begin{verbatim}
sizeof(Student)
    =
sizeof(int)
+ sizeof(double)
+ sizeof(char)
+ sizeof(Address)
+ sizeof(Birthday)
+ padding (if required)
\end{verbatim}

    \item The compiler may insert \textbf{padding bytes} between members so that
    each data member begins at an address suitable for efficient access by the CPU.

    \item Therefore, the value returned by \verb|sizeof(Student)| is often
    \textbf{larger} than simply adding the sizes of the individual data members.

    \item For example:

\begin{verbatim}
#include <iostream>
using namespace std;

int main()
{
    cout << "sizeof(Address)  = "
         << sizeof(Address) << endl;

    cout << "sizeof(Birthday) = "
         << sizeof(Birthday) << endl;

    cout << "sizeof(Student)  = "
         << sizeof(Student) << endl;

    return 0;
}
\end{verbatim}

    \item The exact output depends on:
    \begin{itemize}
        \item The compiler being used.
        \item The target architecture (32-bit or 64-bit).
        \item The alignment rules of the system.
    \end{itemize}

    \item In general,

\begin{verbatim}
sizeof(class) =
sum of the sizes of its data members
+ any compiler-added padding
\end{verbatim}

    \item Remember that \verb|sizeof| measures only the memory occupied by the
    object's data. It does \textbf{not} include memory used by member functions,
    nor does it include memory dynamically allocated through pointers.
\end{itemize}

\newpage

So:
\begin{verbatim}
Typical output (64-bit system):

sizeof(Address)  = 8
sizeof(Birthday) = 12
sizeof(Student)  = 40
\end{verbatim}

\subsubsection*{Explanation of the output}

\begin{itemize}
    \item Assume typical sizes:
    \begin{itemize}
        \item \verb|int| = 4 bytes
        \item \verb|double| = 8 bytes
        \item \verb|char| = 1 byte
    \end{itemize}

    \item \textbf{Address struct:}
\begin{verbatim}
struct Address
{
    int houseNumber;      // 4 bytes
    char streetInitial;   // 1 byte
};
\end{verbatim}

    Memory layout:
\begin{verbatim}
houseNumber     4 bytes
streetInitial   1 byte
padding         3 bytes
-----------------------
Total           8 bytes
\end{verbatim}

    \item \textbf{Birthday class:}
\begin{verbatim}
class Birthday
{
public:
    int day;
    int month;
    int year;
};
\end{verbatim}

    Memory layout:
\begin{verbatim}
day     4 bytes
month   4 bytes
year    4 bytes
----------------
Total   12 bytes
\end{verbatim}

    \item \textbf{Student class:}
\begin{verbatim}
class Student
{
private:
    int id;             // 4 bytes
    double gpa;         // 8 bytes
    char grade;         // 1 byte

    Address address;    // 8 bytes
    Birthday birth;     // 12 bytes
};
\end{verbatim}

    Memory layout (typical alignment):
\begin{verbatim}
id              4 bytes
padding         4 bytes   // align double

gpa             8 bytes

grade           1 byte
padding         7 bytes   // align Address

address         8 bytes

birth           12 bytes

padding         4 bytes   // final alignment
-----------------------------------------------
Total          40 bytes
\end{verbatim}

    \item Even though the raw sum of members is:
\begin{verbatim}
4 + 8 + 1 + 8 + 12 = 33 bytes
\end{verbatim}

    \item The extra bytes come from \textbf{padding added by the compiler}
    to satisfy memory alignment rules, which makes access faster on most CPUs.
\end{itemize}

\newpage
\subsection{How struct differs from class}
\begin{itemize}
	\item struct are public by default.
	\item class is private by default.

	\item struct implementation:
	\begin{verbatim}
#include <iostream>
#include <string>

using namespace std;

struct Student {
private:
    string name;
    int age;
    double gpa;

public:
    // default constructor
    Student() {}

    // Constructor to initialize data members
    Student(string n, int a, double g) {
        name = n;
        age = a;
        gpa = g;
    }

    void display() {
        cout << "Name: " << name << endl;
        cout << "Age: " << age << endl;
        cout << "GPA: " << gpa << endl;
    }
};

int main() {
    // Creating objects using constructor
    Student s1("Alice", 20, 3.8);
    Student s2("Bob", 22, 3.5);

    s1.display();
    cout << "-----------------" << endl;
    s2.display();

    return 0;
}
	\end{verbatim}

	\newpage
\item class implementation
	\begin{verbatim}
#include <iostream>
#include <string>

using namespace std;

class Student {
private:
    string name;
    int age;
    double gpa;

public:
	Student() {} // default constructor

    // Constructor to initialize data members
    Student(string n, int a, double g) {
        name = n;
        age = a;
        gpa = g;
    }

    void display() {
        cout << "Name: " << name << endl;
        cout << "Age: " << age << endl;
        cout << "GPA: " << gpa << endl;
    }
};

int main() {
    // Creating objects using constructor
    Student s1("Alice", 20, 3.8);
    Student s2("Bob", 22, 3.5);

    s1.display();
    cout << "-----------------" << endl;
    s2.display();

    return 0;
}
	\end{verbatim}

	\pbreak

\item Lookover \subsubsection*{Class vs Struct} section above for more details.
\end{itemize}

\newpage

\subsection{Declaring a 'friend' of a class}
\begin{itemize}
	\item A class \textbf{does not permit} access to its \textbf{data members and methods} that are declared \textbf{private}.
	\item The rule is waived for \textbf{classes} and \textbf{functions} that are disclosed as \vocab{friend classes} or \vocab{friend functions} using the keyword \vocab{friend}
	\item \fbox{\textbf{Listing 9.14} friend}
\end{itemize}

\begin{verbatim}
class B;

class A
{
private:
    int secret = 99;

    friend class B;
};

class B
{
public:
    void show(A obj)
    {
        std::cout << obj.secret << std::endl; // allowed
    }
};
\end{verbatim}

\newpage

\subsection{union: A Special Data Storage Mechanism}
\begin{itemize}
    \item A \vocab{union} is a special class type where all non-static data members
    share the same memory location.

    \item This means only one member can hold a meaningful value at a time.

    \item A union \textbf{cannot participate in inheritance hierarchies}. It cannot
    be used as a base class or derived class.

    \item The \verb|sizeof| a union is always equal to the size of its \textbf{largest member},
    even if that member is not currently active in a given instance.

    \item Example:

\begin{verbatim}
#include <iostream>
using namespace std;

union Data
{
    int i;
    float f;
    char c;
};

int main()
{
    Data d;

    d.i = 42;
    cout << "i = " << d.i << endl;

    d.f = 3.14f;
    cout << "f = " << d.f << endl;

    d.c = 'A';
    cout << "c = " << d.c << endl;

    cout << "sizeof(Data) = " << sizeof(Data) << endl;

    return 0;
}
\end{verbatim}

    \item In this example, all members share the same memory, so writing to one
		member \textbf{\underline{overwrites}} the \textbf{previous value}.

    \item Even though multiple members exist, only one is valid at any time.

    \item The output of \verb|sizeof(Data)| depends on the largest member:
    \begin{itemize}
        \item \verb|int| might be 4 bytes
        \item \verb|float| might be 4 bytes
        \item \verb|char| might be 1 byte
    \end{itemize}

    \item Therefore:
\begin{verbatim}
sizeof(Data) = size of largest member (usually 4 bytes in this example)
\end{verbatim}

    \item The fact that a larger member exists does not change the size unless
    alignment/padding rules require it.

\item C++17 introduces \verb|std::variant| as an alternative typesafe structure to \verb|union|
\end{itemize}

\newpage

\subsection{Using Aggregate Initialization on Classes and Structs}
\begin{itemize}
    \item \vocab{Aggregate initialization} is a way of initializing objects
    (typically \verb|struct|s and simple \verb|class|es) using curly braces
    \verb|{}| without explicitly calling a constructor.

    \item It allows you to initialize all public data members in order.

    \item This works best with \verb|struct|s or classes that behave like simple data containers
    (no private members or user-defined constructors in older rules).

    \item Example using a \verb|struct|:

\begin{verbatim}
#include <iostream>
using namespace std;

struct Point
{
    int x;
    int y;
};

int main()
{
    Point p1 = {10, 20};   // aggregate initialization

    cout << "x = " << p1.x << endl;
    cout << "y = " << p1.y << endl;

    return 0;
}
\end{verbatim}

    \item The values inside the braces are assigned to data members \textbf{in order}:
    \begin{itemize}
        \item first value → first member (\verb|x|)
        \item second value → second member (\verb|y|)
    \end{itemize}

    \item Aggregate initialization also works with simple classes that have
    \textbf{public data members and no user-declared constructors}:

\begin{verbatim}
#include <iostream>
using namespace std;

class Color
{
public:
    int r;
    int g;
    int b;
};

int main()
{
    Color c1 = {255, 128, 64};

    cout << c1.r << " " << c1.g << " " << c1.b << endl;

    return 0;
}
\end{verbatim}

    \item However, if a class has a \textbf{user-defined constructor}, then
    aggregate initialization is no longer used in the same way (it will prefer constructors instead).

    \item Example difference:

\begin{verbatim}
struct A
{
    int x;
    int y;
};

A a1 = {1, 2};   // OK: aggregate initialization
\end{verbatim}

\begin{verbatim}
class B
{
public:
    B(int a, int b) {}
};

B b1 = {1, 2};   // calls constructor, NOT aggregate initialization
\end{verbatim}

    \item Aggregate initialization is useful because it provides a concise way
    to initialize simple data structures without writing constructors.

    \item It is commonly used in:
    \begin{itemize}
        \item simple structs
        \item POD (plain old data) types
        \item configuration-like data containers
    \end{itemize}
\end{itemize}


\newpage
\subsection{constexpr with Classes and Objects}
\begin{itemize}
	\item We learned \verb|constexpr| in Lesson 3 where it offers a powerful way to \textbf{improve performance} by computing the expression at \textbf{compile-time}.

	\item By marking your \textbf{functions} that operate on \textbf{constants} or \textbf{const-expression} as \verb|constexpr|, we are instructing the compile to evalute those functions and \textbf{insert their result} instead of inserting instructions that compute the result when the applicatin is executed.

	\item This keyword can also be used with \textbf{classes} and \textbf{objects} that evaluate as \textbf{constants} as demonstrated by \fbox{\textbf{Listing 9.18} constexpr with class Human}.

	\item Note the \textbf{compiler} would \textbf{ignore} the \textbf{constexpr} when the function or class is used with entities that are \textbf{not constant}.
\end{itemize}

\newpage
\subsection{QA}
\begin{itemize}
	\item \textbf{Q: What is the difference between the instance of a class and an object of that class?}

		A: Essentially none. When you \textbf{instantiate} a class, you get an \textbf{instance} that can also be called an \textbf{object}.\\

	\item \textbf{Q: What is a better way to access members: using the dot operator (.) or using the pointer operator ($->$)?}

		A: If you have a \textbf{pointer} to an \textbf{object}, the \textbf{pointer operator} would be best suited.
		If you have instantiated an \textbf{object} as a \textbf{local variable} on the stack, then the \textbf{dot operator} is best suited.\\ 

	\item \textbf{Q: Why is the parameter of a copy constructor one that takes the copy source by reference?}
	
		A: For one, the \textbf{copy constructor} is \textbf{expected} by the compiler that way. The reason behind it is that a \textbf{copy constructor} would invoke itself it it accepted the copy source \textbf{by value}, resulting in an endless \textbf{copy loop}.
\end{itemize}

\newpage
\subsection{Quiz}
\begin{enumerate}
\item My class has a raw pointer \verb|int*| that contains a dynamically allocated array of integers. Does \verb|sizeof| report different sizes depending on the number of integers in the dynamic array?

A:
\begin{itemize}
    \item No. The \verb|sizeof| operator reports only the size of the pointer itself,
    not the size of the dynamically allocated memory that it points to.

    \item The dynamically allocated array resides on the \textbf{heap}, while the
    pointer is simply a variable stored as part of the class object.

    \item For example:

\begin{verbatim}
class Numbers
{
private:
    int* data;
};
\end{verbatim}

    \item On a typical 64-bit system:

\begin{verbatim}
sizeof(Numbers) == 8
\end{verbatim}

    because the class contains only a single pointer, which is typically
    8 bytes in size.

    \item Even if the pointer refers to arrays of different lengths:

\begin{verbatim}
Numbers a;   // data points to 5 integers
Numbers b;   // data points to 100 integers
Numbers c;   // data points to 1,000,000 integers
\end{verbatim}

    the result of \verb|sizeof| is the same for every object:

\begin{verbatim}
sizeof(a) == sizeof(b) == sizeof(c)
\end{verbatim}

    \item This is because \verb|sizeof| measures only the memory occupied by the
    class object itself, not any memory that has been dynamically allocated on
    the heap.
\end{itemize}
\end{enumerate}

\newbp

\section{Implementing Inheritance}
\begin{itemize}
	\item \vocab{Object-oriented programming} is based on 4 important aspects:
	\begin{enumerate}
		\item encapsulation
		\item abstraction
		\item inheritance
		\item polymorphism
	\end{enumerate}
\end{itemize}

\pbreak

\subsection{Basics of Inheritance}
\begin{itemize}
    \item \vocab{Inheritance} is an object-oriented programming (OOP) feature that
    allows one class to \textbf{acquire} the data members and member functions of
    another class.

\item The class being \textbf{inherited from} is called the \vocab{base class}
    (also called the \textbf{parent} or \textbf{superclass}).

\item The class that \textbf{inherits from} another class is called the
    \vocab{derived class} (also called the \textbf{child} or \textbf{subclass}).

    \item Inheritance promotes \textbf{code reuse} by allowing a derived class
    to reuse functionality already implemented in its base class.

    \item A derived class may:
    \begin{itemize}
        \item Inherit existing data members.
        \item Inherit existing member functions.
        \item Add new data members.
        \item Add new member functions.
        \item Override inherited virtual functions.
    \end{itemize}

    \item The general syntax for inheritance is:

\begin{verbatim}
class DerivedClass : access-specifier BaseClass
{
    // additional members
};
\end{verbatim}

    \item The access specifier is commonly one of:
    \begin{itemize}
        \item \verb|public|
        \item \verb|protected|
        \item \verb|private|
    \end{itemize}

    \item Public inheritance is the most common form and models an
    \textbf{"is-a"} relationship.

    \item Example:

\begin{verbatim}
#include <iostream>
#include <string>
using namespace std;

class Animal
{
public:
    void Eat()
    {
        cout << "Animal is eating." << endl;
    }
};

class Dog : public Animal
{
public:
    void Bark()
    {
        cout << "Dog is barking." << endl;
    }
};

int main()
{
    Dog pet;

    pet.Eat();   // inherited from Animal
    pet.Bark();  // defined in Dog

    return 0;
}
\end{verbatim}

    \item Output:

\begin{verbatim}
Animal is eating.
Dog is barking.
\end{verbatim}

    \item In the example above:
    \begin{itemize}
        \item \verb|Animal| is the base class.
        \item \verb|Dog| is the derived class.
        \item The \verb|Dog| object inherits the \verb|Eat()| member function.
        \item The \verb|Dog| class also defines its own member function,
        \verb|Bark()|.
    \end{itemize}

    \item Conceptually, inheritance forms a hierarchy:

\begin{verbatim}
Animal
   |
   +---- Dog
\end{verbatim}

    \item Since every \verb|Dog| object is also an \verb|Animal| object,
    a \verb|Dog| can use all accessible members inherited from
    \verb|Animal| in addition to its own members.

    \item Inheritance is one of the fundamental mechanisms that enables
    polymorphism and code reuse in object-oriented programming.
\end{itemize}

\newpage

\subsubsection{Access specific keyword protected}
\begin{itemize}
	\item Here is a mechanism that allows \textbf{derived class members} to modify chosen attributes of the \textbf{base class}, while \textbf{denying access} to the same from everyone else.
	\item \verb|protected| keyword to the rescue.
\end{itemize}

\begin{tcolorbox}
	\begin{itemize}
		\item \textbf{protected}, like \textbf{public} and \textbf{private}, is also an \vocab{access specifier}.

	\item When you delcare a \textbf{class attribute} or \textbf{function} as \textbf{protected}, you are effectively making it accessible to classes that \textbf{derive} (and \textbf{friends}), yet simulateneously making it \textbf{inaccessible} to everyone else \textbf{outside the class}, including \textbf{main()}
	\end{itemize}
\end{tcolorbox}

\pbreak

\begin{itemize}
	\item \verb|protected| is the \textbf{access specifier} you shoulder use if you want a certain attribute in a \textbf{base class} to be \textbf{accessible} to classes that \textbf{derive} from this \textbf{base} class, as demonstrated in \fbox{\textbf{Listing 10.2}}.
\end{itemize}

\newpage 

\subsubsection{Base Class Initialiation --- Passing Parameters to the Base Class}
\begin{itemize}
	\item What if a \textbf{base class} were to contain an \textbf{overloaded constructor} that requires arguments at the time of \textbf{instantiation}?

	\item How would such a \textbf{base class} be \textbf{instantiated} when the \textbf{derived class} is being \textbf{constructed}?

	\item The clue lies in using \textbf{initialization lists} and in invoking the appropriate \textbf{base class constructor} via the \textbf{constructor} of the \textbf{derived class} as shown in the following code:

		\begin{verbatim}
class Base
{
public:
	Base(int someNumber)	// overloaded constructor
	{
		// use someNumber
	}
};

class Derived: public Base 
{
public:
	Derived() : Base(25)	// instantiate Base with argument 25
	{
		// derived class constructor code
	}
};
		\end{verbatim}

	\item This is nice because the base class can esnure \textbf{every derived class} is forced to engage with the specific constructor as shown in \fbox{\textbf{Listing 10.3}}

	\item In this \verb|Fish| now has a \textbf{constructor} that takes a \textbf{default parameter} initialzing \verb|Fish::isFreshWaterFish|.

		Thus, the \textbf{only possibility} to create an \textbf{object} of \verb|Fish| is via providing it a \textbf{parameter} that initialized the \textbf{protected member}.

	\item This way \verb|class Fish| ensures that the \textbf{protected member} doesn't contain a \textbf{random value}, especially if a \textbf{derive class} forgets to set it.

\item The \textbf{derived classes} are now \textbf{forced} to define a \textbf{constructor} that instantiates the \textbf{base class instance} of \verb|Fish| with the right parameter.
\end{itemize}

\newpage

\subsubsection{Derived Class Overriding Base Class's Methods}
\begin{itemize}
	\item if a \verb|class Derived| implements the same functions with the same return values and signatures as in a \verb|class Base| it inherits from, it effectively \textbf{overrides} that method in \verb|class Base| as shown:

		\begin{verbatim}
class Base
{
public:
	void DoSomething() {}
};

class Derived : public Base
{
public:
	void DoSomething() {}
};

		\end{verbatim}

	\item Thus, if \textbf{method} \verb|DoSomething()| were to be invoked using an \textbf{instance} of \verb|Derived|, the it would \textbf{not} invoke the functionality in \verb|class Base|.

	\item \fbox{\textbf{Listing 10.4} Derived classes Overriding Base method}
\end{itemize}

\newpage 

\subsubsection{Invoking Overriden Methods of a Base Class}
\begin{itemize}
	\item In \textbf{Listing 10.4} we saw a \textbf{derived class} \textbf{overriding} a function from the \textbf{base class} by implementing its versio of the same signature.

	\item If you instead want to invoke the base's method, you have to use \vocab{scope resolution operator (::)}

		\begin{verbatim}
myDinner.Fish::Swim();		// invokes Fish::Sim() using instance of tuna
							// Fish being the base class and Tuna being the child class
		\end{verbatim}

	\item \fbox{\textbf{Listing 10.5} Invoking Base class functions from derived class using ::}
\end{itemize}

\newpage

\subsubsection{Derived Class Hiding Base Class's Methods}
\begin{itemize}
\item \textbf{Overriding} can take an extreme form where \verb|Tuna::Swim()| can potentially \textbf{hide} all \textbf{overloaded} verions of \verb|Fish::Swim()| available, even causing compilation failure when the \textbf{overloaded} ones are used (hence, called \vocab{hidden}), as demonstrated by \textbf{\fbox{Listing 10.6}}.

\item We hide the function in the base class from the compiler.

\item So, if you want to invoke a \textbf{base function} like \verb|Fish::Swim(bool)| via an \textbf{instance of Tuna (child class)}, you have the following solutions

	\begin{enumerate}
		\item Solution 1: Use the \textbf{scope resolution operator} in \verb|main()|
		\begin{verbatim}
myDinner.Fish::Swim();
		\end{verbatim}

	\item Solution 2: Use keying \textbf{using} in \verb|class Tuna| to \textbf{unhide} \verb|Swim()| in \verb|class Fish|:
		\begin{verbatim}
class Tuna : public Fish
{
public:
	using Fish::Swim;		// unhide all Swim() methods in class Fish

	void Swim()
	{
		cout << "Tuna swims real fast" << endl;
	}
};
		\end{verbatim}

	\item Solution 3: \textbf{Override} all \textbf{overloaded variants} of \verb|Swim()| in \verb|class Tuna| (invoke methods of \verb|Fish::Swim(...)| via \verb|Tuna::Fish(...)| if you want):

		\begin{verbatim}
class Tuna : public Fish
{
public:
	void Swim(bool isFreshWaterFish)
	{
		Fish::Swim(isFreshWaterFish);
	}

	void Swim()
	{
		cout << "Tuna swims real fast" << endl;
	}
};
		\end{verbatim}
	\end{enumerate}
\end{itemize}


\newpage

\subsection{Order of Construction}
\begin{itemize}
    \item When an object of a derived class is created, C++ constructs the
    \textbf{base class portion first}, followed by the \textbf{derived class portion}.

    \item This ensures that the base class is fully initialized before the
    derived class begins constructing its own data members.

    \item The order of construction is:
    \begin{enumerate}
        \item Base class constructor.
        \item Derived class constructor.
    \end{enumerate}

    \item Example:

\begin{verbatim}
#include <iostream>
using namespace std;

class Animal
{
public:
    Animal()
    {
        cout << "Animal constructor" << endl;
    }
};

class Dog : public Animal
{
public:
    Dog()
    {
        cout << "Dog constructor" << endl;
    }
};

int main()
{
    Dog pet;

    return 0;
}
\end{verbatim}

    \item Output:

\begin{verbatim}
Animal constructor
Dog constructor
\end{verbatim}

    \item Although the object being created is a \verb|Dog|, the
    \verb|Animal| portion must be constructed first because every
    \verb|Dog| object contains an \verb|Animal| subobject.

    \item A useful way to visualize the process is:

\begin{verbatim}
Create Dog object
        |
        v
Construct Animal
        |
        v
Construct Dog
\end{verbatim}
\end{itemize}

\newpage

\subsection{Order of Destruction}
\begin{itemize}
    \item Object destruction occurs in the \textbf{reverse order} of construction.

    \item When a derived object is destroyed, the derived class destructor
    executes first, followed by the base class destructor.

    \item The order of destruction is:
    \begin{enumerate}
        \item Derived class destructor.
        \item Base class destructor.
    \end{enumerate}

    \item Example:

\begin{verbatim}
#include <iostream>
using namespace std;

class Animal
{
public:
    Animal()
    {
        cout << "Animal constructor" << endl;
    }

    ~Animal()
    {
        cout << "Animal destructor" << endl;
    }
};

class Dog : public Animal
{
public:
    Dog()
    {
        cout << "Dog constructor" << endl;
    }

    ~Dog()
    {
        cout << "Dog destructor" << endl;
    }
};

int main()
{
    Dog pet;

    return 0;
}
\end{verbatim}

    \item Output:

\begin{verbatim}
Animal constructor
Dog constructor
Dog destructor
Animal destructor
\end{verbatim}

    \item The derived class is destroyed first because it may depend on
    resources owned by the base class. The base class is destroyed only
    after the derived class has finished cleaning up.

    \item A useful way to visualize destruction is:

\begin{verbatim}
Destroy Dog object
        |
        v
Destroy Dog
        |
        v
Destroy Animal
\end{verbatim}

    \item In summary:
    \begin{itemize}
        \item Construction proceeds from the base class to the derived class.
        \item Destruction proceeds from the derived class back to the base class.
    \end{itemize}
\end{itemize}

\newpage

\subsection{Private Inheritance}
\begin{itemize}
    \item \vocab{Private inheritance} differs from \vocab{public inheritance}
    in that the inherited members become \textbf{private} members of the
    derived class (unless they were already private in the base class).

    \item The general syntax is:

\begin{verbatim}
class Derived : private Base
{
    // ...
};
\end{verbatim}

    \item With private inheritance:
    \begin{itemize}
        \item Public members of the base class become private in the derived class.
        \item Protected members of the base class also become private in the derived class.
        \item Private members of the base class remain inaccessible to the derived class.
    \end{itemize}

    \item Example:

\begin{verbatim}
#include <iostream>
using namespace std;

class Animal
{
public:
    void Eat()
    {
        cout << "Animal is eating." << endl;
    }
};

class Dog : private Animal
{
public:
    void MakeDogEat()
    {
        Eat();     // OK: inherited as private
    }
};

int main()
{
    Dog pet;

    pet.MakeDogEat();

    // Compile error:
    // pet.Eat();

    return 0;
}
\end{verbatim}

    \item Output:

\begin{verbatim}
Animal is eating.
\end{verbatim}

    \item Although \verb|Eat()| was originally a public member of
    \verb|Animal|, it becomes a private member of \verb|Dog| because
    of private inheritance.

    \item Therefore:
    \begin{itemize}
        \item The \verb|Dog| class itself can call \verb|Eat()|.
        \item Code outside of \verb|Dog| cannot call \verb|Eat()| through a
        \verb|Dog| object.
    \end{itemize}

    \item Unlike public inheritance, private inheritance does \textbf{not}
    model an \textbf{"is-a"} relationship.

    \item Instead, private inheritance is sometimes used to express an
    implementation relationship, where the derived class reuses the base
    class's implementation without exposing its interface publicly.

    \item In modern C++, composition is generally preferred over private
    inheritance unless there is a specific reason to inherit privately.
\end{itemize}

\newpage
\subsection{Composition (side note)}
\begin{itemize}
    \item \vocab{Composition} is an object-oriented programming technique in which
    one class contains an object of another class as one of its data members.

    \item Composition models a \textbf{"has-a"} relationship.

	\item Inheritance models a \textbf{"is-a"} relationship.

    \item For example, a \verb|Car| \textbf{has an} \verb|Engine|, so composition
    is an appropriate design choice.

\begin{verbatim}
class Engine
{
public:
    void Start()
    {
        cout << "Engine started." << endl;
    }
};

class Car
{
private:
    Engine engine;    // Car has an Engine

public:
    void StartCar()
    {
        engine.Start();
    }
};
\end{verbatim}

    \item By contrast, inheritance models an \textbf{"is-a"} relationship.

    \item For example, a \verb|Dog| \textbf{is an} \verb|Animal|, so inheritance
    is appropriate.

\begin{verbatim}
class Animal
{
public:
    void Eat() {}
};

class Dog : public Animal
{
};
\end{verbatim}

    \item As a rule of thumb:
    \begin{itemize}
        \item Use \textbf{inheritance} when there is an \textbf{"is-a"} relationship.
        \item Use \textbf{composition} when there is a \textbf{"has-a"} relationship.
    \end{itemize}

    \item This is why modern C++ encourages the design principle:

    \begin{quote}
        \textbf{``Favor composition over inheritance.''}
    \end{quote}

    \item In many situations, composition leads to code that is easier to
    understand, more flexible, and simpler to maintain.
\end{itemize}

\newpage

\subsection{Protected Inheritance}
\begin{itemize}
    \item \vocab{Protected inheritance} is another form of inheritance in which
    the inherited \verb|public| and \verb|protected| members of the base class
    become \verb|protected| members of the derived class.

    \item The general syntax is:

\begin{verbatim}
class Derived : protected Base
{
    // ...
};
\end{verbatim}

    \item With protected inheritance:
    \begin{itemize}
        \item Public members of the base class become protected.
        \item Protected members of the base class remain protected.
        \item Private members of the base class are still inaccessible.
    \end{itemize}

    \item Example:

\begin{verbatim}
#include <iostream>
using namespace std;

class Animal
{
public:
    void Eat()
    {
        cout << "Animal is eating." << endl;
    }
};

class Dog : protected Animal
{
public:
    void MakeDogEat()
    {
        Eat();    // OK
    }
};

class Puppy : public Dog
{
public:
    void Feed()
    {
        Eat();    // OK because Eat() is protected in Dog
    }
};

int main()
{
    Dog pet;
    pet.MakeDogEat();

    // Compile error:
    // pet.Eat();

    Puppy puppy;
    puppy.Feed();

    return 0;
}
\end{verbatim}

    \item Output:

\begin{verbatim}
Animal is eating.
Animal is eating.
\end{verbatim}

    \item In this example:
    \begin{itemize}
        \item The \verb|Dog| class can call \verb|Eat()|.
        \item The \verb|Puppy| class can also call \verb|Eat()| because it inherits
        from \verb|Dog|.
        \item Code outside these classes cannot call \verb|Eat()| through a
        \verb|Dog| or \verb|Puppy| object.
    \end{itemize}

    \item Therefore, protected inheritance allows derived classes to continue
    using inherited functionality while preventing direct access from outside
    the class hierarchy.

    \item The access transformations are summarized below:

\begin{verbatim}
Base member      After protected inheritance
--------------------------------------------
public           protected
protected        protected
private          inaccessible
\end{verbatim}

    \item Protected inheritance is used less frequently than public inheritance.
    It is most useful when a class wishes to reuse a base class's implementation
    while allowing only further derived classes to access the inherited members.

	\pbreak

\item \boxednote{
	Private and protected signifiy a "has-a" relationship}
\end{itemize}

\newpage

\subsection{Multiple Inheritance}
\begin{itemize}
    \item \vocab{Multiple inheritance} allows a class to inherit from
    \textbf{more than one base class}.

    \item This enables a derived class to reuse the functionality of multiple
    independent classes.

    \item The general syntax is:

\begin{verbatim}
class Derived : public Base1, public Base2
{
    // ...
};
\end{verbatim}

    \item Example:

\begin{verbatim}
#include <iostream>
using namespace std;

class Printer
{
public:
    void Print()
    {
        cout << "Printing..." << endl;
    }
};

class Scanner
{
public:
    void Scan()
    {
        cout << "Scanning..." << endl;
    }
};

class AllInOne : public Printer, public Scanner
{
};

int main()
{
    AllInOne device;

    device.Print();
    device.Scan();

    return 0;
}
\end{verbatim}

    \item Output:

\begin{verbatim}
Printing...
Scanning...
\end{verbatim}

    \item In this example:
    \begin{itemize}
        \item \verb|Printer| is a base class.
        \item \verb|Scanner| is another base class.
        \item \verb|AllInOne| inherits from both classes.
    \end{itemize}

    \item As a result, an \verb|AllInOne| object has access to the member
    functions of both base classes.

    \item One issue that can arise with multiple inheritance is
    \vocab{name ambiguity}. For example:

\begin{verbatim}
class A
{
public:
    void Hello()
    {
        cout << "Hello from A" << endl;
    }
};

class B
{
public:
    void Hello()
    {
        cout << "Hello from B" << endl;
    }
};

class C : public A, public B
{
};

int main()
{
    C obj;

    // obj.Hello();   // Compile error: ambiguous

    obj.A::Hello();
    obj.B::Hello();

    return 0;
}
\end{verbatim}

    \item Since both base classes define a member function named
    \verb|Hello()|, the compiler cannot determine which one should be called.

    \item The ambiguity can be resolved using the
    \vocab{scope resolution operator} (\verb|::|), as shown above.

    \item Multiple inheritance is a powerful language feature, but it should be
    used carefully because it can make class hierarchies more complex.

    \item One well-known issue associated with multiple inheritance is the
    \textbf{Diamond Problem}, which is discussed separately using
    \vocab{virtual inheritance}.
\end{itemize}

\newpage
\subsection{Avoid Inheritance via \texttt{final}}
\begin{itemize}
    \item The \verb|final| keyword prevents a class from being used as a
    \textbf{base class}.

    \item Any attempt to derive from a class marked \verb|final| results in a
    compile-time error.

    \item The general syntax is:

\begin{verbatim}
class ClassName final
{
    // ...
};
\end{verbatim}

    \item Example:

\begin{verbatim}
#include <iostream>
using namespace std;

class Animal final
{
public:
    void Speak()
    {
        cout << "Animal speaks." << endl;
    }
};

// Compile error:
// class Dog : public Animal
// {
// };

int main()
{
    Animal pet;
    pet.Speak();

    return 0;
}
\end{verbatim}

    \item Output:

\begin{verbatim}
Animal speaks.
\end{verbatim}

    \item In the example above, \verb|Animal| is declared as
    \verb|final|, so no other class may inherit from it.

    \item Attempting to write:

\begin{verbatim}
class Dog : public Animal
{
};
\end{verbatim}

    results in a compile-time error because \verb|Animal| is a
    \verb|final| class.

    \item The \verb|final| keyword is commonly used to:
    \begin{itemize}
        \item Prevent further inheritance.
        \item Protect a class from unintended extension.
        \item Indicate that the class represents a complete implementation.
    \end{itemize}

    \item The \verb|final| keyword can also be applied to
    \textbf{virtual member functions}.

    \item When a virtual function is declared \verb|final|, derived classes
    may inherit the function, but they cannot override it.

\begin{verbatim}
class Base
{
public:
    virtual void Print() final
    {
        cout << "Base::Print()" << endl;
    }
};

class Derived : public Base
{
    // Compile error:
    // void Print() override { }
};
\end{verbatim}

    \item In summary:
    \begin{itemize}
        \item Applying \verb|final| to a class prevents inheritance.
        \item Applying \verb|final| to a virtual function prevents overriding.
    \end{itemize}
\end{itemize}

\newbp

\section{Polymorphism}
\subsection{Basics of Polymorphism}
\begin{itemize}
	\item \textbf{"Poly"} is Greek for \textit{\textbf{many}} and \textbf{"morph"} means \textbf{\textit{form}}.
\item \vocab{Polymorphism} is the feature of \textbf{object-oriented} programming that allows \textbf{objects} of different types to be accessed through a common interface, with each object providing its own implementation of the same operation.

\item We go over \textbf{polymorphic} behavior that can be implemented in C++ via the \textbf{inheirtance hierarchy} known as \vocab{subtype polymorphism}.
\end{itemize}

\pbreak

\subsubsection{Need for Polymorphic Behavior}
\begin{itemize}
	\item Last section we sa how \verb|Tuna| and \verb|Carp| \textbf{inheirt public methods} \verb|Swim()| from \verb|Fish| (Listing 10.1)

	\item It is, however, possible that both \verb|Tuna| and \verb|Carp| provide their \textbf{own} \verb|Tuna::Swim()| and \verb|Carp::Swim()| methods to make \verb|Tuna| and \verb|Carp| different swimmers.

	\item If a user with an \textbf{instance} of \verb|Tuna| uses the \underline{\textbf{base class type}} to invoke \verb|Fish::Swim()|, he ends up executing only the \textbf{generic part} \verb|Fish::Swim()| and \underline{\textbf{not}} \verb|Tuna::Swim()|, even though that \textbf{base class instance} \verb|Fish| is a part of a \verb|Tuna|.

	\item \fbox{\textbf{Listing 11.1} Invoking Methods Using an Instance of the Base Class Fish That Belongs to a Tuna}

	\item Code shows if we cast it back to a \verb|Fish|, \verb|Tuna| behaves like a generic Fish and not a Tuna.

	\item What the user would \textbf{ideally expect} is that an \textbf{object} of \textbf{type} \verb|Tuna| behaves like a tuna \textbf{even if} the methods invoked is \verb|Fish::Swim()|.

	\item We use \vocab{virtual functions} to fix this.
\end{itemize}

\newpage

\subsubsection{Polymorphic Behavior Implemented Using Virtual Functions}
\begin{itemize}
	\item You have access to an \textbf{object} of \textbf{type} \verb|Fish|, via \textbf{pointer} \verb|Fish*| or \textbf{reference} \verb|Fish&|.

	This \textbf{object} could have been \textbf{instantiated} solely as \verb|Fish| or be part of a \verb|Tuna| or \verb|Carp| that \textbf{inheirts} from \verb|Fish|.

\item You don't know (or don't give a fuck)

\item You invoke \textbf{method} \verb|Swim()| using the \textbf{pointer} or \textbf{reference} like this:
	\begin{verbatim}
pFish->Swim();
myFish.Swim();
	\end{verbatim}

	\item What you expect is that the \textbf{object} \verb|Fish| swims
		\begin{itemize}
			\item as a \verb|Tuna| if it is part of a \verb|Tuna|, 
			\item as a \verb|Carp| if it is part of a \verb|Carp|, 
			\item or an anonymous \verb|Fish| if it wasn't instantiatead as part of a specialized class such as \verb|Tuna| or \verb|Carp|.
		\end{itemize}

	\item You can ensure this by \textbf{declaring function} \verb|Swim()| in the \underline{\textbf{base class}} as a \vocab{virtual function}

	\begin{verbatim}
class Base
{
	virtual ReturnType FunctionName (Parameter List);
};

class Derived
{
	ReturnType FunctionName (Parameter List);
};
	\end{verbatim}

\item Use of the keyboard \verb|virtual| means that the \textbf{compiler} ensures that any \textbf{overriding variant} of the requested \textbf{base class method} is invoked.

\item Thus, if \verb|Swim()| is declared \textbf{virtual}, invoking \verb|myFish.Swim()| (myFish being of \textbf{type} \verb|Fish&|) results in \verb|Tuna::Swim()| being executed (Listing 11.2)

\item This is \vocab{polymorphism}: treating different fishes as a common type \verb|Fish|, yet \textbf{ensuring} that the right implementation of \verb|Swim()| supplied by the \textbf{derived types} is executed
	\end{itemize}

	\newpage

\subsubsection{Need For Virtual Destructors}
\begin{itemize}
	\item What happens when a function calls operator \verb|delete| using a \textbf{pointer} of type \verb|Base*| that actually points to the \textbf{instance of type Derived}?

	\item See \fbox{\textbf{Listing 11.3} A function that invokes operator 'delete' on Base*}

	\item Note in the listing: while \verb|Tuna| and \verb|Fish| were both \textbf{constructed} on the \textbf{free store (heap)} due to \verb|new|, the \textbf{destructor} of \verb|Tuna| \underline{\textbf{was not invoked}} during \verb|delete|, a rather only that of the \verb|Fish|.

	\item This is in \textbf{stark contrast} to the \textbf{construction} and \textbf{destruction} of \textbf{local member} \verb|myDinner| where \textbf{all} the \textbf{constructors} and \textbf{destructors} are invoked.
		
	\item Clearly, for the \textbf{free store} we have amiss. \verb|~Tuna()| needs to be invoked.

		\begin{verbatim}
Allocating a Tuna on the free store:
Constructed Fish
Constructed Tuna
Deleteing the Tuna:
Destructed Fish
Instantiating a Tuna on the stack:
Constructed Fish
Constructed Tuna
Automatic destruction as it goes out of scope:
Destructed Tuna
Destructed Fish
		\end{verbatim}

		\pbreak

	\item The flaw means that the \textbf{destructor} of a \textbf{deriving class} that has been instantiated on the \textbf{free store} using \verb|new| wwould \textbf{not} be invoked if \verb|delete| is called using a \textbf{pointer of type} \verb|Base*|.

	\item This can result in resources not being released, memory leaks, and so on and is a problem that is not to be taken lightly.

	\item \boxednote{\textbf{Listing 11.4} Using virtual destructors to ensure that destructors in derived classes are invoked when deleting a pointer of Type Base*}

		\begin{verbatim}
Allocating a Tuna on the free store:
Constructed Fish
Constructed Tuna
Deleteing the Tuna:
Destructed Tuna
Destructed Fish
Instantiating a Tuna on the stack:
Constructed Fish
Constructed Tuna
Automatic destruction as it goes out of scope:
Destructed Tuna
Destructed Fish
		\end{verbatim}
\end{itemize}

\pbreak

\begin{tcolorbox}
	\begin{itemize}
		\item \textbf{Always declare the \underline{base class destructor} as 'virtual'}

		\begin{verbatim}
class Base
{
public:
	virtual ~Base() {};		// virtual destructor
};
		\end{verbatim}

	\item This ensures that one with a \textbf{pointer} \verb|Base*| \underline{cannot invoke} \verb|delete| in a way that instances of the \textbf{deriving classes} are not correctly destroyed.
	\end{itemize}
\end{tcolorbox}

\newpage

\subsubsection{How do Virtual Functions Work? --- The Virtual Function Table (vtable)}

\begin{itemize}
	\item A \vocab{virtual function} allows \textbf{runtime polymorphism}. When a member function is declared with the \texttt{virtual} keyword, the function that is executed is determined by the \textbf{actual type of the object} rather than the type of the pointer or reference.

	\item To support this behavior, most C++ compilers implement a hidden data structure called the \vocab{Virtual Function Table (vtable)}. Although the C++ standard does \textbf{not} require this implementation, it is the method used by virtually all modern compilers.

	\item The \textbf{vtable} is a table of function pointers. Each class that contains one or more virtual functions has its own vtable.

	\item Every object of a class containing virtual functions secretly stores a hidden pointer called the \vocab{vptr} (virtual pointer). This pointer points to the class's vtable.

	\item During object construction, the compiler automatically initializes the object's \texttt{vptr} so that it points to the correct vtable for that object's type.

	\item When a virtual function is called through a pointer or reference:
	\begin{enumerate}
		\item The program follows the object's hidden \texttt{vptr}.
		\item The \texttt{vptr} points to the appropriate vtable.
		\item The corresponding function pointer is retrieved from the table.
		\item The function pointed to is invoked.
	\end{enumerate}

	\item This process is known as \vocab{dynamic dispatch} or \vocab{late binding}, since the decision of which function to execute is made at \textbf{runtime} instead of compile time.

	\item Without virtual functions, C++ performs \vocab{static binding} (also called \vocab{early binding}), meaning the compiler decides which function to call solely from the pointer or reference type.

	\item Consider the following example:

\begin{verbatim}
class Fish
{
public:
    virtual void Swim()
    {
        cout << "Fish swims\n";
    }
};

class Tuna : public Fish
{
public:
    void Swim() override
    {
        cout << "Tuna swims fast\n";
    }
};

Fish* pFish = new Tuna;
pFish->Swim();
\end{verbatim}

	\item Even though \texttt{pFish} is a \texttt{Fish*}, the object it points to is actually a \texttt{Tuna}. Since \texttt{Swim()} is virtual, the compiler uses the object's \texttt{vptr} to locate Tuna's implementation in the vtable, producing:

\begin{verbatim}
Tuna swims fast
\end{verbatim}

	\item If \texttt{virtual} were omitted, the compiler would instead call:

\begin{verbatim}
Fish swims
\end{verbatim}

	because it would only consider the pointer's declared type (\texttt{Fish*}).

	\item Conceptually, the compiler generates structures similar to the following (greatly simplified):

\begin{verbatim}
Fish vtable
+----------------------+
| &Fish::Swim          |
| &Fish::~Fish         |
+----------------------+

Tuna vtable
+----------------------+
| &Tuna::Swim          |
| &Tuna::~Tuna         |
+----------------------+
\end{verbatim}

	\item Each object contains a hidden pointer:

\begin{verbatim}
          Fish object

      +------------------+
      | vptr ----------+ |
      +----------------|--+
      | data members   |  |
      +----------------+  |
                           |
                           v
                  Fish vtable
               +------------------+
               | Fish::Swim       |
               | Fish::~Fish      |
               +------------------+
\end{verbatim}

	\item If the object is actually a \texttt{Tuna}, its hidden \texttt{vptr} instead points to Tuna's vtable:

\begin{verbatim}
          Tuna object

      +------------------+
      | vptr ----------+ |
      +----------------|--+
      | Fish members    | |
      | Tuna members    | |
      +-----------------+ |
                           |
                           v
                  Tuna vtable
               +------------------+
               | Tuna::Swim       |
               | Tuna::~Tuna      |
               +------------------+
\end{verbatim}

	\item Since the \texttt{vptr} belongs to the object itself, the correct virtual functions are called regardless of whether the object is accessed through a base-class pointer, base-class reference, or directly.

	\item Virtual destructors also use the vtable. When deleting an object through a base-class pointer:

\begin{verbatim}
Fish* pFish = new Tuna;
delete pFish;
\end{verbatim}

	the virtual destructor ensures the call sequence is:

\begin{verbatim}
Tuna::~Tuna()
Fish::~Fish()
\end{verbatim}

	If the destructor were not virtual, only \texttt{Fish::~Fish()} would be called, resulting in \textbf{undefined behavior} and potentially leaking resources owned by the derived class.

	\item Although virtual functions introduce a very small runtime overhead (one extra pointer lookup and one hidden pointer per object), this cost is generally negligible compared to the flexibility provided by runtime polymorphism.

	\item \textbf{Summary:}
	\begin{itemize}
		\item Each polymorphic class has one vtable.
		\item Each object stores one hidden \texttt{vptr}.
		\item Virtual function calls use the \texttt{vptr} to locate the correct function in the vtable.
		\item Function selection occurs at runtime (dynamic dispatch).
		\item Virtual destructors ensure complete destruction of derived objects when deleting through base-class pointers.
	\end{itemize}
\end{itemize}

\newpage

\subsection{Abstract Base Classes and Pure Virtual Functions}
\begin{itemize}
	\item A \textbf{base class} that \textbf{\underline{cannot be instantiated}} is called an \vocab{abstract base class}.

	\item Such \textbf{base class} fulfills only \underline{\textbf{one purpose}}, that of \underline{\textbf{being derived from}}.
	\item C++ allows you to create an \vocab{abstract base class} using \vocab{pure virtual functions}

		\pbreak

	\item A \vocab{virtual method} is said to be \vocab{pure virtual} when it has a \textbf{declaration} as shown in the following:

	\begin{verbatim}
class AbstractBase
{
public:
	virtual void DoSometing() = 0;		// pure virtual method
};
	\end{verbatim}

\item This \textbf{declaration} essentially tells the \textbf{compiler} that \verb|DoSomething()| \underline{\textbf{needs to be implemented}} and by \textbf{class that derives} from \verb|AbstractBase|:
\begin{verbatim}
class Derived : public AbstractBase
{
public:
	void DoSomething()
	{
		cout << "Implemented virtual function" << endl;	
	}
};
\end{verbatim}

\item Thus, what the class \verb|AbstractBase| has done is that it has \underline{\textbf{enforced}} class \verb|Derived| to \textbf{supply an implementation} for \textbf{virtual method} \verb|DoSomething()|.

\item This functionality where a \textbf{base class} can \textbf{enforce} support of methods with a \textbf{specified name and signature} in \textbf{classes} that \underline{\textbf{derive}} from it is that of an \vocab{interface}.

\item So, making \verb|class Fish| an \vocab{abstract base class} with \verb|Swim()| as a \vocab{pure virtual function} \underline{\textbf{ensures}} that \verb|Tuna| that \textbf{derives} from \verb|Fish| implements \verb|Tuna::Swim()| and swims like a Tuna and not like just any Fish.
	
\item \fbox{\textbf{Listing 11.6} class Fish as an Abstract Base Class for Tuna and Carp}
\end{itemize}

\newpage

\subsection{Using 'virtual' Inheritance to Solve the Diamond Problem}

\begin{itemize}
	\item The \vocab{diamond problem} occurs in \textbf{multiple inheritance} when a class inherits from two classes that both inherit from the same base class.

	\item This creates a diamond-shaped inheritance structure:

\begin{verbatim}
        A
       / \
      B   C
       \ /
        D
\end{verbatim}

	\item In this structure:
	\begin{itemize}
		\item Class \texttt{A} is the base class.
		\item Classes \texttt{B} and \texttt{C} both inherit from \texttt{A}.
		\item Class \texttt{D} inherits from both \texttt{B} and \texttt{C}.
	\end{itemize}

	\item The problem arises because \texttt{D} ends up containing \textbf{two copies of \texttt{A}}:
	\begin{itemize}
		\item One inherited through \texttt{B}
		\item One inherited through \texttt{C}
	\end{itemize}

	\item This leads to ambiguity. For example, if \texttt{A} has a member variable or function, calling it from \texttt{D} becomes unclear:

\begin{verbatim}
D obj;
obj.A::function();  // ambiguous: which A?
\end{verbatim}

	\item It also wastes memory since two separate copies of \texttt{A}'s data exist inside \texttt{D}.

	\item C++ solves this problem using \vocab{virtual inheritance}.

	\item When \texttt{B} and \texttt{C} virtually inherit from \texttt{A}, they share a single common instance of \texttt{A} inside \texttt{D}.

\begin{verbatim}
        A
       / \
   (virtual) (virtual)
     B       C
       \     /
         D
\end{verbatim}

	\item This is written in code as:

\begin{verbatim}
class A
{
public:
    void foo() {}
};

class B : virtual public A {};
class C : virtual public A {};

class D : public B, public C {};
\end{verbatim}

	\item Now, \texttt{D} contains only \textbf{one shared instance of A}, eliminating duplication and ambiguity.

	\item Without virtual inheritance, the compiler cannot decide which \texttt{A} subobject to use when accessing members through \texttt{D}, resulting in a compilation error.

\item With \vocab{virtual inheritance}, the \underline{\textbf{most derived class}} (\texttt{D}) is \textbf{responsible} for constructing the \textbf{shared base class} \texttt{A}.

	\item Example:

\begin{verbatim}
class A
{
public:
    int value;
};

class B : virtual public A {};
class C : virtual public A {};

class D : public B, public C {};

int main()
{
    D obj;
    obj.value = 10;  // OK: only one A exists
}
\end{verbatim}

	\item \textbf{Key idea:}
	\begin{itemize}
		\item Virtual inheritance ensures that only one copy of a shared base class exists in the final derived object.
		\item It resolves ambiguity and prevents duplicate state.
	\end{itemize}

	\item \textbf{Tradeoff:}
	\begin{itemize}
		\item Virtual inheritance introduces a small runtime and memory overhead because the shared base class is accessed indirectly.
		\item However, it is necessary when designing complex inheritance hierarchies.
	\end{itemize}

	\pbreak

\item \boxednote{
	{\textbf{Listing 11.8} Demonstrating How 'virtual' keyword in Inheritance Heirarchy helps restrict the number of instances of base class Animal to One}}
\end{itemize}

\pbreak

\begin{tcolorbox}
\begin{itemize}
	\item In C++, \texttt{public virtual Animal} and \texttt{virtual public Animal} are \textbf{equivalent}. The order of the specifiers does not affect the meaning.

	\item Both forms indicate:
	\begin{itemize}
		\item \textbf{public inheritance} (controls access level)
		\item \textbf{virtual inheritance} (ensures a single shared base class instance in diamond inheritance)
	\end{itemize}

\item The C++ grammar allows \underline{\textbf{multiple valid orders of inheritance specifiers}}  for flexibility, so both are accepted by the compiler.

	\item There is no difference in behavior or memory layout between the two forms.

	\item \textbf{Style convention:} many modern C++ guidelines prefer \texttt{public virtual Base} because it reads as ``publicly inherits from Base'' with virtual inheritance as an additional detail.

		\begin{verbatim}
class B : public virtual A {};   // ok
class B : virtual public A {};   // ok
class B : protected virtual A {}; // ok
class B : virtual protected A {}; // ok
class B : private virtual A {};   // ok
class B : virtual private A {};   // ok
		\end{verbatim}
\end{itemize}
\end{tcolorbox}

\newpage
\subsection{Specifier 'override' to Indicate Intention to Override}

\begin{itemize}
	\item The \texttt{override} specifier in C++ is used to explicitly indicate that a member function is intended to \textbf{override a virtual function} from a base class.

	\item It improves code safety by letting the compiler verify that:
	\begin{itemize}
		\item The base class actually has a matching virtual function.
		\item The function signature in the derived class correctly matches the base class version.
	\end{itemize}

	\item If there is no matching virtual function in the base class, the compiler will produce an error.

	\item This helps prevent subtle bugs caused by accidental mismatches in function signatures.

	\item Example without \texttt{override} (bug-prone):

\begin{verbatim}
class Animal
{
public:
    virtual void speak()
    {
        cout << "Animal sound\n";
    }
};

class Dog : public Animal
{
public:
    void speak(int x)   // does NOT override Animal::speak()
    {
        cout << "Bark\n";
    }
};
\end{verbatim}

	\item In this case, \texttt{Dog::speak(int)} is a completely new function, not an override, because the signature does not match.

	\item Using \texttt{override} catches this mistake:

\begin{verbatim}
class Dog : public Animal
{
public:
    void speak(int x) override   // compiler error
    {
        cout << "Bark\n";
    }
};
\end{verbatim}

	\item The compiler will reject this because there is no matching virtual function in the base class with the same signature.

	\item Correct usage:

\begin{verbatim}
class Dog : public Animal
{
public:
    void speak() override
    {
        cout << "Bark\n";
    }
};
\end{verbatim}

	\item \textbf{Benefits of \texttt{override}:}
	\begin{itemize}
		\item Prevents accidental signature mismatches.
		\item Improves code readability by making intent explicit.
		\item Helps the compiler catch polymorphism errors early.
	\end{itemize}

	\item \textbf{Key idea:} \texttt{override} does not change runtime behavior — it is purely a \textbf{compile-time safety check}.
\end{itemize}

\newpage
\subsection{Using 'final' to Prevent Function Overriding}

\begin{itemize}
\item The \texttt{final} specifier in C++ is used to \underline{\textbf{prevent further overriding}} of a \textbf{virtual function} or \underline{\textbf{prevent inheritance}} from a \textbf{class}.

	\item When applied to a virtual function, \texttt{final} ensures that no derived class can override that function.

	\item Example:

\begin{verbatim}
class Animal
{
public:
    virtual void speak() final
    {
        cout << "Animal sound\n";
    }
};
\end{verbatim}

	\item In this case, any attempt to override \texttt{speak()} in a derived class will result in a compiler error:

\begin{verbatim}
class Dog : public Animal
{
public:
    void speak() override   // error: cannot override final function
    {
        cout << "Bark\n";
    }
};
\end{verbatim}

	\item \textbf{Key idea:} marking a function as \texttt{final} "locks" its implementation in the inheritance hierarchy.

	\item \texttt{final} can also be applied to an entire class to prevent inheritance altogether:

\begin{verbatim}
class Animal final
{
public:
    void speak()
    {
        cout << "Animal sound\n";
    }
};
\end{verbatim}

	\item In this case, any attempt to derive from the class is illegal:

\begin{verbatim}
class Dog : public Animal   // error: cannot inherit from final class
{
};
\end{verbatim}

	\item \textbf{Benefits of \texttt{final}:}
	\begin{itemize}
		\item Prevents unintended modification of critical behavior.
		\item Allows the compiler to optimize virtual dispatch in some cases.
		\item Clearly communicates design intent to other developers.
	\end{itemize}

	\item \textbf{Key idea:} \texttt{final} is a compile-time restriction that strengthens class and interface design by limiting extensibility.
\end{itemize}

\newpage
\subsection{Virtual Copy Constructors?}
\begin{itemize}
	\item The question mark here is because it is technically impossible in C++ to have \textbf{virtual copy constructors}.

	\item \textbf{Virtual copy constructors} are \textbf{not possible} because the concept of \texttt{virtual} functions in C++ relies on \textbf{runtime polymorphism}, where a function call is resolved based on the \textbf{dynamic type of an already-constructed object}.

	\item A constructor, however, is fundamentally different: it is responsible for \textbf{creating an object from scratch}. At the time a constructor runs, there is \textbf{no fully constructed object yet}, and therefore no vtable-based dispatch can occur.

	\item In other words, \texttt{virtual} functions require an existing object with a \textbf{vptr} (virtual table pointer), but during construction the object is still being built, so polymorphic dispatch cannot apply.

	\item Additionally, constructors are \textbf{not inherited in the same way as normal member functions} and are always called based on the \textbf{static type being constructed}, not the dynamic type.

	\item \textbf{Copy constructors specifically} must create a new object of a known type, so their signature and target type are fixed at compile time. There is no meaningful way to "choose" a derived constructor at runtime.

	\item For this reason, C++ does not allow constructors (including copy constructors) to be declared as \texttt{virtual}.

	\item \textbf{Key idea:} virtual dispatch requires an existing object and runtime type information, while construction is a compile-time process that produces the object in the first place. These two mechanisms are fundamentally incompatible.

		\pbreak

	\item Having said that, there is a nice workaround in the form of defining your own clone function that allows you to do just that:

		\begin{verbatim}
class Fish
{
public:
	virtual Fish* Clone() const = 0; // pure virtual function
};

class Tuna : public Fish 
{
// ... other members
public:
	Tuna * Clone() const	// virtual clone function
	{
		return new Tuna(*this);	// return new Tuna that is a copy of this
	}
};
		\end{verbatim}

	\item Thus, \textbf{virtual function} \verb|Clone| is a \vocab{smulated virtual copy constructor} that needs to be explicitly invoked, as shown in \fbox{\textbf{Listing 11.9}}
\end{itemize}

\newpage
\subsection{QA}

\newpage
\subsection{Quiz}

\newbp

\section{Operator Types and Operator Overloading}
\begin{itemize}
	\item In addition to encapsulating data and methods, classes can also encapsulate \textbf{operators} that make it easy to operate on \textbf{instances} of this class.
	\item You can use these \textbf{operators} to perform operations such as \textbf{assignment} or \textbf{addition} on \textbf{class objects} similar to those on integers that you saw in Lesson 5.

	\item Just like \textbf{functions}, \fbox{\textbf{operators}} can also be overloaded.
\end{itemize}

\pbreak

\subsection{What are Operators in C++?}
\begin{itemize}
	\item On a syntactical level, there is very little that differentiates an \textbf{operator} from a \textbf{function}, save for the use of the keyword \verb|operator|.

\item Here is an \vocab{operator declaration}:

	\begin{verbatim}
return_type operator operator_symbol (...parameter list...);
	\end{verbatim}

\item The \verb|operator_symbol| in this case could be any of the several \textbf{operator types} that the programmer can define.

\item It could be \begin{itemize}
		\item + (addition)
		\item \&\& (logical AND)
		\item and so on...
\end{itemize}

\item The \vocab{operands} help the compiler distinguish \textbf{one operator} from \textbf{another}.

\item So, why does C++ provide operators when functions are supported?

	\pbreak

\item Consider a utility class \verb|Date| that encapsulates the day, month, and year:
	\begin{verbatim}
Date holiday(12, 25, 2016);	// Dec 26, 2016
	\end{verbatim}

\item Assuming that you want to \textbf{add a day} and get the \textbf{instance} to contain the next day --- Dec 26, 2016 --- which of the following two options would be more intuitive?

	\begin{enumerate}
		\item Option 1 (using the \textbf{increment operator ++}):
				\begin{verbatim}
++holiday;
				\end{verbatim}
		\item Option 2 (using a \textbf{member function} \verb|Increment()|):
		\begin{verbatim}
holiday.Increment(); // Dec 26, 2016
		\end{verbatim}
	\end{enumerate}

\item Clearly, \fbox{\textbf{Option 1}} scores over method \verb|Increment()|

\item The \fbox{\textbf{operator-based mechanism} facilitates consumption by supplying ease of use and intuitiveness.}

	\pbreak

\item Implementing \textbf{operator $<$} in class \verb|Date| would help you compre \textbf{two instances} of \verb|class Date| like this:

	\begin{verbatim}
if (date1 < date2)
	do something;
else
	do something else;
	\end{verbatim}

\item \textbf{Operators} can be used in more situations than just \textbf{classes} that manage dates.

\item An \textbf{addition operator (+)} in a \textbf{string utility class} such as \verb|MyString| introduced in Listing 9.9 would facilitate easy \textbf{concatentation}:

	\begin{verbatim}
MyString sayHello("Hello ");
MyString sayWorld(" world");
MyString sumThem (sayHello + sayWorld);			// if operator+ were supported by MyString
	\end{verbatim}

\item The effor in implementing relevant operators will be rewarded by the ease of consumption of the class.

\item On a broad level, \textbf{operators} in C++ can be classified into two types:
	\begin{enumerate}
		\item \vocab{Unary operators}
		\item \vocab{Binary Operators}
	\end{enumerate}
\end{itemize}

\newpage

\subsection{Unary Operators}
\begin{itemize}
	\item \textbf{Operators} that function on a \textbf{single operand} are called \vocab{unary operaetors}.

	\item A \textbf{unary operator} that is implemented in the \textbf{global namespace} or as a \textbf{static member function} of a class uses the following structure:
	\begin{verbatim}
return_type operator operator_type (parameter type)
{
	// ... implementation
}
	\end{verbatim}

\item A \textbf{unary operator} that is a \textbf{(non-static)} member of a class has a similar structure but is lacking in parameters, b/c the \textbf{single parameter} that it works upon is the \underline{\textbf{instance of the class itself (*this)}}:

	\begin{verbatim}
return_type operator operator_type ()
{
	// ... implementation
}
	\end{verbatim}
\end{itemize}

\newpage

\subsubsection{Types of Unary Operators}
\begin{tabular}{|p{0.25\textwidth}|p{0.65\textwidth}|}
    \hline
    \textbf{Operator} & \textbf{Name} \\
    \hline
    ++ & Increment \\
    -- & Decrement \\
	* & Pointer dereference \\
	\verb|->| & Member selection \\
	! & Logical NOT \\
	\& & Address-of \\
	- & One's complement \\ 
	+ & Unary plus \\ 
	- & Unary negation \\ 
	Conversion operators & Conversion into other types \\
    \hline
\end{tabular}

\newpage
\subsubsection{Programming a Unary Increment/Decrement Operator}
\begin{itemize}
	\item A \fbox{\textbf{unary prefix increment operator (++)}} is the following syntax \textbf{within the class declaration}:
		\begin{verbatim}
Date& operator ++ (int)
{
	// store a copy of the current state of the object before incrementing day
	Date copy (*this);

	// increment implementation code

	// return state before increment (because, postfix)
	return copy;
}
		\end{verbatim}

	\item The \textbf{prefix} and \textbf{postfix} \underline{\textbf{decrement operators}} have a similar syntax as the \textbf{increment operators}, just that the \textbf{declaration} would contain a $--$ where you see a ++.

	\item \fbox{\textbf{Listing 12.1} Allowing Incrementing and Decrementing Days}

	\item The \textbf{postfix} differs from \textbf{prefix} operator in \textbf{return-type} and \textbf{parameters}.
\end{itemize}

\newpage

\subsubsection{Programming Conversion Operators}
\begin{itemize}
	\item doing \verb|cout << holiday;| will produce an error in absence of conversion operator

		\begin{verbatim}
error: binary '<<' : no operator found which takes a right-hand operand of type 'Date' 
(or there is no acceptable conversion)
		\end{verbatim}

	\item \verb|cout| didn't accept the type of the instance of the class.

	\item We know that \verb|cout| c an work well with a \verb|const char*|:
		\begin{verbatim}
std::cout << "Hello world";	// const char* works!
		\end{verbatim}

	\item So, getting \verb|cout| to work with an \textbf{instance} of type \verb|Date| might be as simple as adding an \textbf{operator} that returns a \verb|const char*| version:
		\begin{verbatim}
operator const char*()
{
	// operator implementation that returns a char*
}
		\end{verbatim}

	\item \fbox{\textbf{Listing 12.2} implementing conversion operator const char* for class Date}

	\item The \textbf{operator} opens up new possibilities toward consuming \verb|class Date|.

	It allows you to even assign an \textbf{instance} of a \verb|Date| directly to a \verb|string|.
\end{itemize}

\pbreak
\begin{tcolorbox}
	\begin{itemize}
	\item Note that such assigments cause \textbf{implicit conversions}, that is, the compiler has used the available \textbf{conversion operator} (in this case \verb|const char*|) there by permitting unintended assignments that get compiled without error.

	\item To \textbf{avoid implicit conversions}, use the keyword \verb|explicit| at the begging of an \textbf{operator declaration}, as follows:

		\begin{verbatim}
explicit operator const char*)_
{
	// conversion code here
}
		\end{verbatim}

\item Using \vocab{explicit} would force the programmer to assert his intention to convert using a \textbf{cast}:
	\begin{verbatim}
// explicit
string strHoliday(static_cast<const char*>(Holiday));
strHoliday = static_cast<const char*>(Date(11, 11, 2016));	

// implicit
string strHoliday (holiday);
strHoliday = Date(11, 11, 2016);
	\end{verbatim}

\item Casting, including \verb|static_cast| is dicussined in Lesson 13
	\end{itemize}
\end{tcolorbox}

\pbreak

\begin{tcolorbox}
	\begin{itemize}
		\item Program as many \textbf{operators} as you think your class would be used with.

			If your application needs an itneger representation of a \verb|Date|, then you may program it as follows:

		\begin{verbatim}
explicit operator int()
{
	return day + month + year;
}
		\end{verbatim}

	\item This would allow an \textbf{instance} of \verb|Date| to be used or transacted as an \textbf{integer}:

	\begin{verbatim}
 FuncTakesInt(static_cast<int>(Date(12, 25, 2016)));
	\end{verbatim}

\item \fbox{\textbf{Listing 12.8} later on also demonstrates \vocab{conversion operators} used with a \textbf{string class}}.
	\end{itemize}
\end{tcolorbox}


\newpage

\subsubsection{Programming Dereference Operator (*) and Member Selection Operator ($->$)}
\begin{itemize}
\item The \vocab{dereference operator (*)} and \vocab{member selection operator} (\verb|->|) are most frequently used in the programming of \vocab{smart pointer classes}.
\item \vocab{Smart pointers} are \textbf{utility classes} that wrap regular pointers and simply memory management by resolving \textbf{ownership} and \textbf{copy issues} using \textbf{operators}. (Lesson 26)

\item This lesson takes a brief look at how \textbf{overloading operators} helps in making \textbf{smart pointers} work.

\pbreak

\item Analyze the use of the \verb|std::unique_ptr| in \fbox{\textbf{Listing 12.3}} and understand how it uses \textbf{operator (*)} and \textbf{operator} \verb|(->)| to help you use the \textbf{smart pointer class} like any normal pointer.

	\pbreak

\item This example demonstrates how a \vocab{smart pointer} allows you to use \textbf{normal pointer syntax} as shown.

\item It's like the variables are just the regular \verb|int*| and \verb|Date*|.

\item \boxednote{
	The \textbf{\fbox{secret lies}} in the \textbf{pointer class} \textbf{std::unique\_ptr} that is \textbf{\underline{smart}} b/c it implements \textbf{operator (*)} and \textbf{operator} \textbf{($->$)}.
}
\item Smart pointers can do much more like de-allocate memory as they go out of scope.

\end{itemize}

\newpage

\subsection{Binary Operators}
\begin{itemize}
	\item \textbf{Operators} that function on \textbf{two operands} are called \vocab{binary operators}.

	\item The \textbf{definition} of a \textbf{binary operator} implemented as a \textbf{global function} or a \textbf{static member function} is the following:

	\begin{verbatim}
return_type operator_type (parameter1, parameter2);
	\end{verbatim}

\item The \textbf{definition} of a \textbf{binary operator} implemented as a \textbf{class member} is 

\begin{verbatim}
return_type operator_type (paramater);
\end{verbatim}

\item The reason the \textbf{class member version} of a \textbf{binary operator} accepts \textbf{only one parameter} is that the \textbf{second parameter} is usually derived from the attributes of the class itself.
\end{itemize}

\newpage

\subsubsection{Type of Binary Operators}
\begin{itemize}
	\item Here is a table of \textbf{binary operators} that can be \textbf{overloaded} or \textbf{redefined}:

		\begin{tabular}{p{0.48\linewidth} p{0.48\linewidth}}
			\hline
			Operator & Name\\
			\hline
			, & Comma \\ 
			!=	&	Inequality \\
			\%	&	Modulus\\
			\%=	&	Modulus/Assignment\\
			\&	&	Bitwise AND\\
			\&\&	&	Logical AND\\
			\&=	&	Bitwise AND/assignment\\
			*	&	Multiplication\\
			*=	&	Multiplication/assignment\\
			+	&	Addition\\
			+=	&	Addition/assignment\\
			-	&	Subtraction\\
			-=	&	Subtraction/assignment\\
			\verb|->|*	&	Pointer-to-member selection\\
			/	&	Division\\
			/=	&	Division/assignment\\
			\verb|<|	& Less than\\
			\verb|<<|	& Left shift\\
			\verb|<<=|	& Left shift/assignment\\
			\verb|<=|	& Less than or equal to\\ 
			=	&	Assignment, Copy Assignment and Move Assignment\\
			==	&	Equality\\
			\verb|>|	& Greater than\\
			\verb|>=|	& Greater than or equal to\\ 
			\verb|>>|	& Right shift\\
			\verb|>>=|	& Right shift/assignment\\
			\verb|^|	&	Exclusive OR\\
			\verb|^=|	&	Exclusive OR/assignment\\
			\verb+|+      & Bitwise inclusive OR \\
			\verb+|=+     & Bitwise inclusive OR/assignment \\
			\verb+||+     & Logical OR \\
			\verb|[]|	&	Subscript operator\\
			\hline
		\end{tabular}
\end{itemize}

\newpage

\subsubsection{Programming Binary Addition (a + b) and Subtraction (a - b) Operators}
\begin{itemize}
	\item Similar to the \textbf{increment/decrement operators}, the \textbf{binary plus} and \textbf{minus}, when defined, enable you to add or subtract the value of a supported data type from an \textbf{object} of the class that \textbf{implements these operators}.

	\item Say for our \verb|class Date| object we want to increase the date by say 5 says instead of just 1 (increment previously).

		\fbox{\textbf{Listing 12.4} Binary Addition Operator}

	\pbreak

\item Note: now we have the \textbf{binary operator} mechanism available to us we can make use of adding strings/instances together that are of class \verb|MyString| type:

	\begin{verbatim}
MyString Hello("Hello ");
MyString World(" World");
MyString HelloWorld(Hello + World);
	\end{verbatim}

\item We do so be \textbf{defining operator (+)}
	\begin{verbatim}
MyString operator+ (const MyString& addThis)
{
	MyString newString;

	if (addThis.buffer != NULL)
	{
		newString.bugger = new char[GetLength() + strlen(addThis.buffer) + 1];
		strcpy(newString.buffer, buffer);
		strcat(newString.buffer, addThis.buffer);
	}	

	return newString;
}
	\end{verbatim}
\end{itemize}

\newpage

\subsubsection{Implementing Addition Assignment (+=) and Subtraction Assignment (-=) Operators}
\begin{itemize}
	\item The \textbf{addition assignment operators} allows syntax such as "a += b;" that allows the programmer to increment the value of an \textbf{object} \verb|a| by an amount \verb|b|.

	\item In doing this, the utility of the \textbf{addition assigment operator} is that it can be \textbf{overloaded} to accept different types of parameter \verb|b|.

	\item \fbox{\textbf{Listing 12.5}} that follows allows you to add an \textbf{integer value} to a \verb|Date| \textbf{object}.

		\pbreak

\item Your \verb|class Date| now allows users to \textbf{add} or \textbf{remove} days from it as if they are dealing with \textbf{integers} using \textbf{addition} or \textbf{subtraction assignment operators} that take an \verb|int| \textbf{\underline{as a parameter}}.

	\item You can even provide \textbf{overloaded verions} of the \textbf{addition assignment operator (+=)} that work with an \textbf{instance} of a \textbf{fictitious class} \verb|Days|:

	\begin{verbatim}
void operator+= (const Days& daysToAdd)
{
	day += daysToAdd.GetDays();
}
	\end{verbatim}
\end{itemize}

\newpage

\begin{tcolorbox}
	The
	\begin{itemize}
		\item multiplication assignment *=
		\item division assignment /=
		\item modulus assignment \%=
		\item subtraction assignment -=
		\item left-shift assignment $<<=$
		\item right-shift assignment $>>=$
		\item XOR assigment \verb|^=|
		\item bitwise inclusive OR assignment \verb+|=+
		\item bitwise AND assignment \&=
	\end{itemize}

	have a similar syntax to the \textbf{addition assignment operator} shown in \textbf{Listing 12.5}
\end{tcolorbox}

\newpage

\subsubsection{Overloading Equality (==) and Inequality (!=) Operators}
\begin{itemize}
	\item What do you expect when the user of \verb|class Date| compares one \textbf{instance} to another:
		\begin{verbatim}
if (date1 == date2)
{
	// do something
}
else {}
		\end{verbatim}

\item In the \textbf{absence of an equality operator ==}, the compiler simply performs a \textbf{binary comparion} of the \textbf{two objects} and returns \verb|true| when they are \textbf{exactly identical}.

	\item A binary comparison checks whether the \textbf{raw bytes} of the two objects are identical in memory.

	\item This is different from \textbf{logical equality}: two objects can represent the same value but still have different byte representations (e.g., if they contain pointers or padding bytes).

	\item For a simple class such as \verb|Date| containing only \verb|int| members, two objects with the same values will typically have the same byte representation and therefore compare as identical. However that isnt always the case for other objects.

	\item Now if you don't even supply the \textbf{equality operator ==} and you compare two objects, then you will get a compiler error (for modern c++ compilers) just like all the other operators on objects. 

		\pbreak

	\item So, This \textbf{binary comparison} will work for instances of classes containing simple data types (like the \verb|Date| class as of now), but it will \textbf{not work} if the class in question has a \textbf{\underline{non-static}} string member (\verb|char*|), such as \verb|MyString| in \textbf{Listing 9.9}.

	\item Wehn \textbf{two instances} of \verb|class MyString| are compared, a \textbf{binary comparison} of the \textbf{member attributes} would actually compare the \textbf{member string pointer values} (\verb|MyString::buffer|).

\item These would \textbf{not be equal} even when the \textbf{strings are identical} in \textbf{content}.

\item Comparisons involving \textbf{two instances} of \verb|MyString| would return \verb|false| consistently.

\item We solve this by \textbf{defining} \textbf{\fbox{comparison operators}} like so:

	\begin{verbatim}
bool operator== (const ClassType& compareTo)
{
	// comparison code here, return true if equal else false
}
	\end{verbatim}

\item The \textbf{\fbox{inequality operator}} can reuse that structure:
\begin{verbatim}
bool operator!= (const ClassType& compareTo)
{
	// comparison code here, return true if not equal else false
}
\end{verbatim}

\item The \fbox{\textbf{inequality operator}} can be the \textbf{\fbox{inverse (logical NOT)}} of the result of the \textbf{equality operator} (obviously)

\item \fbox{\textbf{Listing 12.6} Operators == and !=}

	\end{itemize}

	\newpage

\subsubsection{Overloading $<$, $>$, $<=$, and $>=$ Operators}
\begin{itemize}
	\item Similarly from \textbf{Listing 12.6}, we can change this to the $<$, $>$, $<=$, and $>=$ operators for comparing objects.

	\item We can do those conditional checks respectively akin to:
	\begin{verbatim}
if (date1 < date2) { // ... }
	or
if (date1 <= date2) { // ... }
	or
if (date1 > date2) { // ... }
	or
if (date1 >= date2) { // ... }
	\end{verbatim}

\item \fbox{\textbf{Listing 12.7} Overloading $<$, $>$, $<=$, and $>=$ Operators}
\end{itemize}

\newpage

\subsubsection{Overloading Copy Assignment Operator (=)}
\begin{itemize}
	\item There are times when you want to \textbf{assign the contents} of an \textbf{instance} of a class to another like this:
		\begin{verbatim}
Date holiday(12, 25, 2016);
Date anotherHoliday(1, 1, 2017);
anotherHoliday = holiday;	// uses copy assignment operator
		\end{verbatim}

	\item The \textbf{assignment} invokes the \textbf{default copy assignment operator} that the compiler has built in to your class when you have \textbf{not supplied one}.

	\item Depending on the nature of the class, the \textbf{default copy assignment opeartor} might be inadequate, especially if your class is managing a resource that will not be copied.

	\item Similar problem to the \textbf{default copy constructor} and not actually copying its own resource.

	\item So, to ensure \textbf{deep copies} as with the \textbf{copy constructor}, you need to specify an accompanying \textbf{copy assignment operator}:

\begin{verbatim}
ClassType& operator= (const ClassType& copySource)
{
	if (this != &copySource)	// protection against copy into self
	{
		// copy assignment operator implementation
	}
	return *this;
}
\end{verbatim}

\item \textbf{Deep copies} are important if your class \textbf{encapsulates} a \textbf{raw pointer}, such as \verb|class MyString| shown in \textbf{Listing 9.9}.

\item To \textbf{ensure deep copy} during \textbf{assignments}, define a \textbf{\fbox{copy assignment operator}} as shown in \textbf{\fbox{Listing 12.8}} 

\item \fbox{\textbf{Listing 12.8} A Better 'class MyString' from \textbf{Listing 9.9} with a Copy Assignment Operator =}

	\newpage

\item Notice we did a \vocab{type conversion} (next lesson), which means we can go from\\
\verb|MyString| $\rightarrow$ \verb|cont char*| $\rightarrow$ \verb|output| like so:

\begin{verbatim}
class MyString
{
    char buffer[100];

public:
    operator const char*()
    {
        return buffer;
    }
};

// we can do

MyString str;
printf("%s", str);

// Even though printf expects a const char*, C++ automatically calls:

str.operator const char*()
\end{verbatim}

\item The above \verb|operator const char*| replaced the \verb|GetString()| from before and we used this operator when doing
	\begin{verbatim}
cout << string1 << string2 << endl;
	\end{verbatim}

\pbreak

\item The \textbf{copy assignment operator} is similar \textbf{in function} to a \textbf{copy constructor} and performs a starting check to ensure that the same object is not both the copy source and destination.

\item Look at it for more important details.

\end{itemize}
	\begin{tcolorbox}
		\begin{itemize}
			\item To create a class that \textbf{cannot be copied}, delcare the \textbf{copy constructor} \underline{and} \textbf{copy assignment operator} \verb|private|.

			\textbf{Declaration as } \verb|private| \underline{without implementation} is \textbf{sufficient} for the compiler to \textbf{throw error} on all attempts at \textbf{copying this class} via passing to a \textbf{function by value} or \textbf{assigning one instance} into another.
		\end{itemize}
	\end{tcolorbox}

\newpage

\subsubsection{Subscript Operator ([])}
\begin{itemize}
	\item The \textbf{operator} that allows \textbf{array-style} \verb|[]| \textbf{access} to a \textbf{\underline{class}} is called \vocab{subscript operator}.

	\item The typical syntax of a \fbox{\textbf{subscript operator}} is:

		\begin{verbatim}
return_type& operator [] (subscript_type& subscript);
		\end{verbatim}

	\item So when creating a class such as \verb|MyString| that encapsulates a \textbf{dynamic array} class of characters in a \verb|char* buffer|, a \textbf{subscript operator} makes it really easy to randomly access individual characters in the buffer:

		\begin{verbatim}
class MyString
{
	// ... other members
public:
	/* const */ char& operator [] (int index) /* const */
	{
		// return the char at position index in buffer
	}
};
		\end{verbatim}

	\item \textbf{Listing 12.9} shoes how the \textbf{subscript operator ([])} helps the user in iterating through the chars contained in an \textbf{instance} of \verb|MyString| using \textbf{normal array semantics}.

	\item \fbox{\textbf{Listing 12.9}}

		\pbreak

	\item Friendly recall:
	\begin{verbatim}
	// recall: const here tells us
	// "This member function promises not to modify the object it belongs to."
	int GetLength() const 
	{

	}
	\end{verbatim}

\item Using the keyword \fbox{\textbf{const}} is important when programming \textbf{operators}.

	In \textbf{Listing 12.9} restricted the \textbf{return value} of \textbf{subscript operator []} to \textbf{const char\&}

	\begin{verbatim}
	const char& operator[] (int index) const
	{
		if (index < GetLength())
			return buffer[index];

		return buffer[0]; // failed
	}
	\end{verbatim}

\item The program works and compiles even without the \textbf{const} keywords, yet the reason you have it there is to avoid this code:

	\begin{verbatim}
MyString sayHello("Hello World");
sayHello[2] = 'k';	// error: operator[] is const
	\end{verbatim}

\item By using \textbf{const} you are protecting interal member \verb|MyString::buffer| (so the rightmost \verb|const|) from \textbf{direct modifications} from the outisde via \textbf{operator []}

\begin{itemize}
    \item The \textbf{leftmost} \verb|const| applies to the \textbf{return type}, making the returned object read-only.

    \item The \textbf{rightmost} \verb|const| applies to the \textbf{member function}, indicating that it does not modify the object (\verb|*this|).

    \item The two are \textbf{independent}: a function may have either, both, or neither, depending on its purpose.
\end{itemize}

\begin{verbatim}
char& operator[](int index);                  // neither
const char& operator[](int index);            // left only
char& operator[](int index) const;            // right only
const char& operator[](int index) const;      // both
\end{verbatim}

\item In addition to classifying the \textbf{return value} as \textbf{const}, you can also declare the \textbf{member function} as \textbf{const} (rightmost \verb|const|) to ensure that it \textbf{cannot modify} the class's member attributes.

\end{itemize}

\newpage

\begin{itemize}
	\item When implementing \textbf{\fbox{subscript operators}}, you can imporve on the version shown in \textbf{Listing 12.9}.
	\item The one is an implementation of a \textbf{single subscript operator} that works for both \textbf{reading from} and \textbf{writing to} the slots in the \textbf{dynamic array}.

	\pbreak

	\item You can, however, implement \textbf{two subscript operators} --- one as a \textbf{const} function and the other as a \textbf{non-const} one:
		\begin{verbatim}
char& operator[] (int index);	// use to write / change buffer at index
char& operator[] (int index) const;	// used only for accessing char at index
		\end{verbatim}

	\item The compiler will invoke the \textbf{const function} for \textbf{\underline{read operations}} and the \textbf{non-const} version for operations that \textbf{\underline{write}} into the \verb|MyString| object.

	\item Thus, you can (if you want to) have separate functinalities in the two \textbf{subscript operations}.

		\pbreak

	\item There are other \textbf{binary operators} (listed in Table 12.2) that can be \textbf{redefined} or \textbf{overloaded}, but this book does not dicuss them futher, but it's definition/logic is similar.
\end{itemize}

\pbreak

\begin{itemize}
	\item Other \textbf{operators}, such as the \fbox{\textbf{logical operators}} and the \textbf{\fbox{bitwiseop operators}}, need to be programmed if the purpose of the class would be \textbf{enhanced} by having them.
\end{itemize}

\newpage

\subsection{Function Operator ()}
\begin{itemize}
	\item The \textbf{operator ()} that make \textbf{objects} behave like a \textbf{function} is called a \vocab{function operator}.

	\item They find application in the \textbf{standard template library (STL)} and are typically used in \textbf{STL algorithms}.

	\item Their usage can include \textbf{making decisions}; such \textbf{function objects} are typically called \vocab{unary} or \vocab{binary predicates}, depending on the \textbf{number of operands} they work on.

	\item \fbox{\textbf{Listing 12.10} A Function Object Created Using Operator ()}

	\pbreak

\item Note how the compiler allows the use of \verb|object displayFuncObj| as a \verb|function| in Line 19 ish by \textbf{implicitly converting} what looks like a \textbf{function call} to a \textbf{call} to \textbf{operator()}.
\end{itemize}

\newpage

\subsection{Move Constructor and Move Assignment Operator for High Performance Programming}
\begin{itemize}
	\item The \vocab{move constructor} and the \vocab{move assignment operators} are performance optimization features that have become a part of the \textbf{standard} in C++11, ensuring that \textbf{temporary values} (\vocab{rvalues} that don't exist beyong the statement) are not wastefully copied.

	\item This is particularly useful when handling a \textbf{class} that manages a \textbf{dynamically allocated} resource, such as a \textbf{dynamic array class} or a \textbf{string class}.
\end{itemize}

\pbreak

\subsubsection{The Problem of Unwanted Copy Steps}
\begin{itemize}
	\item Take a look at the \fbox{\textbf{addition operator+}} as implemented in \textbf{Listing 12.4}.

		Notice that is actually \textbf{creates a copy} and \textbf{returns it}, the following lines of code would be valid exmaples of easy string concatenation:
		\begin{verbatim}
MyString Hello("Hello ");
MyString World("World");
MyString CPP(" of C++");
MyString sayHello(Hello + World + CPP);	// operator+, copy constructor
MyString sayHelloAgain("overwrite this");
sayHelloAgain = Hello + World + CPP;	// operator+, copy constructor, copy assignment operator=
		\end{verbatim}

		\item This simple construct that makes concatenating three strings easy uses \textbf{binary addition operator+}:

			\begin{verbatim}
MyString operator+ (const MyString& addThis)
{
	MyString newStr;

	if (addThis.buffer != NULL)
	{
		// copy into newStr
	}

	return newStr;	\\ return copy by value, invoke copy constructor
}
			\end{verbatim}

		\item While making it easy to concatenate the strings, the \textbf{\fbox{addition operator+}} can cause \underline{\textbf{performance problems}}

		\item The creation of \verb|sayHello| requires the execution of the \textbf{addition operator twice}.

		\item Each execution of \verb|operator+| results in the creation of a \textbf{temporary copy} as a \verb|MyString| is returned \textbf{by value}, thus causing the execution of the \textbf{copy constructor}.

		\item The \textbf{copy constructor} executes a \textbf{deep copy} --- to a \textbf{temporary value} that does \textbf{not exist after the expression}.

		\item Thus, the expression results in \textbf{temporary copies} (\fbox{\textbf{rvalues}}, for the pursits) that are \textbf{not ever required after} the statement and hence are performance bottleneck forced by C++. Well, until recently.
\end{itemize}

\newpage
Another explanation of the problem:

\begin{itemize}
    \item The expression \verb|Hello + World + CPP| is evaluated as \verb|(Hello + World) + CPP|, so the addition operator \verb|operator+| is called \textbf{twice}.

\item Each call to \verb|operator+| creates a \textbf{local} \verb|MyString| object (e.g., \verb|newStr|) to hold the concatenated result. Since it is returned \textbf{by value}, a temporary result is produced for use in the surrounding expression.

    \item Since \verb|operator+| returns its result \textbf{by value}, the temporary object is copied (via the \textbf{copy constructor}) to the caller.

    \item The \textbf{copy constructor} performs a \textbf{deep copy}, allocating new memory and copying the string data, which can be expensive.

    \item These temporary objects exist only for the duration of the expression and are immediately destroyed afterwards, making the deep copies unnecessary work.

    \item Consequently, expressions involving repeated \verb|operator+| calls can create multiple temporary objects and deep copies, leading to reduced performance.

\end{itemize}

Note
\begin{itemize}
    \item The object \verb|newStr| is a \textbf{local variable} and is destroyed when \verb|operator+| returns.

    \item Since the function returns \textbf{by value}, the caller cannot receive the local object directly. Instead, a \textbf{new object} is created from \verb|newStr| using the \textbf{copy constructor} (or the move constructor/copy elision in modern C++).
\end{itemize}

\newpage

\begin{itemize}
	\item This problem has now been finally resolved in C++11 in which the compiler specifically rezognizies \textbf{temporaries} and uses \vocab{move constructors} and \vocab{move assignment operators}, where supplied by the \textbf{programmer}.
\end{itemize}

\newpage

\subsubsection{Declaring a Move Constructor and Move Assignment Operator}
\begin{itemize}
	\item The syntax of the \fbox{\textbf{move constructor}} is as follows:

		\begin{verbatim}
class Sample
{
private:
	Type* ptrResource;

public:
	Sample(Sample&& moveSource)	// move constructor, note &&
	{
		ptrResource = moveSource.ptrResource;	// take ownership, start move
		moveSource.ptrResource = NULL;
	}

	Sample& operator= (Sample&& moveSource)	// move assignment operator, note &&
	{
		if (this != &moveSource)
		{
			delete[] ptrResource;	// free own resource
			ptrResource = moveSource.ptrResource;	// take ownership, start move
			moveSource.ptrResource = NULL;	// copy assignment
		}
	}

	Sample();	// default constructor
	Sample(const Sample& copySource);	// copy contructor
	Sample& operator= (const Sample& copySource);	// free move source of ownership
};
		\end{verbatim}
\end{itemize}

\newpage

\subsection*{Quick recall on virtual methods (not for this move constructor/operator section}

\subsubsection*{Virtual Methods}
\begin{itemize}
    \item A \textbf{virtual method} is a member function in a base class that can be overridden by derived classes.

    \item When called through a base class pointer or reference, the correct derived class implementation is selected at \textbf{runtime}.
\end{itemize}

\begin{verbatim}
class Animal
{
public:
    virtual void Speak()
    {
        cout << "Animal sound\n";
    }
};

class Dog : public Animal
{
public:
    void Speak() override
    {
        cout << "Woof\n";
    }
};

int main()
{
    Animal* a = new Dog();
    a->Speak();   // Calls Dog::Speak()
}
\end{verbatim}


\subsubsection*{Virtual Destructor}
\begin{itemize}
    \item A \textbf{virtual destructor} ensures that when a derived object is deleted through a base class pointer, the derived class destructor is called first.

    \item A base class intended for inheritance should usually have a virtual destructor.
\end{itemize}

\begin{verbatim}
class Animal
{
public:
    virtual ~Animal()
    {
        cout << "Animal destroyed\n";
    }
};

class Dog : public Animal
{
public:
    ~Dog()
    {
        cout << "Dog destroyed\n";
    }
};

int main()
{
    Animal* a = new Dog();
    delete a;
}
\end{verbatim}

Output:

\begin{verbatim}
Dog destroyed
Animal destroyed
\end{verbatim}


\subsubsection*{Virtual Constructor}
\begin{itemize}
    \item C++ does \textbf{not support virtual constructors}.

    \item Constructors create objects, and virtual dispatch requires an object to already exist.

    \item Instead, the \textbf{virtual constructor pattern} is implemented using virtual methods such as \verb|clone()|.
\end{itemize}

\begin{verbatim}
class Shape
{
public:
    virtual Shape* clone() = 0;
};

class Circle : public Shape
{
public:
    Shape* clone() override
    {
        return new Circle(*this);
    }
};
\end{verbatim}


\subsubsection*{Pure Virtual Functions and Abstract Classes}
\begin{itemize}
    \item A \textbf{pure virtual function} has no implementation in the base class and is declared using \verb|= 0|.

    \item A class containing at least one pure virtual function is called an \textbf{abstract class}.

    \item Abstract classes cannot be instantiated directly.
\end{itemize}

\begin{verbatim}
class Shape
{
public:
    virtual double Area() = 0;
};

class Circle : public Shape
{
private:
    double radius;

public:
    Circle(double r) : radius(r) {}

    double Area() override
    {
        return 3.14 * radius * radius;
    }
};
\end{verbatim}

This is invalid:

\begin{verbatim}
Shape s;   // ERROR
\end{verbatim}

This is valid:

\begin{verbatim}
Shape* s = new Circle(5);
cout << s->Area();
\end{verbatim}


\subsubsection*{Summary of Virtual Features}

\begin{center}
\begin{tabular}{|c|c|c|}
\hline
Feature & Syntax & Purpose \\
\hline
Virtual method & \verb|virtual void func();| &
Allows derived override at runtime \\
\hline
Override & \verb|void func() override;| &
Explicitly replaces base method \\
\hline
Virtual destructor & \verb|virtual ~Class();| &
Ensures proper cleanup through base pointer \\
\hline
Virtual constructor & Not supported &
Use \verb|clone()| pattern instead \\
\hline
Pure virtual method & \verb|virtual void func() = 0;| &
Forces derived classes to implement \\
\hline
Abstract class &
Contains pure virtual methods &
Cannot create objects directly \\
\hline
\end{tabular}
\end{center}


\subsubsection*{Common Polymorphic Base Class Pattern}

\begin{verbatim}
class Base
{
public:
    virtual ~Base() {}
    virtual void Function() = 0;
};

class Derived : public Base
{
public:
    ~Derived() {}

    void Function() override {}
};
\end{verbatim}

\newpage
(back to move section)
\begin{itemize}
	\item Thus, the \textbf{declaration} of the \textbf{\fbox{move constructor}} and \fbox{\textbf{assignment operator}} are different from the \textbf{regular copy constructor} and \textbf{copy assignment operator} in that the \textbf{input parameter} is a \textbf{type} \verb|Sample&&|.

	\item Additionally, as the \textbf{input parameter} is the \textbf{move-source}, it \textbf{cannot be} a \verb|const| parameter as it is \textbf{modified}.

	\item \textbf{Return values} remain the same, as these are \textbf{overloaded} versions of the \textbf{constructor} and the \textbf{assignment operator}, respectively. 
\end{itemize}

\pbreak

\begin{itemize}
	\item C++11 compilant compilers ensure that for \textbf{rvalue temporaries} the \textbf{move constrcturos} is used instead of the \textbf{copy constructor} and the \textbf{move assignment operator} is invoked instead of the \textbf{copy assignment operator}.

	\item In your implementation of these two, you ensure that instead of \textbf{copying}, you are simply \textbf{moving} the resource from the \textbf{source} to the \textbf{destination}.

	\item \boxednote{
	\textbf{Listing 12.11} 'class MyString' with Move Constructor and Move Assignment Operator in addition to Copy Constructor and Copy Assignment Operator}.
\end{itemize}

\newpage

\begin{itemize}
	\item Note from \textbf{Listing 12.11} C++17 made \vocab{copy elision mandatory} in \textbf{many return-by-value cases}.

	So, we only see as output:
	\begin{verbatim}
Constructor called for: Hello
Constructor called for: World
Constructor called for:  of C++
Constructor called for: overwrite this
operator+ called:
Default constructor called
operator+ called:
Default constructor called
Move assignment op. moves: Hello World of C++
	\end{verbatim}

\item Notice there is no Move constructor cout here.

	\item The reason the \textbf{move constructor} is not called is because C++17 performs \vocab{copy elision} in this situation. The compiler constructs the returned \verb|MyString| object directly in its final destination instead of creating a temporary object and moving it.

\item Therefore, the statement \verb|return newStr;| does not require a copy or move constructor call. The local object \verb|newStr| is optimized away, and its storage is used directly for the returned object.

\item The \textbf{move assignment operator} is still called because the final temporary result of \verb|Hello + World + CPP| must be assigned to the already existing object \verb|sayHelloAgain|.
\end{itemize}

\newpage

\subsection{User Defined Literals}
\begin{itemize}
	\item \textbf{\fbox{Literal constants}} were introduced in Lesson 3. Here are some examples of a few:
	\begin{verbatim}
int bankBalance = 10000;
double pi = 3.14;
char firstAlphabet = 'a';
const char* sayHello = "Hello!";
	\end{verbatim}

\item In the preceding code, 10000, 3.14, 'a', and "Hello!" are all \textbf{literal constants}.

\item C++11 extended literals by allowed you to define your own literals.

\item Example, you want all your temperatures to be stored and operated using a scale called Kevin. You may now declare all your temperatures using a syntax similar to the following:
\begin{verbatim}
Temperature k1 = 32.15_F;
Temperature k2 = 0.0_C;
\end{verbatim}

\item Using \textbf{literals} \verb|_F| and \verb|_C| that \underline{\textbf{you have defined}}, you have made your application a lot simpler to read and therefore maintain.

	\pbreak

\item To \textbf{define your own literal}, you \textbf{define} \verb|operator ""| like this:
\begin{verbatim}
ReturnType operator"" YourLiteral(ValueType value)
{
	// conversion code here
}
\end{verbatim}
\end{itemize}

\pbreak

\begin{tcolorbox}
	\begin{itemize}
		\item Dependeing on the nature of the \textbf{user defined literal}, the \verb|ValueType| parameter would be restricted to one of the following:
			\begin{itemize}
				\item \verb|unsigned long int| for \textbf{integer literal}
				\item \verb|long double| for \textbf{floating point literal}
				\item \verb|char, wchar_t, char16_T, char32_t| for \textbf{character literal}
				\item \verb|const char*| for \textbf{raw string literal}
				\item \verb|const char*| together with \verb|size_t| for \textbf{string literal}
				\item \verb|wchar_t*| together with \verb|size_t| for \textbf{string literal}
				\item \verb|char16_t*| together with \verb|size_t| for \textbf{string literal}
				\item \verb|char32_t*| together with \verb|size_t| for \textbf{string literal}
			\end{itemize}
	\end{itemize}
\end{tcolorbox}

\begin{itemize}
	\item \fbox{\textbf{Listing 12.12} shows a user defined literal that converts types}.
\end{itemize}

\newpage

\subsection{Operators That Cannot Be Overloaded}
\begin{itemize}
	\item \fbox{\textbf{Table 12.3} Operators that CANNOT be Overloaded or Redefined}\\

	\begin{tabular}{p{0.48\linewidth} p{0.48\linewidth}}
		\hline
		Operator & Name\\
		\hline
		.	& Member selection\\
		.*	& Pointer-to-member selection\\
		::	& Scope resolution\\
		? :	& Conditional ternary operator\\
		sizeof	& Gets the size of an object/class type	\\
		\hline
	\end{tabular}

	\pbreak

\item \textbf{DO} program as many operators as would help making using your class easy, but not more.
\item \textbf{DO} mark \textbf{conversion operators} as \verb|explicit| to avoid \textbf{implicit conversions}.
\item \textbf{DO} always program a \textbf{copy assignment operator} (with a \textbf{copy constructor} and \textbf{destructor}) for a class that contains \textbf{raw pointer members}.
\item \textbf{DO} always program a \textbf{move assignment operator} (and \textbf{move constructor}) for classes that manage \textbf{dynamically allocated resources}, such as an array of data, when using a C++11-compliant compiler.

	\pbreak

\item \textbf{DON'T} forget that the compiler provides a \textbf{default copy assignment operator} and \textbf{copy constructor} if you don't supply these, and they \textbf{won't ensure deep copies} of any \textbf{raw pointers} contained within the class.
\item \textbf{DON'T} forget that if you don't supply a \textbf{move assignment operator} or \textbf{move constructor}, the compiler \textbf{does not create these for you}, but isntead falls back on the \textbf{regular copy assingment operator} and \textbf{copy constructor}.
	
\end{itemize}

\newpage

\section{Casting Operators}
\begin{itemize}
\item \textbf{Casting} is the process of converting an object or value from one data type to another, changing how the compiler interprets and uses that value.

\item \textbf{Operators} that change the \textbf{interpretation} of an \underline{\textbf{object}} are called \vocab{casting operators}.
\end{itemize}

\pbreak

\subsection{The Need for Casting}
\begin{itemize}
    \item Old programs and lots of modules working together inevitably call for the need to cast values.

    \item A common example is when using C library functions such as \verb|malloc()|, \verb|calloc()|, and \verb|realloc()|, which return a \verb|void*| pointer. Since \verb|void*| has no associated data type, it must be cast to the appropriate pointer type before use in C++.

\begin{verbatim}
int* arr = (int*)malloc(10 * sizeof(int));

double* values = (double*)calloc(20, sizeof(double));

char* buffer = (char*)realloc(arr, 20 * sizeof(char));
\end{verbatim}

    \item Although explicit casts are required in these cases, C++ programs should generally prefer \verb|new| and \verb|delete| over the C memory allocation functions, as they correctly invoke constructors and destructors for class objects.

	\item C++ devs don't like
		\begin{verbatim}
int* intArray = (int*)staticStr;
		\end{verbatim}

	because this would typically give an error compiled with a C++ compiler:

	\begin{verbatim}
char* staticStr = "Hello World!";
int* intArray = staticStr;	// error: cannot convert char* to int*
	\end{verbatim}
\end{itemize}

\newpage
\subsection{Examples of C++ Casts}

\begin{itemize}
    \item C++ provides four explicit cast operators: 
		\begin{itemize}
			\item \verb|static_cast|, 
			\item \verb|dynamic_cast|, 
			\item \verb|const_cast|, and 
			\item \verb|reinterpret_cast|. 		
		\end{itemize}

		Each is intended for a specific purpose.
\end{itemize}

\newpage 
\paragraph{1. \texttt{static\_cast}}
Used for compile-time conversions between related types.

\begin{verbatim}
class Animal
{
public:
    virtual ~Animal() {}
};

class Dog : public Animal
{
public:
    void Bark()
    {
        cout << "Woof!" << endl;
    }
};

Animal* animal = new Dog();

Dog* dog = static_cast<Dog*>(animal);
dog->Bark();
\end{verbatim}

\newpage 

\paragraph{2. \texttt{dynamic\_cast}}
Used for safe runtime downcasting. Returns \verb|nullptr| if the cast fails.

\begin{verbatim}
Animal* animal = new Dog();

Dog* dog = dynamic_cast<Dog*>(animal);

if (dog != nullptr)
    dog->Bark();
\end{verbatim}

\newpage 

\paragraph{3. \texttt{const\_cast}}
Adds or removes the \verb|const| qualifier.

\begin{verbatim}
const int value = 10;

int* ptr = const_cast<int*>(&value);
\end{verbatim}

\newpage 
\paragraph{4. \texttt{reinterpret\_cast}}
Reinterprets the bits of one pointer type as another.

\begin{verbatim}
int number = 65;

char* ptr = reinterpret_cast<char*>(&number);

cout << *ptr << endl;
\end{verbatim}

\begin{itemize}
    \item \verb|static_cast| performs compile-time conversions between compatible types.

    \item \verb|dynamic_cast| safely converts pointers or references within an inheritance hierarchy by checking the object's type at runtime.

    \item \verb|const_cast| adds or removes \verb|const| or \verb|volatile| qualifiers.

    \item \verb|reinterpret_cast| performs low-level reinterpretation of memory and should be used only when absolutely necessary.
\end{itemize}

\newpage

\subsection{The C++ Casting Operators}
\begin{itemize}
	\item Casting issues are the consequence of long history of code being developed by many different people and is just a natural consequence of not knowing the cast of something by another developer (like in the malloc case).

	\item But what is changed is again the \textbf{interpretation} of the data, which can be dangerous.

	\item C++ supplies a \textbf{casting operator} specific to \textbf{inheritance-based} scenarios \\ 
		(meaning \textbf{\underline{mainly useful for inheritance-based scenarios}}) that did not exist with C programming.

	\item As seen, the \textbf{4 C++ casting operators} are:
		\begin{enumerate}
			\item \verb|static_cast|
			\item \verb|dynamic_cast|
			\item \verb|reinterpret_cast|
			\item \verb|const_cast|
		\end{enumerate}

	\item The usage synatx of the \textbf{casting operators} is consistent:
		\begin{verbatim}
destination_type result = cast_operator<destination_type> (object_to_cast);
		\end{verbatim}
\end{itemize}

\newpage

\subsubsection{Using static\_cast}
\begin{itemize}
	\item \verb|static_cast| is a mechanism that can be used to \textbf{convert pointers} between related types, and perform \textbf{explicit type conversions} for standard data types that would otherwise happen automatiicaly or \textbf{implicitly}.

	\item \verb|static_cast| implements a basic \textbf{\underline{compile-time check}} to ensure that the \textbf{pointer} is being \textbf{cast} to a \textbf{related type}.

	\item C-style casting allows a \textbf{pointer} to \textbf{one object} to be casted to an absolutely \textbf{unrelated type} \underline{without complaint} by the compiler.

		\pbreak

	\item Using a \verb|static_cast|, a \textbf{pointer} can be \fbox{\textbf{upcasted}} to the \fbox{\textbf{base type}} or can be \fbox{\textbf{downcasted}} to the \fbox{\textbf{derived type}}:

		\begin{verbatim}
Base* objBase = new Derived();
Derived* objDer = static_cast<Derived*>(objBase);	// ok!

// class unrelated is not related to Base
Unrelated* notRelated = static_cast<Unrelated*>(objBase);	// error
// The cast is not permitted as types are unrelated
		\end{verbatim}


		\begin{itemize}
    \item The following example demonstrates both an \textbf{upcast} and a \textbf{downcast}:

\begin{verbatim}
Derived* objDer = new Derived();

Base* objBase = objDer;
// upcast: Derived* -> Base*

Derived* objDer2 = static_cast<Derived*>(objBase);
// downcast: Base* -> Derived*
\end{verbatim}

    \item An \textbf{upcast} converts a derived class pointer into a base class pointer. This is safe because every derived object is also a base object.

    \item A \textbf{downcast} converts a base class pointer into a derived class pointer. This is only safe if the object being pointed to is actually an instance of the derived class.

    \item A \verb|static_cast| does not perform a runtime check during a downcast. If the object is not actually of the derived type, using the resulting pointer can cause undefined behavior. For safer runtime-checked downcasts, use \verb|dynamic_cast| with polymorphic base classes.
\end{itemize}
\end{itemize}

\pbreak

\begin{tcolorbox}
	\begin{itemize}
		\item Casting a \verb|Derived*| to a \verb|Base*| is called \textbf{\textit{upcasting}} and \underline{can be done} without any \textbf{explicit casting operator}:

			\begin{verbatim}
Derived objDerived;
Base* objBase = &objDerived;	// ok!
			\end{verbatim}

		\item Casting a \verb|Base*| to a \verb|Derived*| is called \textbf{\textit{downcasting}} and \underline{cannot} be done without usage of \textbf{explicit casting operator}:

	\begin{verbatim}
Derived objDerived;
Base* objBase = &objDerived;	// upcast -> ok!!!
Derived* objDer = objBase;	// error: downcast needs explicit cast
	\end{verbatim}


	\end{itemize}
\end{tcolorbox}

\pbreak

\begin{itemize}
	\item However, note that \verb|static_cast| verifies only that the \textbf{pointer types} are \textbf{\underline{related}}.

	It does \textbf{\underline{not}} perform any runtime checks.

\item So, with \verb|static_cast|, a programmer could still get away with this bug:

\begin{verbatim}
Base* objBase = new Base();
Derived* objDer = static_cast<Derived*>(objBase);	// still no errors!
\end{verbatim}

\item Here, \verb|objDer| actually \textbf{points} to a \textbf{\underline{partial Derived object}} because the \textbf{object} being \textbf{pointed to} is actually a \verb|Base()| type.

\item Because \verb|static_cast| performs \textbf{only a compile-time check} of verifying that the types in question \textbf{are related} and \textbf{does not} perform a \textbf{runtime check}, a call to \verb|objDer->DerivedFunction()| would get compiled, but probably result in unexpected behavior during \textbf{runtime}.
\end{itemize}

\newpage

\begin{itemize}
	\item Apart from helping in \textbf{upcasting} or \textbf{downcasting}, \verb|static_cast| can, in many cases, help make \textbf{implicit casts explicit} and bring them to the attention of the programmar or reader:

	\begin{verbatim}
double Pi = 3.14;
int num = static_cast<int>(Pi);	// making an otherwise implicit cast, explicit
	\end{verbatim}

\item This makes things explicit to the reader, and lower the likelihood of reader error from implicit conversions.

	\pbreak

\item You would also need to use \verb|static_cast| when using \textbf{conversion operators} or \textbf{constructors} that have been declared using keyword \verb|explicit|.

\item \textbf{Avoiding implicit conversions} via keyword \verb|explicit| was discussed in \textbf{Lesson 9} and \textbf{Lesson 12}.
\end{itemize}

\newpage

\subsubsection{Using dynamic\_cast and Runtime Type Identification}
\begin{itemize}
	\item \vocab{Dynamic\_casting} is the \textbf{opposite} of \textbf{static\_casting} and executes the \textbf{cast} at \underline{\textbf{runtime}}---that is, at \textbf{application execution time}.

	\item The result of a \verb|dynamic_cast| \textbf{operation} can be checked to see whether the attempt at \textbf{casting succeeded}.

	\item Here is typical usage:
		\begin{verbatim}
destination_type* Dest = dynamic_cast<class_type*>(Source);

if (Dest)			// Check for success of the casting operation
	Dest->CallFunc();
		\end{verbatim}

	\item For example:

		\begin{verbatim}
Base* objBase = new Derived();

// perform a downcast
Derived* objDer = dynamic_cast<Derived*>(objBase);

if (objDer)	// check for success of the cast	
	objDer->CallDerivedFunction();
		\end{verbatim}

	\item As shown, given a \textbf{pointer} to a \textbf{base-class object}, we can resort to a \verb|dynamic_cast| to \textbf{verify the type} of the \textbf{destination object} being \textbf{pointed to} before proceeding to use the pointer as such.

\item Note that in the code snippet it is apparent that the \textbf{destination object} is a \verb|Derived| \textbf{type}.

	So, the sample if of demonstrative value only.

\item Yet, this is not always the case---when a \textbf{pointer} of \textbf{type} \verb|Derived*| is passed to a functino that accepts \textbf{type} \verb|Base*|.

\item The function can use \verb|dynamic_cast| given a \textbf{base-class pointer type} to detected type and then perform operations specific to the types detected.

\item Thus, \verb|dynamic_cast| helps determine the \textbf{type} at \textbf{\underline{runtime}} and use a \textbf{casted pointer} when it is safe to do so.

\item Therefore, this mechanism of identifying the \textbf{type} of the \textbf{object} at \textbf{\underline{runtime}} is called \\ 
	\vocab{runtime identification (RTTI)}.

\item \fbox{\textbf{Listing 13.1} Using Dynamic Casting to Tell Whether a Fish Object is a Tuna or a Carp}

\item Note: the \textbf{return value} of a \verb|dynamic_cast| \textbf{operation} should always be checked for \textbf{validity}.

\item It is \textbf{NULL} when the \textbf{\underline{cast fails}}.

\end{itemize}

\newpage

\subsubsection{Using reinterpret\_cast}
\begin{itemize}
	\item \verb|reinterpret_cast| is the \textbf{closest} a C++ \textbf{casting operator} gets to the C-style cast.

	\item It forces a \textbf{reinterpretation of type} using syntax like below (similar to the \textbf{unrelated casting} option C):
		\begin{verbatim}
Base *objBase = new Base();
Unrelated* notRelated = reinterpret_cast<Unrelated*>(objBase);
// The code above compiles, but is not good programming!
		\end{verbatim}

	\item The \textbf{cast} actually \textbf{forces} the compiler to accept situations that \verb|static_cast| would normally not permit.

	\item Sometimes used in cases like in drivers where data needs to be converted to a simple type than an API can accept---like some OS-level APIs require data to be sent as a \textbf{BYTE array}, that is, \verb|unsigned char*|):

		\begin{verbatim}
SomeClass* object = new SomeClass();
// Need to send the object as a byte-stream...
unsigned char* bytesFoAPI = reinterpret_cast<unsigned char*>(object);
		\end{verbatim}

	\item  The \textbf{cast} used in the preceding code has \textbf{\underline{not changed the binary representation}} of the \textbf{source object} and has effectively cheated the compiler into allowing the programmer to peek into \textbf{individual bytes} contained by an \textbf{object} of \textbf{type} \verb|SomeClass|.

	\item B/c this comprises \textbf{type safetly}, \verb|reinterpret_cast| is a \textbf{last resort} in performing an otherwise unsafe (and nonportable) conversion.
\end{itemize}

\newpage

\subsubsection{Using \texttt{const\_cast}}
\begin{itemize}
    \item \vocab{const\_cast} enables you to add or remove the \verb|const| qualifier from an object or pointer.

    \item It is primarily used when interfacing with legacy code or APIs that do not correctly use \verb|const| qualifiers.

    \item Removing \verb|const| is only safe if the original object was \textbf{not} declared as \verb|const|. Modifying an object that was originally declared \verb|const| results in \textbf{undefined behavior}.

    \item Example:
\end{itemize}

\begin{verbatim}
void Print(char* str)
{
    cout << str << endl;
}

const char* message = "Hello";

// Remove constness to call a legacy function
Print(const_cast<char*>(message));
\end{verbatim}

\begin{itemize}
    \item Since \verb|Print()| does not modify the string, this use of \verb|const_cast| is safe. However, if \verb|Print()| attempted to modify the string, the program would have undefined behavior.
\end{itemize}

\newpage

\subsection{Problems with the C++ Casting Operators}
\begin{itemize}
	\item Not all C++ loves the casting operators given by C++: they can be cumbersome like example:

	\begin{verbatim}
double Pi = 3.14;

// C++ style cast: static_cast
int num = static_cast<int>(Pi);	// result: num is 3

// C-style cast
int num2 = (int)Pi;	// result: num is 3

// leave casting to the compiler
int num3 = Pi;	// result: num is 3
	\end{verbatim}

\item In all three cases, the programmer achieved the same result.

\item In practice the second option is most prevalent, followed by the third.

\item The first makes things difficult to read.

\pbreak

\item Another case is where C can handle its casting:

\begin{verbatim}
// using static_cast
Derived* objDer = static_cast<Derived*>(objBase);

// But, this works just as well...
Dervied* objDerSimple = (Derived*)objBase;
\end{verbatim}

\item Thus, the adv of using \verb|static_cast| is overshadowed by its syntax handle.

\item Etc for the others, except the \verb|dynamic_cast| is benefical for modern C++ applications.

\item Best practice is to \textbf{avoid casting} as far as possible.
\end{itemize}

\newpage

\begin{itemize}
	\item \textbf{DO} remember that casting a \verb|Dervied*| to a \verb|Base*| is called \vocab{upcasting} and this is \textbf{\underline{safe}}.
	\item \textbf{DO} remember that casting a \verb|Base*| to a \verb|Dervied*| is called \vocab{downcasting} and this is \textbf{\underline{unsafe}} unless you use \verb|dynamic_cast| and \textbf{check for success}.
	\item \textbf{DO} remember that the objective of creating an \textbf{inheritance hierarchy} is typically having \textbf{virtual functions} that when invoked using \textbf{base class pointers} ensure that the avilable \textbf{derived class} versions are invoked (like the \textbf{virtual destructor} most importantly).

	\pbreak

	\item \textbf{DON'T} forget to \textbf{check the pointer} for \textbf{validity} after using \verb|dynamic_cast|.
	\item \textbf{DON'T} design your application around \fbox{\textbf{RTTI}} using \verb|dynamic_cast|.

\end{itemize}

\newbp

\section{An Introduction to Macros and Templates}
\subsection{The Preprocessor and the Compiler}
\begin{itemize}
	\item \textbf{Lesson 2} had introduced the \fbox{\textbf{preprocessor}}, which is what that runs \textbf{before} the compiler starts.

	\item In other words, the \textbf{preprocessor} actually decides what is compiled on the basis of how you instruct it.

	\item \textbf{Preprocessor directives} all start with the \verb|#| sign:

		\begin{verbatim}
// instruct preprocessor to insert contents of iostream here
#include <iostream>

// define a macro constant
#define ARRAY_LENGTH 25
int numbers[ARRAY_LENGTH];	// array of 25 integers

// define a macro function
#define SQUARE(x) ((x) * (x))
int TwentyFive = SQUARE(5);
		\end{verbatim}

	\item The lesson focuses on \textbf{two types} of \textbf{preprocessor directives}:
		\begin{itemize}
			\item one using \verb|#define| to define a \textbf{constant}
			\item another using \verb|#define| a \textbf{macro function}
		\end{itemize}

	\item Both these \textbf{directives} tell the \textbf{preprocessor} to replace every instance of the \textbf{macro} (ARRAY\_LENGTH or SQUARE) with the value they define (aka text substitute)
\end{itemize}

\newpage

\subsection{Using Macro \#define to Define Constants}
\begin{itemize}
	\item The syntax of using \#define to compose a \textbf{constant} is simple:
		\begin{verbatim}
#define indentifier value
	
// Example:
#define ARRAY_LENGTH 25
		\end{verbatim}

	\item So after the \textbf{preprocessor} runs we can see text substitutes like
	\begin{verbatim}
int numbers [ARRAY_LENGTH] = {0}; 

// gets converted to

int numbers [25] = {0}; 
	\end{verbatim}

\item I am not going over code for this trivial thing, but note that the \textbf{preprocessor} \underline{does not always check} for \textbf{correctness} of the substitution (but the compiler always does).

\end{itemize}

\newpage

\subsubsection{Using Macros to Protection against Multiple Inclusion}
\begin{itemize}
	\item C/C++ programmers typically \underline{\textbf{declare}} their classes and functions in \vocab{.h} filess called \vocab{header files}.

	\item The respective functions are \underline{\textbf{defined}} in \vocab{.cpp files} called \vocab{source files}, and they typically include the \textbf{header files} using the \verb|#include <header>| \textbf{preprocessor directive}.

	\item A tricky part for any \textbf{circular dependency} of \textbf{header files} would entail \verb|forward declarations| (it's best to not have circular dependency):

	\item A.h
	\begin{verbatim}
#ifndef A_H
#define A_H

class B;  // Forward declaration

class A {
public:
    B* b;   // OK: pointer or reference
};

#endif
	\end{verbatim}

\item B.h
	\begin{verbatim}
#ifndef B_H
#define B_H

class A;  // Forward declaration

class B {
public:
    A* a;
};

#endif
	\end{verbatim}

\item A.cpp

	\begin{verbatim}
#include "A.h"
#include "B.h"
	\end{verbatim}

\item B.cpp

	\begin{verbatim}
#include "B.h"
#include "A.h"
	\end{verbatim}
\end{itemize}


\newpage

\begin{itemize}
	\item So with many header files potentially including many others and having \textbf{multiple inclusions} (and not necessarily as the circular dependency case, but more and very likely in general), we use \textbf{macros} in conjunction with \textbf{preprocessor directives} \verb|#ifndef| and  \verb|#endif|

	\item header1.h that includes \verb|<header2.>| looks like the following:
		\begin{verbatim}
#ifndef HEADER1_H_		// multiple inclusion guard
#define HEADER1_H_		// preprocessor will read this and following lines once
#include <header2.h>

class Class1 
{
	// class members
};
#endif	// end of header1.h
		\end{verbatim}

	\item header2.h looks similar, but is flipped:
		\begin{verbatim}
#ifndef HEADER2_H_		// multiple inclusion guard
#define HEADER2_H_		// preprocessor will read this and following lines once
#include <header1.h>

class Class2 
{
	// class members
};
#endif	// end of header2.h
		\end{verbatim}

	\item the \verb|#ifndef| is read as "if-not-defined" and is a \textbf{preprocessing conditional command}, instructing the \textbf{preprocessor} to continue only if the identifier \textbf{has not been defined}.

	\item \verb|#endif| marks the \textbf{end} of this \textbf{conditional processing instruction} for the \textbf{preprocessor}.

\item It prevents the header file from being processed more than once during a single compilation, avoiding duplicate declarations/definitions and helping prevent errors caused by multiple inclusion.

\item This mechanism is simply the \underline{\textbf{one of the most frequently used macro-based functionalities}} in the world of C and C++ programming.

\end{itemize}

\newpage

\subsection{Using \#define to Write Macro Functions}
\begin{itemize}
	\item The capability of the \textbf{preprocessor} to simply replace text elements identified by a \textbf{macro} often results it in being used to write \textbf{simple functions}, for example:

		\begin{verbatim}
#define SQUARE(x) ((x) * (x))

// or 

#define PI 3.1416
#define AREA_CIRCLE(r) (PI*(r)*(r))
		\end{verbatim}

	\item \textbf{Macro functions} are expanded \textbf{\underline{inline}} before compilations and hence can help improve code performance in certain cases.

\item \fbox{\textbf{Listing 14.2} Using Macro Functions}
\end{itemize}

\newpage
\subsubsection{Why All the Parentheses?}

\begin{itemize}
    \item For example:
\begin{verbatim}
#define AREA_CIRCLE(r) (PI*(r)*(r))
\end{verbatim}

    \item The reason for the many parentheses is that a \textbf{macro} is expanded by the \textbf{preprocessor} using simple \textbf{text substitution}; it does not understand operator precedence.

    \item For example, if we call
\begin{verbatim}
AREA_CIRCLE(a + b)
\end{verbatim}
    then the expansion becomes
\begin{verbatim}
(PI*(a + b)*(a + b))
\end{verbatim}
    which correctly computes $\pi(a+b)^2$.

    \item Without the parentheses,
\begin{verbatim}
#define AREA_CIRCLE(r) PI*r*r
\end{verbatim}
    the same call expands to
\begin{verbatim}
PI*a + b*a + b
\end{verbatim}
    which is interpreted according to operator precedence and produces a completely different result.

	\pbreak

\item These brackets automatiicaly result in the calculation of an accurate area, making your \textbf{macro code} independent of \textbf{operator precedence} and the effects thereof.
\end{itemize}

\newpage

\subsubsection{Using Macro 'assert' to Validate Expressions}
\begin{itemize}
	\item The \vocab{assert} \textbf{macro} enables you to \textbf{check that a condition is true at runtime}; if the condition evaluates to \verb|false|, the program is terminated and a diagnostic message indicating the failed assertion is displayed.

	\item In large applications this is a decent way to test functionality.

	\item To use \verb|assert| you include \verb|<assert.h>| and the syntax is as follows:
		\begin{verbatim}
assert (expression that evaluations to true or false);
		\end{verbatim}

	\item A sample use of \verb|assert()| that validates the contents of a \textbf{pointer} is:
		\begin{verbatim}
#include <assert.h>
int main()
{
	char* sayHello = new char[25];
	assert(sayHello != NULL);	// throws a message if pointer is NULL;

	// other code

	delete[] sayHello;
	return 0;
}
		\end{verbatim}

	\item \verb|assert()| ensures that you are notified if the \textbf{pointer} is \textbf{invalid}.

	\item This makes it a handy debugging feature to say check for validty of the parameters as well.

	\item This is \underline{\textbf{highly recommended}} and helps you improve the quality of your code over the long term.
\end{itemize}

\pbreak

\begin{tcolorbox}
	\begin{itemize}
		\item \verb|assert()| is typically \textbf{disabled} in \vocab{release modes} and provides you with an \textbf{error message} or information only in the \vocab{debug build} of most development environments.

		\item Additionally, some environments have implemented this as a \textbf{function}, \underline{not} as a \textbf{macro}.
	\end{itemize}
\end{tcolorbox}

\pbreak

\begin{tcolorbox}
	\begin{itemize}
		\item An \textbf{assert} does \underline{not} make it to the \textbf{release build}, it is important to ensure that checks are critical to the functioning fo your application (for ex, the \textbf{return value} of a \verb|dynamic_cast| operation) are still performed using an \textbf{if-statement}.

		\item an \textbf{assert} assits you in problem detection; it's \underline{not} something to \textbf{replace pointer checks} in the code.
	\end{itemize}
\end{tcolorbox}

\newpage

\subsubsection{Adv and Disadv of Using Macro Functions}
\begin{itemize}
    \item \textbf{Advantages}
    \begin{itemize}
        \item \textbf{Eliminates function call overhead.} A macro is expanded by the \textbf{preprocessor} before compilation, so the generated code is inserted directly at each point of use rather than invoking a function at runtime.

        \item \textbf{May improve performance for small operations.} Since there is no function call, stack frame creation, or parameter passing, macro expansion can be more efficient for very simple operations.

        \item \textbf{Can operate on different data types.} Because a macro performs only \textbf{text substitution}, it is not restricted to a particular type (although this also introduces disadvantages).

        \item \textbf{Useful for compile-time constants and conditional compilation.} Macros are commonly used to define symbolic constants, enable or disable debugging code, and control platform-specific compilation using preprocessing directives.

        \item \textbf{Reduces repetitive code.} Frequently used expressions can be written once and reused throughout a program, improving consistency and maintainability when used appropriately.
    \end{itemize}

    \item \textbf{Disadvantages}
    \begin{itemize}
        \item \textbf{No type checking.} The \textbf{preprocessor} simply performs text substitution and has no knowledge of C/C++ types, allowing type-related errors to pass unnoticed until compilation or execution.

        \item \textbf{Arguments may be evaluated multiple times.} A macro parameter appearing more than once in the macro body causes the corresponding argument to be substituted multiple times, potentially producing unintended side effects.
        \begin{verbatim}
#define SQUARE(x) ((x) * (x))

SQUARE(i++)   // expands to ((i++) * (i++))
        \end{verbatim}

        \item \textbf{Operator precedence can produce incorrect results.} Since macro expansion is purely textual, failing to surround parameters and the entire expression with parentheses may lead to unexpected evaluations.
        \begin{verbatim}
#define SQUARE(x) x*x

SQUARE(a + b)   // expands to a + b * a + b
        \end{verbatim}

        \item \textbf{More difficult to debug.} The compiler only sees the expanded code rather than the original macro invocation, making compiler errors and debugger output harder to interpret.

        \item \textbf{Code size may increase.} Every macro invocation expands into a copy of its replacement text, potentially increasing the size of the generated executable.

        \item \textbf{Lack of scope.} Unlike variables and functions, macros are processed before compilation and are not governed by C++ namespaces or normal scope rules, increasing the risk of name collisions.

        \item \textbf{Modern alternatives are usually preferable.} In modern C++, \textbf{inline functions}, \textbf{constexpr}, templates, and `const` variables provide similar performance while offering type safety, proper scoping, and better debugging support.
    \end{itemize}
\end{itemize}

\newpage

\subsection*{Do's and Don't's}
\begin{tcolorbox}
\begin{itemize}
	\item \textbf{DO} program your own \textbf{macro functions} as \underline{infrequently} as possible.
	\item \textbf{DO} use \verb|const| variables where you can \underline{instead of} \textbf{macros}
	\item \textbf{DO} remember that \textbf{macros} are \underline{not type safe} and the \textbf{preprocessor} performs \textbf{no type checking}.

	\pbreak

\item \textbf{DON'T} forget to envelop every variable in a \textbf{macro function definition} with brackets
\item \textbf{DON'T} forget to insert multiple \textbf{inclusion guards} using \verb|#ifndef|, \verb|#define|, and \verb|#endif| in your \textbf{header files}.
\item \textbf{DON'T} forget to sprinkle your code with \textbf{generous} number of \verb|assert()| statements---these don't make it to the \vocab{release version} are are good at improving \textbf{quality} of your code.
\end{itemize}
\end{tcolorbox}

\newpage

\subsection{An Introduction to Templates}
\begin{itemize}
	\item You begin the \textbf{declaration} of a \vocab{template} using the keyword \verb|template| followed by a \textbf{type parameter list}:

		\begin{verbatim}
template <parameter list>
template function / class declaration...
		\end{verbatim}

	\item The keyword \verb|template| marks the \textbf{start} of a \textbf{template declaration} and is followed by the \textbf{template parameter list}.

	\item This \textbf{parameter list} contains the keyword \vocab{typename} that \textbf{defines} the \textbf{template parameter} \verb|objType|, making it a \underline{\textbf{placeholder}} for the \underline{\textbf{type}} of the \textbf{object} that the \textbf{template} is being instantiated for.

		\begin{verbatim}
// a template function
template <typename T1, typename T2 = T1>
bool TemplateFunction(const T1& param1, const T2& param2);

// a template class
template <typename T1, typename T2 = T1>
class MyTemplate 
{
private:
	T1 member1;
	T2 member2;

public:
	T1 GetObj1() { return member1; }
	// ... other members
};
		\end{verbatim}

	\item What you see is a \vocab{template function} and a \vocab{template class}, each taking \textbf{two template parameters} \verb|T1| and \verb|T2|, where \verb|T2| has been given the \textbf{default type} as that of \verb|T1|.
\end{itemize}

\newpage

\subsubsection{The Different Types of Template Declarations}
\begin{itemize}
	\item A \textbf{\fbox{template declaration}} can be 
		\begin{itemize}
			\item a declaration of defintion of a function
			\item a declaration of defintion of a class
			\item a definition of a member function or a member class of a class template
			\item a definition of a statid data member of a class template
			\item a definition of a static data member of a class nested within a class template
			\item a definition of a member template of a class or class template
		\end{itemize}
\end{itemize}

\newpage

\subsubsection{Template Functions}
\begin{itemize}
\item Imagine a \textbf{function} that would adapt itself to suit parameters of \textbf{different types}.
\item Such a function is possible using \textbf{templates}.

\item Let's look at a sample \textbf{template declaration} that returns the greater of two supplied parameters:

	\begin{verbatim}
template <typename objType>
const objType& GetMax(const objType& value1, const objType& value2)
{
	if (value1 > value2)
		return value1;
	else
		return value2;
}

// usage:

int num1 = 25;
int num2 = 40;

int maxVal = GetMax<int>(num1, num2);

double double1 = 1.1;
double double2 = 1.001;
double maxVal = GetMax<double>(double1, double2);
	\end{verbatim}

\item Note the \verb|<int>| and \verb|<double>|.
\item It effectively \textbf{defines} the \textbf{template parameter} \verb|objType| as \verb|int| or \verb|double| respectively.

\item The proceding code leads to the compiler generating two versions of the \textbf{template function} \verb|GetMax|, which can be visualized as the following:

\begin{verbatim}
const int& GetMax(const int& value1, const int& value2)
{
	// ...
}

const double& GetMax(const double& value1, const double& value2)
{
	// ...
}
\end{verbatim}

\item In reality, however, \textbf{template functions} \underline{don't} necessarily need an \textbf{accompanying type specifier}. So the following function call works perfectly well:

	\begin{verbatim}
int maxVal = GetMax(num1, num2);
	\end{verbatim}

\item Compilers are intelligent enough to understand that the \textbf{template function} is being invoked for the \textbf{integer type} in this case.

\item With \textbf{template classes}, \underline{however}, you need to \underline{\textbf{explicitly specify type}}, as shown in \fbox{\textbf{Listing 14.3}}
\end{itemize}

\newpage

\subsection{Templates and Type Safety}
\begin{itemize}
	\item \textbf{Template functions} \verb|DisplayComparison()| and \verb|GetMax()| shown in \textbf{Listing 14.3} are \underline{\textbf{type safe}}.

	\item This means that they would \underline{not} allow a meaningless call like this one:
		\begin{verbatim}
DisplayComparison(num1, name1);
		\end{verbatim}

	\item This would immediately result in a \textbf{compile failure} (int and a string)
\end{itemize}

\newpage 

\subsubsection{Template Classes}
\begin{itemize}
	\item Recall: \textbf{classes} are design blueprints, and the real-world respresentation of a \textbf{class} is an \textbf{object}.

	\item \textbf{Template classes} are the \textbf{templatized versions} of C++ classes. They are blueprints of blueprints.

	\item When using a \textbf{template class}, you are given the option to specify the \textbf{"type"} you are specializing the \textbf{class} for.

	\item This enables you to create some humans with \textbf{template parameters} \verb|Age| as a \verb|long long|, some with \verb|int|, and some wiht \verb|Age| as an integer of type \verb|short|.

		\pbreak

	\item A simple \textbf{template class} that uses a \textbf{single parameter} \verb|T| to hold a \textbf{member variable} can be written as the following:
		\begin{verbatim}
template <typename T>
class HoldVarTypeT
{
private:
	T value;

public:
	void SetValue (const T& newValue) { value = newValue; }
	T& GetValue() { return value; }
};
		\end{verbatim}

	\item The \textbf{type} of the \textbf{variable} is \verb|T|, and that is assigned at the time the \textbf{template} is used, that is, \vocab{instantiated}.

		\pbreak

	\item So, let's look at a sample usage of this \textbf{template class}:
		\begin{verbatim}
HoldVarTypeT <char> holdInt;	// template instantiation for int
holdInt.SetValue(5);
cout << "The value stored is: " << holdInt.GetValue() << endl;
		\end{verbatim}

	\item You have used thnis \textbf{template class} to hold an retrieve an \textbf{object} of \textbf{type} \verb|int|; that is, the \verb|Template| class is \textbf{instantiated} for a \textbf{template parameter} of \textbf{type} \verb|int|.

	\pbreak

\item You can do this for other types like for \textbf{character strings}:
	\begin{verbatim}
HoldVarTypeT <char*> holdStr;
holdStr.SetValue("Sample string");
cout << "The value stored is: " << holdStr.GetValue() << endl;
	\end{verbatim}

\item Thus, the \textbf{template class} defines a pattern for classes and helps implement that pattern on \textbf{different data types} that the \textbf{template} may be \textbf{instantiated} with.

\end{itemize}

\pbreak

\begin{tcolorbox}
	\begin{itemize}
		\item \textbf{\fbox{Template classes}} can be \textbf{instanstatied} with types other than simple ones like \verb|int| or classes supplied by the standard library.

		\item You may \textbf{instantiate} a \textbf{template} using a \textbf{class} defined by you..

		\item For example:

			\begin{verbatim}
HoldVarTypeT<Human> holdHuman;
holdHuman.SetValue(firstMan);
holdHuman.GetValue().IntroduceSelf();
			\end{verbatim}

			\item For \textbf{class templates}, the template arguments must generally be specified inside \verb|<>| when creating an object, such as \verb|HoldVarTypeT<Human> holdHuman;|, because the compiler needs to know what type should replace the template parameter.

				However, since \textbf{C++17}, \textbf{Class Template Argument Deduction (CTAD)} allows the compiler to infer the template arguments from constructor arguments in some cases, allowing the \verb|<>| to be omitted.
	\end{itemize}
\end{tcolorbox}

\newpage

\subsubsection{Declaring Templates with Mutliple Parameters}
\begin{itemize}
	\item The \vocab{template parameter list} cna be expanded to \textbf{declare multiple parameters} separated by a comma.

	\item So yuo can declare a generic class that holds a pair of \textbf{objects} that can be of \textbf{differing types}:

		\begin{verbatim}
template <typename T1, typename T2>
class HoldsPair
{
private:
	T1 value1;
	T2 value2;

public:
	// constructor that initializes member variables
	HoldsPair(const T1& val1, const T& val2)
	{
		value1 = val1;
		value2 = val2;
	}

	// ... other member functions
};
		\end{verbatim}

	\item In this example, \verb|class HoldsPair| accepts \textbf{two template parameters} named \verb|T1| and \verb|T2|.

	\pbreak

	\item We can use this class to hold \textbf{two objects} of the \textbf{same type} or od \textbf{different types} as you can see here:

		\begin{verbatim}
// a template instantiation that pairs an int with a double
HoldsPair<int, double> pairIntDouble (6, 1.99);

// a template instantiation that pairs an int with an int
HoldsPair<int, int> pairIntInt (6, 500);
		\end{verbatim}
\end{itemize}

\newpage

\subsubsection{Declaring Templates with Default Parameters}
\begin{itemize}
	\item We can modify the previous version of \verb|HoldsPair <...>| to declare \verb|int| as the \\
		\vocab{default template parameter type}:

\begin{verbatim}
template <typename T1=int, typename T2=int>
class HoldsPair 
{ 
	// ... method declarations	
};
\end{verbatim}

\item This is simlilar in \textbf{construction} to \textbf{functions} that \textbf{define default input parameter values} except for the fact that, in this case, we define \textbf{default \textit{types}}.

\item The second usage of \verb|HoldsPair| can thus be compacted to:
	\begin{verbatim}
// pair an int with an int (default type)
HoldsPair<> pairsInts(6, 500);
	\end{verbatim}
\end{itemize}

\newpage

\subsubsection{Sample Template 'class$<>$' HoldsPair}
\begin{itemize}
	\item \fbox{\textbf{Listing 14.4} A Template Class with a Pair of Member Attributes}
\end{itemize}

\newpage

\subsubsection{Template Instantiation and Specilization}
\begin{itemize}
	\item A \textbf{template class} is a \textbf{blueprint} of a \textbf{class} (a blueprint of a blueprint), and therefore doesn't truly exist for the compiler before it has been used in one form or another.

	\item That is, as far as the compiler is concerned, a \textbf{template class} you \textbf{define} but don't consume in code that is simply \underline{ignored}.

	\item However, \underline{\textit{once you instantiate}} a \textbf{template class}, like \verb|HoldsPair|, by \textbf{supplying template arguments} like this:
		\begin{verbatim}
HoldsPair<int, double> pairIntDbl;
		\end{verbatim}

	\item You are instructing the compiler to create a \textbf{class} for you using the \textbf{template} and \textbf{instantiate} it for the \textbf{types specified} as \textbf{template arguments} (\verb|int| and \verb|double| in this case).

	\item Thus, in the case of \textbf{templates}, \vocab{instantiation} is the act of process of creating a \textbf{specific type} using one ore more \textbf{template arguments}.

		\pbreak

	\item On the other hand, there may be situations that require you to \textbf{explicitly define} a (different) \textbf{behavior} of a \textbf{template} when \textbf{instantiated} with a \textbf{specific type}.

	\item This is where you \vocab{specialize a template} (or \textbf{behavior thereof}) \underline{for that \textbf{type}}.

	\item A \vocab{specilization} of a \textbf{template class} \verb|HoldsPair| \underline{when \textbf{instantiated}} with \textbf{template parameters} both of type \verb|int| would look like this:

		\begin{verbatim}
template<> class HoldsPair<int, int>
{
	// implementation code here
}
		\end{verbatim}

	\item Needless to say, code that \textbf{specializes} a \textbf{template} must follow the \textbf{template definition}.

	\item \fbox{\textbf{Listing 14.5} Demonstrates Template Specialization}


		\begin{itemize}
    \item \textbf{Template specialization} allows you to provide a custom implementation of a class template for a specific combination of types. In this example, \verb|HoldsPair<int, int>| has a specialized version that behaves differently from the general template.

	\item \boxednote{

			Specialization is \textbf{useful} when a \textbf{particular type combination} requires \\ 
		\textbf{\underline{different behavior, optimizations, or additional functionality}} that\\ 
		\underline{\textbf{does not apply to all possible template types}}.
	}

    \item The general template \verb|HoldsPair<T1, T2>| provides a reusable implementation for any pair of types, while the specialized version \verb|HoldsPair<int, int>| allows the programmer to override that behavior for integers.

    \item Template specialization avoids duplicating an entire class for every type combination while still allowing precise control when a specific case needs custom handling.

    \item In this example, the specialized \verb|GetFirstValue()| function prints a message before returning the integer value, demonstrating how specialization can modify the behavior of individual template instances.
\end{itemize}
\end{itemize}

\newpage

\subsubsection{Template Classes and 'static' Members}
\begin{itemize}
	\item We learned how code in \textbf{templates} \textit{begins to exist} for the \textbf{compiler} \underline{when used} and \underline{not otherwise}.

	\item So, how would \verb|static| \textbf{member attributes} functions within a \textbf{template class}?

	\item In \textbf{lesson 9}, we've learned that \textbf{declaring} a \textbf{class member} \verb|static| results in the \underline{\textbf{member being shared}} across \textbf{all instances} of a \textbf{class}.

		\pbreak

\item It's similar with a \textbf{template class} too, save for the fact that a \textbf{static member} is \underline{\textbf{shared across all objects}} of a \textbf{template class} with the \textbf{\underline{same template instantiation}}.

\item So, a \verb|static| member \verb|X| within a \textbf{template class} is \underline{static} within \textbf{all instances} of the \textbf{class instantiated} \underline{for} \verb|int|.

\item Similarly, \verb|X| is also static within all instances of the class \textbf{specialized} for \verb|double|, \underline{\textbf{independent}} of the other \textbf{template instantiation} for \verb|int|.

	\pbreak

\item In other words, you can visualize it as the compiler creating \textbf{two versions} of the \textbf{static member variable} in a \textbf{template class}: 
	\begin{itemize}
\item \verb|X_int| for \textbf{template instantiation} as \verb|int| and
\item for \textbf{template instantiations} as \verb|double|
	\end{itemize}

	\pbreak

	\begin{itemize}
    \item A \textbf{template class} does not generate actual code for the compiler until a specific type is used. The compiler creates a separate class definition for each \textbf{template instantiation}.

    \item This means that a \verb|static| member in a template class is not shared across all possible template types; instead, it is shared only among objects that use the \textbf{same template instantiation}.

    \item For example, \verb|MyClass<int>| and \verb|MyClass<double>| are treated as two separate classes by the compiler. Therefore, each one receives its own independent copy of the static member variable.

    \item If a template class contains a static member \verb|X|, then all objects of \verb|MyClass<int>| share the same \verb|X_int|, while all objects of \verb|MyClass<double>| share a separate \verb|X_double|.

    \item Conceptually, the compiler creates separate versions of the static member for each type:
    \begin{itemize}
        \item \verb|X_int| for the \verb|int| template instantiation.
        \item \verb|X_double| for the \verb|double| template instantiation.
    \end{itemize}

    \item Thus, static members in templates are \textbf{shared within each template instantiation}, but are \textbf{independent between different template instantiations}.
\end{itemize}

\item \fbox{\textbf{Listing 14.6} The Effect of Static Variables on Template Class and Instances Thereof}

\end{itemize}

\begin{tcolorbox}
	\begin{itemize}
		\item \textbf{Static member instantation} syntax for a \textbf{template class} is \textbf{not to be missed} here:
		\begin{verbatim}
// static member initialization
template<typename T> int TestStatic<T>::staticVal;
		\end{verbatim}

	\item This follows the pattern:

		\begin{verbatim}
template<template parameters> StaticType
ClassName<Template Arguments>::StaticVarName;
		\end{verbatim}
	\end{itemize}
\end{tcolorbox}

\newpage

\begin{verbatim}
#include <iostream>
using namespace std;

template <typename T>
class TestStatic 
{
public:
	static int staticVal;
};

// static member initialization
template<typename T> int TestStatic<T>::staticVal;
\end{verbatim}

\begin{itemize}
    \item The class \verb|TestStatic| is a \textbf{class template}, meaning the compiler creates a separate version of the class for each type used as the template argument.

    \item The declaration:
\begin{verbatim}
static int staticVal;
\end{verbatim}
    declares a \textbf{static member variable} that is shared among all objects of the same template instantiation.

    \item For example, all objects of \verb|TestStatic<int>| share one copy of \verb|staticVal|, while all objects of \verb|TestStatic<double>| share a separate copy. The static variables are independent because \verb|TestStatic<int>| and \verb|TestStatic<double>| are different classes created by the compiler.

    \item The definition:
\begin{verbatim}
template<typename T> int TestStatic<T>::staticVal;
\end{verbatim}
    is required because the declaration inside the class only tells the compiler that the static member exists. This line actually allocates storage for the static member for each template instantiation.

    \item The compiler conceptually creates separate static variables for each type:
    \begin{itemize}
        \item \verb|TestStatic<int>::staticVal| for the \verb|int| instantiation.
        \item \verb|TestStatic<double>::staticVal| for the \verb|double| instantiation.
    \end{itemize}

    \item Therefore, a static member inside a template class is \textbf{shared among objects of the same instantiated type}, but it is \textbf{separate between different template instantiations}.

		\pbreak

	\item It's very similar to normal static declaration:
		\begin{verbatim}
class Test
{
public:
    static int staticVal;
};

// definition of static member
int Test::staticVal;
		\end{verbatim}
\end{itemize}

\newpage

\subsubsection{Variable Templates, Also Called Variadic Templates}
\begin{itemize}
	\item Let's assume you want to write a \textbf{generic function} that adds two values.

	\textbf{Template function} \verb|Sum()| achieves just that:
	\begin{verbatim}
template <typename T1, typename T2, typename T3>
void Sum(T1& res, T2 num1, T3 num2)
{
	res = num1 + num2;
	return;
}
	\end{verbatim}

\item This is simple, however, if you were required to write \textbf{one single function} that would be capable of \textbf{adding any number of values}, each passed an \textbf{argument}, you would need to make use of \vocab{variable templates} in defining such a function.

\item \textbf{Variable templates} or \textbf{variadic templates} have been part of C++ since C++14.

\item \fbox{\textbf{Listing 14.7} Function Using Variadic Templates Demonstrates Variable Arguments}

\end{itemize}

\pbreak

	\begin{tcolorbox}
		\begin{itemize}
			\item Note the \textbf{ellipis mark ...}
			\item \textbf{Ellipis} in C++ used with \textbf{templates} tell the compiler that the \textbf{template class} or \textbf{function} may accept an \underline{\textbf{arbitrary number of template arguments}} of \underline{any type}.
		\end{itemize}
	\end{tcolorbox}

	\pbreak

	\begin{itemize}
		\item \textbf{Variable templates} are powerful addition to C++.
		\item Programmers using \textbf{variable templates} save themselves the repetitive effot of implementing functions that perform a task in various \textbf{overloaded versions}, creating code that is shorter and simpler to maintain.
	\end{itemize}

	\pbreak
	\begin{tcolorbox}
		\begin{itemize}
			\item Note: C++14 supplies you with an \textbf{operator} that would tell the \textbf{number of template arguments} passed in a call to a \textbf{variable templated}.

			\item In \textbf{Listing 14.7}, you could use this \textbf{operator} inside a function like \verb|Sum()|, like this:

				\begin{verbatim}
int arrNums[sizeof...(Rest)];
// length of array evaluated using sizeof...() at compile time
				\end{verbatim}

			\item You \underline{must not confuse} \verb|sizeof...()| with \verb|sizeof(Type)|.

			\item The \textbf{latter} returns the \textbf{size} of a \textbf{type}; the \textbf{former} returns the \textbf{number of template arguments} sent to a \textbf{variadic template}.
		\end{itemize}
	\end{tcolorbox}

	\newpage

	\begin{itemize}
		\item The support of \textbf{variable templates} has ushered in standard support of \vocab{tuples}: \verb|std::tuple|
		\item \verb|std::tuple| is the \textbf{class template} that implements the tuple and it may be \textbf{instantiated} with a \underline{\textbf{varying number of member elements and types}} thereof.

		\item These may be \textbf{individually accessed} using standard library function \verb|std::get|

		\item \fbox{\textbf{Listing 14.8} Instantiating and Using a 'std::tuple'}
	\end{itemize}

	\newpage

\newpage

\subsubsection{Using 'static\_assert' to Perform Compile-Time Checks}
\begin{itemize}
    \item The \verb|static_assert| keyword allows you to perform \textbf{compile-time checks} by verifying that a condition is true while the program is being compiled.

    \item Unlike the \verb|assert| macro, which checks conditions at \textbf{runtime}, \verb|static_assert| checks conditions before the program is executed. If the condition evaluates to \verb|false|, compilation fails and an error message is displayed.

    \item \verb|static_assert| is useful for catching programming errors early, enforcing requirements, and ensuring that assumptions about types, sizes, or values remain valid.

    \item The syntax is:
\begin{verbatim}
static_assert(condition, "error message");
\end{verbatim}

    \item The condition must be a \textbf{constant expression} that the compiler can evaluate during compilation.

    \item For example:
\begin{verbatim}
#include <iostream>

template <typename T>
class Array
{
private:
    T data[10];

public:
    Array()
    {
        static_assert(sizeof(T) <= 4,
                      "Type is too large for this array");
    }
};

int main()
{
    Array<int> intArray;       // Compiles successfully

    // Array<double> doubleArray;
    // Compilation error if sizeof(double) > 4

    return 0;
}
\end{verbatim}

    \item In this example, \verb|static_assert| ensures that the template type satisfies a requirement before generating the class. If a type violates the condition, the compiler stops and reports the provided error message.

    \item \verb|static_assert| is especially useful with \textbf{templates}, because templates can accept many different types, and compile-time checks can prevent invalid template instantiations from being created.

		\pbreak

\item \fbox{\textbf{Listing 14.9} static\_assert}
\end{itemize}

\newpage

\begin{itemize}
	\item \textbf{DO} use \textbf{templates} for implementation of \textbf{generic concepts}.
	\item \textbf{DO} choose \textbf{templates} over \textbf{macros}.
		\pbreak

	\item \textbf{DON'T} forget to use the principles of \verb|const| correctness when programming \textbf{template functions} and \textbf{classes}.
	\item \textbf{DON'T} forget that a \textbf{static member} contained within a \textbf{template class} is \textbf{static} for \underline{every type-specilization} of the \textbf{class}.
\end{itemize}

\newbp

\section{Lambda Expressions}
\begin{itemize}
	\item \textbf{Lambda expressions} are a compact way to \textbf{define} and \textbf{construct} function objects \underline{without a name}.

	\item Introduced in C++11.
\end{itemize}

\pbreak

\subsection{What is a Lambda Expression?}
\begin{itemize}
	\item A \vocab{lambda expression} can be visualized as a \textbf{compact} version of an \textbf{unnamed} \verb|struct| (or \verb|class|) with a \textbf{public} \verb|operator()|.

	\item In that sense, a \textbf{lambda expressions} is a \textbf{function object}.

	\item Let's look into what is a \vocab{function object}:

		\begin{verbatim}
// struct that behaves as a unary function
template <typename elementType>
struct DisplayElement
{
	void operator() (const elementType& element) const
	{
		cout << element << ' ';
	}
};
		\end{verbatim}

	\item THe \textbf{function object} displays an \textbf{object} \verb|element| on the screen using \verb|cout| and is typically used in algos such as \verb|std::for_each()|:

		\begin{verbatim}
// display every integer contained in a vector
for_each (numsInVec.cbegin(),		// start of range
		numsInVec.cend(),			// end of range
		DisplayElement<int>());		// unary function object
		\end{verbatim}
\end{itemize}

\pbreak

\begin{itemize}
	\item A \textbf{lambda expression} compacts the entire code including the \textbf{definition} of the function object into three lines:

		\begin{verbatim}
for_each (numsInVec.cbegin(),		// start of range
		numsInVec.cend(),			// end of range
		[](const int& element) { cout << element << ' '; });
		\end{verbatim}

	\item When the compiler sees the \textbf{lambda expression}, in this case
		\begin{verbatim}
			[](const int& element) { cout << element << ' '; }
		\end{verbatim}

	it automatically \textbf{expands} this expression into a representation that is similar to \verb|struct DisplayElement<int>|:
		\begin{verbatim}
struct NoName
{
	void operator() (const int& element) const
	{
		cout << element << ' ';
	}
};
		\end{verbatim}

	\item Note: \textbf{lambda expressions} are also called \vocab{lambda functions}.
\end{itemize}

\newpage

\subsection{How to Define a Lambda Expression}
\begin{itemize}
	\item The \textbf{definition} of a \textbf{lambda expression} has to start with \textbf{square brackets []}  
	\item These \textbf{brackets} tell the compiler that the \textbf{lambda expression} has started.
	\item They are followed by a \textbf{parameter list}, which is the \underline{same} as the \textbf{parameter list} you would supply your implementation of \verb|operator()| if you were not using a \textbf{lambda expression}.
\end{itemize}

\newpage
\subsection{Lambda Expressions for a Unary Function}
\begin{itemize}
	\item The \textbf{lambda} version of a \textbf{unary operator(Type)} that takes \textbf{one parameter} would be the following:
		\begin{verbatim}
			[](Type paramName) { // lambda expression code here; }
		\end{verbatim}

	\item Note that you can pass the \textbf{parameter by reference} if you so wish:
		\begin{verbatim}
			[](Type& paramName) { // lambda expression code here; }
		\end{verbatim}

	\item \fbox{\textbf{Listing 22.1} for\_each() and lambda expression instead of function object}
\end{itemize}


\newpage

\subsection{Lambda Expression for a Unary Predicate}
\begin{itemize}
	\item A \vocab{predicate} helps make decisions.

	\item A \vocab{unary predicate} is a \textbf{unary expression} that returns a \verb|bool|: true or false.

	\item \textbf{Lambda expressions} can \textbf{return values too}.

	\item Example: \textbf{lambda expression} that returns true for even nums:
		\begin{verbatim}
			[](int num) { return (n & 1) == 0; }
		\end{verbatim}

	\item The nature of the \textbf{return value} in this case tells the compiler that the \textbf{lambda expressions} returns a \verb|bool|

	\item You can use a \textbf{lambda expression} that is a \textbf{unary predicate} in algorithms, such as \verb|std::find_if()|, to fine even numbers in a collection.

	\item \fbox{\textbf{Listing 22.2} std::fine\_if()}
\end{itemize}

\newpage

\subsection{Lambda Expression with State via Capture Lists [...]}
\begin{itemize}
	\item In \textbf{Listing 22.2} we created a \textbf{unary predicate} that returned true if an int was even.

	\item What if we want to make it more generic to see if the int is divisible by a divisor of the user's choice---you need to maintain that "state" (the divisor) in the expression?

		\begin{verbatim}
int divisor = 3;	// initi value
auto element = find_if(numsInVec.cbegin(),
			numsInVec.cend(),
			[divisor](int dividend) { return (dividend % divsor) == 0; });

			or

int divisor = userChoice;

auto pred = [divisor](int x) {
    return x % divisor == 0;
};
		\end{verbatim}

		\begin{itemize}
			\item \textbf{State} in a lambda means the values it \textbf{captures} from its \underline{\textbf{surrounding scope.}}
    \item In this example, the lambda captures \texttt{divisor}:
    \begin{verbatim}
[divisor](int dividend) {
    return (dividend % divisor) == 0;
}
    \end{verbatim}
    \item The captured \texttt{divisor} is stored inside the lambda, allowing it to remember the user's chosen divisor each time it is called.
\end{itemize}

	\item \fbox{\textbf{Listing 22.3} Lambda Expressions that hold state}

\end{itemize}

\newpage
\subsection{The Generic Syntax of Lambda Expressions}
\begin{itemize}
	\item A \textbf{lambda expressions} always starts with \textbf{square brackets []} and can be configured to take \textbf{multiple state variables} separated using commas in a \textbf{capture list [...]}:

		\begin{verbatim}
			[stateVar1, stateVar2](Type& param) { // lambda code here; }
		\end{verbatim}

	\item If you want to ensure that these \textbf{state variables} are \textbf{\underline{modified within a lambda}}, you add keyword \verb|mutable|:
		\begin{verbatim}
			[stateVar1, stateVar2](Type& param) mutuable { // lambda code here; }
		\end{verbatim}

		\pbreak

	\item Note here the variables supplied in the \textbf{capture list []} are \textbf{modifiable} within the lambda, but changes \underline{do not take effect outside of it}.

	\item If you want to ensure that modifications made to the \textbf{state variables} within the lambda are valid \textbf{outside it}, then you use \fbox{\textbf{references}}:
		\begin{verbatim}
			[&stateVar1, &stateVar2](Type& param) { // lambda code here; }
		\end{verbatim}

		\pbreak

	\item \textbf{Lambdas} can take \textbf{multiple input parameters}, separated by commas:
		\begin{verbatim}
			[stateVar1, stateVar2](Type1& var1, Type2& var2) { // lambda code here; }
		\end{verbatim}

		\pbreak

	\item If you want to mention the \textbf{return type} and not leave the disambiguation to the compiler use the $\leftarrow$:
		\begin{verbatim}
			[stateVar1, stateVar2](Type1 var1, Type2 var2) -> ReturnType { // lambda code here; }
		\end{verbatim}

		\pbreak

	\item Finally, the \textbf{compound statement {}} can hold multiple statements sep by a \verb|;| (obviously)

		\begin{verbatim}
			[stateVar1, stateVar2] (Type1 var1, Type2 var2) -> ReturnType
			{
				Statement 1;
				Statement 2;
				return (value or expression);	
			}
		\end{verbatim}

	\item So if you lambda spans multiple lines, you are \textbf{required} to supply an \textbf{explicit return type}.
\end{itemize}

\pbreak

\begin{itemize}
	\item Thus, a \textbf{lambda function} is a compact, fully functional replacement of a \textbf{function object} such as the following:
		\begin{verbatim}
		template<typename Type1, typename Type2>
		struct IsNowTooLong
		{
			// state variables
			Type1 var1;
			Type2 var2;

			// constructor
			IsNowTooLong(const Type1& in1, Type2& in2) : var1(in1), var2(in2) {};

			// the actual purpose
			ReturnType operator()
			{
				Statement 1;
				Statement 2;
				return (value or expression);
			}
		};
		\end{verbatim}
\end{itemize}

\newpage

\subsection{Lambda Expression for a Binary Function}
\begin{itemize}
	\item a \vocab{binary function} takes \textbf{two parameters} and optionally returns a value.

	\item a \vocab{lambda expression} equivalent of the same would be

		\begin{verbatim}
			[...](Type1& param1Name, Type& param2Name) { // lambda code here; }
		\end{verbatim}

	\item \fbox{\textbf{Listing 22.4} Lambda Expression as a Binary Function}
\end{itemize}

\newpage

\subsection{Lambda Expression for a Binary Predicate}
\begin{itemize}
	\item A \textbf{binary function} that returns \verb|true| or \verb|false| to help make a decision is called a \verb|binary predicate|.

	\item These \textbf{\fbox{predicates}} find use in sort algos, such as \verb|std::sort()|, that invoke the \textbf{binary predicate} for any two values in a container to know which one should be placed after the other.

	\item Generic syntax of a \textbf{binary predicate} is:
		\begin{verbatim}
			[...](Type1& param1Name, Type2& param2Name) { // return bool expression; }
		\end{verbatim}

	\item \fbox{\textbf{Listing 22.5} Lambda Expression as a Binary Predicate}
\end{itemize}

\newpage

\begin{itemize}
	\item \textbf{DO} remember that \textbf{lambda expressions} always start with \verb|[]| or \verb|[state1, state2, ...]|
	\item \textbf{DO} remember that unless specified, \textbf{state variables} supplied within a \textbf{capture list} \verb|[]| are \textbf{\underline{not modifiable}} unless you use the keyword \verb|mutable|

	\pbreak

\item \textbf{DON'T} forget that \textbf{lambda expressions} are \textbf{unnamed representations} of a \textbf{class} or a \textbf{struct} with \textbf{operator()}.
\item \textbf{DON'T} forget to use \verb|const| corrent parameter types when writing your \textbf{lambda expressions}:
	\begin{verbatim}
		[](const T& value) { // code; }
	\end{verbatim}
\item \textbf{DON'T} forget to explicitly mention return type when \textbf{lambda expressions} includes \textbf{multiple statements} within the statement block {}
\item \textbf{DON'T} choose lambda expressions over a \textbf{function object} when the lambda gets extremely long and spans multiple statements, for these are redefined in evrey use and do not assist code resuability.
\end{itemize}

\newpage
\section{Understanding Smart Pointers}
\begin{itemize}
	\item C++ can use \textbf{smart pointers} when managing memory on the heap (or free store).
\end{itemize}

\pbreak

\subsection{What are Smart Pointers?}
\begin{itemize}
	\item A \vocab{smart pointer} in C++ is a \textbf{class} with \textbf{overloaded operators}, which \underline{behaves} like a conventional pointer.

	\item Yet, it supplies additional value  by ensuring proper and timely \textbf{destruction} of \textbf{dynamically allocated data} and facilitates a \textbf{well-defined} object lifecycle.
\end{itemize}

\pbreak

\subsubsection{The Problem with Using Conventional (Raw) Pointers}
\begin{itemize}
	\item C++ supplies full flexibility to the programmer in memory allocation, deallocation, and mangement.

	\item One side of this, it makes C++ a powerful language; the other side: it allows the programmer to create memory-related problems, such us memory leaks, when dynamically allocated objects are not correctly released.

	\item Example:
		\begin{verbatim}
SomeClass* ptrData = anObject.GetData();
	// so here should caller or callee invoke delete/free here?
	// no idea.
ptrData->DoSomething();
		\end{verbatim}

	\item This can easily be solved, but it can also easily cause errors.
\end{itemize}

\newpage
\subsubsection{How Do Smart Pointers Help?}
\begin{itemize}
	\item The programmer can choose a smarter way to allocate and manage dynamic data by adopting the use of \textbf{smart pointers}:
		\begin{verbatim}
smart_pointer<SomeClass> spData = anObject.GetData();

// use a smart pointer like a conventional pointer!
spData->Display();
(*spData).Display();

// don't have to worry about deallocation
// (the smart pointer's destructor does it for you!)
		\end{verbatim}

	\item Thus, \textbf{smart pointers} behave like \textbf{conventional (or raw) pointers} but supply useful features via their \underline{\textit{\textbf{overloaded operations}}} and \underline{\textbf{\textit{destructors}}} to ensure that \textbf{dynamically allocated data} is destroyed in a timely manner.
\end{itemize}

\newpage
\subsection{How Are Smart Pointers Implemented?}

\begin{itemize}
	\item Smart pointers, like above, can mimic raw pointers b/c of their \textbf{overloaded derefencing operator (*)} and \textbf{member selection operator }(\verb|->|).
	\item \textbf{Operator overloading} was discussed in Lesson 12.

	\item Additionally, almost all good \textbf{smart pointer classes} are \underline{\textbf{template classes}} that containa generic implementation of their functionality (so that you can manage a type of your choice on the heap).

	\item Being \textbf{templates}, they are versatile and can be \textbf{specialized} to manage an \textbf{object} of a \textbf{type} of your choice.

	\item \fbox{\textbf{Listing 26.1} Minimal Essential Components of a Smart Pointer Class}

		\pbreak

	\item The listing displays the implemention of the \textbf{two operators} * and \verb|->| and assists in having this template class behave like a conventional pointer.

	\item For instance, you would use the \textbf{smart pointer} on an \textbf{object} of type \verb|class Tuna|, you would instantiate it like this:
		\begin{verbatim}
smart_pointer<Tuna> smartTuna(new Tuna);
smart_pointer->Swim();
(*smart_pointer).Swim();
		\end{verbatim}

	\item This is a very minimal example and don't have great benefits over a conventional pointer, however, it's nice to have the \textbf{constructor} accept a pointer that is saved as the \textbf{internal pointer object} in the \textbf{smart pointer class}.

		And the \textbf{destructor} automatically freeing this pointer for \textbf{automatic memory release}.

\end{itemize}
		\begin{tcolorbox}
			\begin{itemize}
				\item The implementation that makes a \textbf{smart pointer} really "smart" is the implementation of the \textbf{copy constructor}, the \textbf{assignment operator}, and the \textbf{destructor}.

				\item They determine the behavior of the \textbf{smart pointer object} when it is \textbf{passed across functions}, when is \textbf{assigned}, or when it \textbf{goes out of scope} (that is, gets \textbf{destructed})

				\item So, before looking at a \textbf{complete smart pointer implementation}, you should understand some \textbf{smart pointer types}.
			\end{itemize}
		\end{tcolorbox}
\newpage

\subsection{Types of Smart Pointers}
\begin{itemize}
	\item The \textbf{management} of the \textbf{memory resource} (that is, the \textbf{\underline{ownership model implemented}}) is what sets \textbf{smart pointer classes} apart.

	\item \textbf{Smart pointers} decide what they do with the resource when they are \textbf{copied} and \textbf{assigned} to.

	\item Classification of \textbf{smart pointers} is actually a classification of their \textbf{memory resource management} strategies. These are
		\begin{itemize}
		\item Deep copy
		\item Copy on Write (COW)
		\item Reference counted
		\item Reference linked
		\item Destructive copy
		\end{itemize}

	\item Let's look into these before studying C++'s \verb|std::unique_ptr| (just an example of one of C++'s smart pointers)
\end{itemize}

\newpage
\subsubsection{Deep Copy}
\begin{itemize}
	\item In a \textbf{smart pointer} that implements \textbf{deep copy}, every \textbf{smart pointer instance} holds a \textbf{complete copy} of the \textbf{object} that is begin managed.

	\item Whenever the \textbf{smart pointer} is \textbf{copied}, the \textbf{object pointed to} is \underline{also \textbf{copied}} (thus, \textbf{deep copy}).

	\item When the \textbf{smart pointer} goes out of scope, it \textbf{releases the memory} it points to (via the \textbf{destructor})

		\pbreak

	\item Although the \textbf{deep-copy-based smart pointer} does not seem to render any value over \textbf{passing objects by values}, its advantage becomes apparent in the treatment of \textbf{polymorphic objects} as seen in the following, where it can \textbf{avoid \vocab{slicing}}.

		\begin{verbatim}
// examaple of Slicing when passing polymorphic objects by value 
// fish is a base class for Tuna and Carp, Fish::Swim() is virtual
void MakeFishSwim(Fish aFish)	// note parameter type
{
	aFish.Swim();	// virtual function
}

// ... some function
Carp freshWaterFish;
MakeFishSwim(freshWaterFish);	// carp will be 'sliced' to Fish
// slicing: only the Fish part of Carp is sent to MakeFishSwim()

Tuna marineFish;
MakeFishSwim(marineFish);	// slicing again

// so its an issue cause passing a polymorphic object by value 
// slices off the derived-class portion, 
// so polymorphism is lost and virtual functions 
// behave as if the object were only the base class.
		\end{verbatim}

		\pbreak
		
	\item \textbf{Slicing} issues are resolved when the programmer chooses a \textbf{deep-copy smart pointer}, as shown in \fbox{\textbf{Listing 26.2}}

	\item \verb|deepcopy_smart_ptr| implements a \textbf{copy constructor} that allows a \textbf{deep copy} of the \textbf{polymorphic object} via a \verb|Clone()| function that the \textbf{class} needs to implement.

	\item Similarly, it implements a \textbf{copy assignment operator}.

	\item Again, for the sake of simplicity, it is taken for granted in this ex that the \textbf{virtual function} implemented by the \textbf{base class} \verb|Fish| is called \verb|Clone()|.

	\item Typically \textbf{smart pointers} that implement \textbf{deep-copy models} have this function supplied either as a \textbf{template parameter} or a \textbf{function object}.

		\pbreak

	\item Thus, when the \textbf{smart pointer} itself is passed as a \textbf{pointer} to \textbf{base class} type \verb|Fish|:
		\begin{verbatim}
deepcopy_smart_ptr<Carp> freshWaterFish(new Carp);
MakeFishSwim(freshWaterFish);	// carp will not be 'sliced'
		\end{verbatim}

		\pbreak

	\item The \textbf{deep copy} implemented in the \textbf{smart pointer's constructor} kicks in to ensure that the \textbf{object} being \textbf{passed} is \underline{not sliced}, even though syntactically only the \textbf{base part} of it is required by the destination function \verb|MakeFishSwim()|

		\pbreak

	\item The disadvantage of the \textbf{deep-copy-based} mechanism is performance.
	\item Instead the programmer may want to pass a \textbf{base type pointer (conventional pointer, Fish*)} to functions such as \verb|MakeFishSwim()|

\end{itemize}

\newpage
\subsubsection{Copy on Write Mechanism}
\begin{itemize}
	\item \vocab{Copy on Write (COW)} attempts to \textbf{optimize performance} of \textbf{deep-copy smart pointers} by \textbf{sharing pointers} until the first attemp at writing to the \textbf{object} is made.

	\item One the \textbf{first attempt} at invoking a \textbf{non-const function}, a \fbox{\textbf{COW pointer}} typically \textbf{creates a copy} of the \textbf{object} on which the \textbf{non-const function} is invoked, whereas other \textbf{instances} of the \textbf{pointer} \underline{continue \textbf{sharing} the \textbf{source object}}.

		\pbreak

	\item \textbf{COW} has its fair share of fans.
	\item For those that swear by \textbf{COW}, implementing \textbf{operators (*)} and \textbf{(}\verb|->|\textbf{)} in their \verb|const| and non-\verb|const| versions is key to the functionality of the \textbf{COW pointer}.

	\item The latter \textbf{creates a copy}.

		\pbreak

	\item Know your shit or you will copy too little or too much
\end{itemize}

\newpage
\subsubsection{Reference-Counted Smart Pointers}
\begin{itemize}
	\item \textbf{Reference counting} in general is a mechanism that keeps a count of the \textbf{number of users} of an \textbf{object}.

	\item When the count reduces to \textbf{zero}, the \textbf{object is released}.

	\item So, \textbf{reference counting} makes a very good mechanism for \textbf{sharing objects} without having to \textbf{copy} them.

		\pbreak

	\item Such \textbf{smart pointers}, when \textbf{copied}, need to have the \textbf{reference count} of the \textbf{object} in question \textbf{incremented}.

	\item There are at least \textbf{two} popular ways to keep this \textbf{count}:
		\begin{enumerate}
			\item Reference count maintained in the \textbf{object being pointed to}.
			\item Reference count maintained by the \textbf{pointer class} in a \textbf{shared object}.
		\end{enumerate}
\end{itemize}

\pbreak

\begin{itemize}
	\item The \textbf{first variant} where the \textbf{reference count} is maintained in the \textbf{object} is called \vocab{intrusive reference counting} because the \textbf{object} needs to be modified.

	\item The \textbf{object} in this case maintains, increments, and supplies the \textbf{reference count} to any \textbf{smart pointer class} that magnages it.

	\item Incidentally, this is the approach chosen by \textbf{COM}.

		\pbreak

	\item The \textbf{second variant} where the \textbf{reference count} is maintained in a \textbf{shared object} is a mechanism where the \textbf{smart pointer class} can keep the reference count on the free store (a dynamically allocated integer, for example) and \textbf{when copied}, the \textbf{copy constructor} increments this value.
\end{itemize}

\pbreak

\begin{itemize}
	\item Therefore, the \textbf{reference-counting} mechanism makes it pertinent that the programmer works with the \textbf{smart pointers} only when using the object.

	\item A \textbf{smart pointer} managing the \textbf{object} and a \textbf{raw pointer} pointing to it is a bad idea because the \textbf{smart pointer} (smartly) releases the object when the count maintained by it goes down to zero, but the \textbf{raw pointer} continues poiting to the part of the memory that no longer belongs to your application.

	\item Similarly, \textbf{reference counting} can cause issues peculiar to their situation: \textbf{2 objects} that \textbf{hold a pointer to each other} are never released b/c their \textbf{cyclic dependency} holds their reference counts at a minimum of 1.
\end{itemize}

\newpage
\subsubsection{Reference-Linked Smart Pointers}
\begin{itemize}
	\item \vocab{Reference-linked smart pointers} are ones that don't proactively count th enumber of \textbf{references} of using the \textbf{object}; rather, they just need to know when the \textbf{number comes down to zero} so that the \textbf{object} can be released.

	\item They are called \vocab{reference-linked} because their implementation is based on a \textbf{double-linked list}.

	\item When a \textbf{new smart pointer} is created by \textbf{copying} an existing one, it is \textbf{appended} to the list.

	\item When a \textbf{smart pointer} goes \textbf{out of scope} or is \textbf{destroyed}, the destructor \textbf{de-indexes} the \textbf{smart pointer} from the list.

\item \textbf{Reference counting} also suffers from the problem caused by \textbf{cylic dependency}, as applicable to \textbf{reference-counted pointers}.
\end{itemize}

\newpage
\subsubsection{Destructive Copy}
\begin{itemize}
\item \vocab{Destructive copy} is a mechanism where a \textbf{smart pointer}, when \textbf{copied}, transfers \textbf{complete ownership} of the object being handled to the destination and resets itself:

	\begin{verbatim}
destructive_copy_smartptr<SampleClass> smartPtr(new SampleClass());

SomeFunc(smartPtr);	// ownership transferred to SomeFunc
// don't use smartPtr in the caller any more!
	\end{verbatim}

\item Although this mechanism is obviously intuitive to use, the advantage supplied by \textbf{destructive copy smart pointers} is that they ensure that at any point in time, \textbf{only one active pointer points} to an \textbf{object}.

\item So, they make good mechanisms for \textbf{returning pointers} from functions, and are of use in scenarios where you can use their \textbf{"destructive"} properties to your advantage.

\item The implementation of \textbf{destructive copy pointers} deviate from standard, recommended C++ programming techniques, as shown in \textbf{Listing 26.3}.

\item \fbox{\textbf{Listing 26.3} A Sample Destructive Copy Smart Pointer}
\end{itemize}

\pbreak


\begin{tcolorbox}
	\begin{itemize}
		\item The \textbf{copy constructor} and \textbf{copy assignment operators} that are critical to the implementation of \textbf{destructive copy smart pointers} as shown in \textbf{Listing 26.3} also attract maximum criticism.

		\item Unlike most C++ classes, this \textbf{smart pointer class} \underline{cannot} have the \textbf{copy constructor} and \textbf{assignment operator} acept \verb|const| \textbf{references}, as it needs to invalidate the \textbf{source} after copying it.

		\item This is not only a deviation from traditional copy-constructor and assignment-operator semantics, but also makes using the smart pointer class counter intuitive.

		\item Few expect the \textbf{copy source} or the \textbf{assignment source} to be damaged after a \textbf{copy} or \textbf{assignment} step. counter intuitive.

		\pbreak
		
	\item The fact that such \textbf{smart pointers destroy} the \textbf{source} also makes them unsuitable for the use in the \textbf{STL containers}, such as \verb|std::vector| or any other dynamic collection class that you might use.

	\item These \textbf{containers} need to \textbf{copy} your content internally and end up invalidating the pointers.

		\pbreak

	\item So, for more than one reason, you are \textbf{advised to avoid using \underline{destructive copy smart pointers}} in your programs.
	\end{itemize}
\end{tcolorbox}

\newpage
\subsubsection{Using std::unique\_ptr}
\begin{itemize}
	\item \verb|std::unique_ptr| \underline{does not allow} \textbf{copy} or \textbf{assignment}
	\item Use \verb|#include <memory>| to use \verb|std::unique_ptr|

\item The \verb|std::unique_ptr| is a \textbf{simple smart pointer} similar to what's shown in \textbf{Listing 26.1}, but with a \textbf{private copy constructor} and \textbf{assignment operator} to \underline{disallow copy} via passing as an \textbf{argument to a function by value} or \textbf{copy via assignment}

\item \fbox{\textbf{Listing 26.4} Using std::unique\_ptr}
\end{itemize}

\newpage
\begin{tcolorbox}
	\begin{itemize}
		\item \textbf{Listing 26.4} demonstrated that the \verb|unique_ptr| \textbf{\underline{doesn't support copy}}:

			\begin{verbatim}
copySmartFish = smartFish;		// error: operator= is private
			\end{verbatim}

		\item It however \textbf{does support \fbox{move semantics}}.

			Therefore, an option that would work is

			\begin{verbatim}
unique_ptr<Fish> sameFish(std::move(smartFish));
// smartFish is empty henceforth
			\end{verbatim}

		\item If you were ever to write a \textbf{lambda expression} that would need to capture an \verb|unique_ptr|, then you would use \verb|std::move()| in your \textbf{lambda capture} as supported by C++14.

			\pbreak

			\begin{verbatim}
std::unique_ptr<char> alphabet(new char);
*alphabet = 's';
auto lambda = [capture = std::move(alphabet)]() 
{ std::cout << *capture << endl; }

// alphabet is empty henceforth as contents have been 'moved'

lambda();
			\end{verbatim}
	\end{itemize}
\end{tcolorbox}

\newpage
\subsection{Smart pointers std::}
\subsubsection{std::unique\_ptr}
\begin{itemize}
    \item A smart pointer that owns a dynamically allocated object exclusively.
    \item Automatically deletes the object when the \texttt{std::unique\_ptr} goes out of scope (RAII).
    \item Cannot be copied, only moved, ensuring a single owner at all times.
    \item Commonly created with \texttt{std::make\_unique} (C++14+).
    \item Useful for managing resources without manual \texttt{delete} calls.

    \begin{verbatim}
#include <iostream>
#include <memory>

int main() {
    auto ptr = std::make_unique<int>(42);

    std::cout << *ptr << std::endl;   // 42

    // Ownership can be transferred
    auto ptr2 = std::move(ptr);

    if (!ptr)
        std::cout << "ptr is now empty\n";

    std::cout << *ptr2 << std::endl;  // 42
}
    \end{verbatim}

    \item Managing a user-defined object:

    \begin{verbatim}
#include <iostream>
#include <memory>

class Dog {
public:
    Dog() { std::cout << "Dog created\n"; }
    ~Dog() { std::cout << "Dog destroyed\n"; }

    void bark() { std::cout << "Woof!\n"; }
};

int main() {
    auto dog = std::make_unique<Dog>();
    dog->bark();
} // Dog is automatically destroyed here
    \end{verbatim}

	\item Transferring ownership with \texttt{std::move()}:

\begin{verbatim}
#include <iostream>
#include <memory>

void takeOwnership(std::unique_ptr<int> p) {
    std::cout << *p << std::endl;
} // p is destroyed here

int main() {
    auto ptr = std::make_unique<int>(100);

    takeOwnership(std::move(ptr));

    if (!ptr)
        std::cout << "ptr no longer owns the object\n";
}
\end{verbatim}
\end{itemize}

\newpage
\subsubsection{std::shared\_ptr}
\begin{itemize}
    \item A smart pointer that allows multiple owners of the same dynamically allocated object.
    \item Uses reference counting to track how many \texttt{std::shared\_ptr}s own the object.
    \item The object is automatically deleted when the last owning \texttt{std::shared\_ptr} is destroyed or reset.
    \item Can be copied freely, increasing the reference count.
    \item Commonly created with \texttt{std::make\_shared}, which is more efficient than using \texttt{new}.

    \begin{verbatim}
#include <iostream>
#include <memory>

int main() {
    auto ptr1 = std::make_shared<int>(42);
    auto ptr2 = ptr1;    // Shared ownership

    std::cout << *ptr1 << std::endl;          // 42
    std::cout << ptr1.use_count() << std::endl; // 2
}
    \end{verbatim}

    \item Managing a user-defined object:

    \begin{verbatim}
#include <iostream>
#include <memory>

class Dog {
public:
    Dog() { std::cout << "Dog created\n"; }
    ~Dog() { std::cout << "Dog destroyed\n"; }

    void bark() { std::cout << "Woof!\n"; }
};

int main() {
    auto dog1 = std::make_shared<Dog>();
    auto dog2 = dog1;   // Both point to the same Dog

    dog1->bark();

    std::cout << dog1.use_count() << std::endl; // 2
} // Dog destroyed after both pointers go out of scope
    \end{verbatim}

    \item Sharing ownership by passing to a function:

\begin{verbatim}
#include <iostream>
#include <memory>

void printValue(std::shared_ptr<int> p) {
    std::cout << *p << std::endl;
    std::cout << p.use_count() << std::endl;
}

int main() {
    auto ptr = std::make_shared<int>(100);

    std::cout << ptr.use_count() << std::endl; // 1

    printValue(ptr);   // Copies the shared_ptr

    std::cout << ptr.use_count() << std::endl; // 1
}
\end{verbatim}

    \item Releasing ownership:

\begin{verbatim}
#include <iostream>
#include <memory>

int main() {
    auto ptr1 = std::make_shared<int>(10);
    auto ptr2 = ptr1;

    std::cout << ptr1.use_count() << std::endl; // 2

    ptr2.reset();   // Releases one owner

    std::cout << ptr1.use_count() << std::endl; // 1
}
\end{verbatim}

\end{itemize}

\newpage
\subsubsection{std::weak\_ptr}
\begin{itemize}
    \item A smart pointer that observes an object managed by a \texttt{std::shared\_ptr}.
    \item Does not own the object and does not increase the reference count.
    \item Used to prevent circular references between \texttt{std::shared\_ptr}s.
    \item Must be converted to a \texttt{std::shared\_ptr} using \texttt{lock()} before accessing the object.
    \item Can check whether the object still exists using \texttt{expired()}.

    \begin{verbatim}
#include <iostream>
#include <memory>

int main() {
    auto shared = std::make_shared<int>(42);
    std::weak_ptr<int> weak = shared;

    std::cout << shared.use_count() << std::endl; // 1

    if (auto ptr = weak.lock()) {
        std::cout << *ptr << std::endl; // 42
    }
}
    \end{verbatim}

    \item Detecting when an object has been destroyed:

    \begin{verbatim}
#include <iostream>
#include <memory>

int main() {
    std::weak_ptr<int> weak;

    {
        auto shared = std::make_shared<int>(100);
        weak = shared;
    } // shared is destroyed here

    if (weak.expired()) {
        std::cout << "Object no longer exists\n";
    }
}
    \end{verbatim}

    \item Safely accessing the object with \texttt{lock()}:

\begin{verbatim}
#include <iostream>
#include <memory>

void printValue(std::weak_ptr<int> weak) {
    if (auto ptr = weak.lock()) {
        std::cout << *ptr << std::endl;
    } else {
        std::cout << "Object has been destroyed\n";
    }
}

int main() {
    auto shared = std::make_shared<int>(50);

    printValue(shared);

    shared.reset();

    printValue(shared);
}
\end{verbatim}

\end{itemize}

\end{document}

